%\includeonly{sensitivity,backgrounds}
%\includeonly{software}
%\includeonly{theory,mueconv}
\documentclass{thesis}
\usepackage{dcolumn}
%\usepackage[top=2cm,bottom=2cm,left=1.5cm,right=1.5cm]{geometry}
\usepackage{feynmp}
\usepackage{graphicx}
%%\usepackage{dblfloatfix}
%\usepackage{minted}
\usepackage{xspace}
%\usepackage[heightadjust=all,valign=t]{floatrow,fr-subfig}
\usepackage{siunitx}
\usepackage{eval}
%\usepackage{ragged2e}
%\usepackage{array}
\usepackage[printonlyused]{acronym}
\usepackage[strings]{underscore}
\usepackage{babel}
\usepackage{pbox}
%\usepackage{pythontex}
%%\usepackage{amsmath}
%%\usepackage{amssymb}
%\usepackage{xfrac}
%\usepackage{xcolor}
\usepackage[sharp]{easylist}
\usepackage{multirow}
\usepackage{listings}
\usepackage{hyperref}
%\usepackage[caption=false]{subfig}
%\hypersetup{pdfborder={0 0 0}, urlcolor=blue}

\hypersetup{%
  colorlinks=true,% hyperlinks will be coloured
  linkcolor=blue,% hyperlink text will be green
  urlcolor=blue,% hyperlink text will be green
  linkbordercolor=red% hyperlink border will be red
}

\newcolumntype{d}[1]{D{.}{\cdot}{#1} }
\newcolumntype{a}{D{-}{\le\theta<}{4} }

\lstdefinestyle{customc}{
  belowcaptionskip=1\baselineskip,
  breaklines=true,
  frame=L,
  xleftmargin=\parindent,
  language=sh,
  showstringspaces=false,
  basicstyle=\footnotesize\ttfamily,
  keywordstyle=\bfseries\color{green!40!black},
  commentstyle=\itshape\color{purple!40!black},
  identifierstyle=\color{blue},
  stringstyle=\color{orange},
}

\lstdefinestyle{customasm}{
  belowcaptionskip=1\baselineskip,
  frame=L,
  xleftmargin=\parindent,
  language=[x86masm]Assembler,
  basicstyle=\footnotesize\ttfamily,
  commentstyle=\itshape\color{purple!40!black},
}

\lstset{escapechar=@,style=customc}

\makeatletter
\Hy@AtBeginDocument{%
  \def\@pdfborder{0 0 1}% Overrides border definition set with colorlinks=true
  \def\@pdfborderstyle{/S/U/W 0}% Overrides border style set with colorlinks=true
                                % Hyperlink border style will be underline of width 1pt
}
\makeatother

\makeatletter
\@ifpackageloaded{hyperref}{%
\hypersetup{%
pdftitle = {},
pdfsubject = {PhD Thesis for Ben Krikler},
pdfkeywords = {COMET, mu-e conversion, CLFV, physics},
pdfauthor = {\textcopyright\ Ben Krikler},
}
}{}
\makeatother

\input{figs/feynman/fixFeynmpForPdf}
%%%%%%%%%%%%%%%%%%%%%%%%%%%%%%%
% Values
%%%%%%%%%%%%%%%%%%%%%%%%%%%%%%%
\newcommand{\senseSindrum}{\ensuremath{7\times10^{-13}}\xspace}
\newcommand{\sensePI}{\ensuremath{3\times10^{-15}}\xspace}
\newcommand{\sensePII}{\ensuremath{3\times10^{-17}}\xspace}
\newcommand{\senseMuEEE}{\ensuremath{\times10^{-16}}\xspace}
\newcommand{\muStopsPI}{\ensuremath{2\times10^{19}}\xspace}
%\newcommand{\maxRatesPI}{\ensuremath{6\times10^{-6}}\xspace}
\newcommand{\limitMEG}{\ensuremath{5.7\times10^{-13}}\xspace}
 

%%%%%%%%%%%%%%%%%%%%%%%%%%%%%%%
% decay expressions
%%%%%%%%%%%%%%%%%%%%%%%%%%%%%%%
\newcommand{\mueconv}{\ensuremath{\mu$-$e\textrm{ conversion}}\xspace}
\newcommand{\muec}{\ensuremath{\mu^{-}+N \rightarrow e^{-}+N}\xspace}
\newcommand{\mue}{\ensuremath{\mu^{-}-e^{-}}\xspace}
\newcommand{\mutoe}{\ensuremath{\mu^{-}$$-$$e^{-}}\xspace}
\newcommand{\mueg}{\ensuremath{\mu\rightarrow e\gamma}\xspace}
\newcommand{\muegamma}{\ensuremath{\mu^+\rightarrow e^+\gamma}\xspace}
\newcommand{\muThreeE}{\ensuremath{\mu^+\rightarrow e^+e^-e^+}\xspace}

%%%%%%%%%%%%%%%%%%%%%%%%%%%%%%%
% Experiment names
%%%%%%%%%%%%%%%%%%%%%%%%%%%%%%%
\newcommand{\sindrumII}{\mbox{SINDRUM-II}\xspace}
\newcommand{\phaseI}{\mbox{Phase-I}\xspace}
\newcommand{\phaseII}{\mbox{Phase-II}\xspace}
\newcommand{\alcap}{AlCap\xspace}
%\newcommand{\COMET}{COMET\xspace}
\newcommand{\COMET}{{COMET}\xspace}
%\newcommand{\COMET}{{\fontfamily{phv}\selectfont{}COMET}\xspace}

%%%%%%%%%%%%%%%%%%%%%%%%%%%%%%%
% short-cuts
%%%%%%%%%%%%%%%%%%%%%%%%%%%%%%%
\newcommand {\etal}{\mbox{et al.}\xspace} %et al. - no preceding comma
\newcommand {\ie}{\mbox{i.e.}\xspace}     %i.e.
\newcommand {\eg}{\mbox{e.g.}\xspace}     %e.g.
\newcommand {\etc}{\mbox{etc.}\xspace}     %etc.
\newcommand {\vs}{\mbox{\sl vs.}\xspace}      %vs.
\newcommand {\mdash}{\ensuremath{\mathrm{-}}} % for use within formulas
\newcommand {\CHECK}[1]{{\bf \emph{((CHECK: #1))} }\xspace}
\newcommand {\code}[1]{\mintinline{c++}{#1}}
\newcommand {\doxygenUrl}{http://www.hep.ph.ic.ac.uk/~bek07/comet/}
\newcommand {\clfv}{cLFV\xspace}
\newcommand {\degree}{\ensuremath{^\circ}\xspace}
\newcommand {\dEdx}{\ensuremath{dE/dx}\xspace}

%%%%%%%%%%%%%%%%%%%%%%%%%%%%%%%
% structuring commands
%%%%%%%%%%%%%%%%%%%%%%%%%%%%%%%
\newcommand{\Fig}[1]{Fig.~\ref{fig:#1}\xspace }
\newcommand{\fig}[1]{Fig.~\ref{fig:#1}\xspace }
\newcommand{\tab}[1]{Table~\ref{tab:#1}\xspace }
\newcommand{\Tab}[1]{Table~\ref{tab:#1}\xspace }
\newcommand{\Eq}[1]{Eq.~\ref{eq:#1}\xspace }
\newcommand{\eq}[1]{Eq.~\ref{eq:#1}\xspace }
\newcommand{\figlabel}[1]{\label{fig:#1}\xspace }
\newcommand{\tablabel}[1]{\label{tab:#1}\xspace }
\newcommand{\eqlabel}[1]{\label{eq:#1}\xspace }

\newcommand{\heading}[1]{\section{#1}}
\newcommand{\subheading}[1]{\subsection{#1}}
\newcommand{\subsubheading}[1]{\subsubsection*{#1}}
\newcommand{\subsubsubheading}[1]{\subsubsection*{\normalsize{#1}}\addcontentsline{toc}{subsubsection}{\emph{#1}}}
% disable subsubsections in the TOC (Need this approach for revtex)
\makeatletter
\def\l@subsubsection#1#2{}
\makeatother


%\input{figs/figures.tex}

\renewcommand{\textfraction}{0.01}
\renewcommand{\topfraction}{0.9} 
\renewcommand{\bottomfraction}{0.9} 
\renewcommand{\floatpagefraction}{0.85}

%% Define the thesis title and author
\title{Sensitivity and Background Estimates for \phaseII of the \COMET Experiment}
\author{Benjamin Edward Krikler}

\begin{document}
% The VarUnits and VarNum commands are declared in the eval.sty file
% They create new commands of the form \Var<arg1>

%%%%%%% SIGNAL
\VarUnits{MomThreshold}{104.2}{MeV/c}
\VarUnits{MomThresholdHigh}{105.5}{MeV/c}

% Signal acceptance
\VarNum{AcceptanceGeom}{0.22}
\VarNum{AcceptanceMom}{0.70}
\VarNum{AcceptanceTime}{0.53}
\VarUnits{TotalSignalAcceptance}{5.7}{\%}
\VarNum{TotalSignalAcceptanceDec}{0.057}

% Muon stopping
\VarNum{MuStopsPerPOT}{1.61e-3}
\VarNum{TotalMuStops}{1.10e18}

% Signal Sensitivity
\VarNum{TotalPOTCDRRunTime}{8.70E+20}
\newcommand{\VarCDRRunTime}[1][1]{\scinot[#1]{2.00e7}}
\VarNum{PredictedSESCDRRunTime}{2.04e-17}
%\VarNum{PredictedSESPerYear}{1.29e-17}

\VarNum{TotalPOT}{6.83E+20}
\newcommand{\VarRunTime}[1][3]{\scinot[#1]{1.57e7}}
\VarNum{PredictedSES}{2.6e-17}
\VarNum{PredictedSESPerYear}{1.29e-17}
\VarNum{PredictedSESPerSecond}{4.08e-10}

%%%%%%% BACKGROUNDS
%\VarNum{ExtinctionFactor}{1e-12}
\newcommand{\VarExtinctionFactor}[1][1]{\scinot[#1]{1.00e-11}}
% DIO
\VarNum{DIOPerMuStop}{6.20e-20}
\VarNum{DIOPerPOT}{9.92E-23}
\VarNum{DIOPerSecond}{4.31E-09}
\VarNum{DIOTotal}{0.068}

%RMC
%\VarNum{RMCPerMuStop}{1.05E-19}
%\VarNum{RMCPerPOT}{1.70E-22}
%\VarNum{RMCPerSecond}{7.37E-09}
%\VarNum{RMCTotal}{0.12}
\VarNum{RMCPerMuStop}{5.26E-20}
\VarNum{RMCPerPOT}{8.48E-23}
\VarNum{RMCPerSecond}{3.69E-09}
\VarNum{RMCTotal}{0.058}
			
% RPC
\VarNum{PiStopsPerPOT}{4.33e-7}
\VarNum{DetectedEsPerRPC}{1.05e-5}
\VarUnits{RPCLifetime}{18.6}{ns}
\VarNum{RPCTimingEfficiency}{9.50e-15}
\VarNum{RPCPromptPerPOT}{1.82E-24}
\VarNum{RPCDelayedPerPOT}{1.73e-27}
\VarNum{RPCPromptPerSec} {7.91E-11}
\VarNum{RPCDelayedPerSec}{7.511E-14}
\VarNum{RPCPromptTotal} {1.24E-03}
\VarNum{RPCDelayedTotal}{1.18E-06}
%\VarUnits{RPCLifetime}{21.3}{ns}
%\VarNum{RPCTimingEfficiency}{5.49E-13}
%\VarNum{RPCPromptPerPOT}{1.82E-24}
%\VarNum{RPCDelayedPerPOT}{9.98E-26}
%\VarNum{RPCPromptPerSec} {7.91E-11}
%\VarNum{RPCDelayedPerSec}{4.34E-12}
%\VarNum{RPCPromptTotal} {1.24E-03}
%\VarNum{RPCDelayedTotal}{6.82E-05}


% Antiprotons
\VarNum{BgAntiprotonsPerPOT}{4.34e-22}
\VarNum{BgAntiprotonsPiDelayedPerPOT}{1.95E-30}
\VarNum{BgAntiprotonsPiPromptPerPOT}{3.56E-29}

\VarNum{BgAntiprotonsPerSec}         {1.89E-08}
\VarNum{BgAntiprotonsPiDelayedPerSec}{8.49E-17}
\VarNum{BgAntiprotonsPiPromptPerSec} {1.55E-15}

\VarNum{BgAntiprotonsTotal}         {0.296}
\VarNum{BgAntiprotonsPiDelayedTotal}{1.33E-09}
\VarNum{BgAntiprotonsPiPromptTotal} {2.43E-08}





% Beam backgrounds
\VarNum{BeamBgGeometric}{8.93e-09}
\VarNum{BeamBgMomentum}{3.13e-05}
\VarNum{BeamBgTiming}{5.25e-12}
\VarNum{BeamBgPromptPerPOT}{2.80E-24}
\VarNum{BeamBgDelayedPerPOT}{1.47E-24}
\VarNum{BeamBgPromptPerSec}{1.22e-10}
\VarNum{BeamBgDelayedPerSec}{6.39E-11}
\VarNum{BeamBgPromptTotal} {1.91E-03}
\VarNum{BeamBgDelayedTotal}{1.00E-03}

% cosmics
\VarNum{CosmicRatePerSecond}{1.87e-8}
\VarNum{CosmicRateTotal}{0.294}

%% Total rates
\VarNum{TotalBgPerSecond}{4.59E-08}
\VarNum{TotalBgPhasII}{0.720}


%% Define the un-numbered front matter (cover pages, rubrik and table of contents)
\begin{frontmatter}
  %% Title
%!!UPDATE THIS
\titlepage[of Imperial College London]%
{A dissertation submitted to Imperial College London\\
  for the degree of Doctor of Philosophy}

%% Abstract
\begin{abstract}%[\smaller \thetitle\\ \vspace*{1cm} \smaller {\theauthor}]
  %\thispagestyle{empty}
\end{abstract}


%% Declaration
\begin{declaration}
  This dissertation is the result of my own work, except where explicit
  reference is made to the work of others, and has not been submitted
  for another qualification to this or any other university. This
  dissertation does not exceed the word limit for the respective Degree
  Committee.
  \vspace*{1cm}
  \begin{flushright}
	Benjamin Edward Krikler
  \end{flushright}
\end{declaration}


%% Acknowledgements
\begin{acknowledgements}

\end{acknowledgements}


%% Preface
%\begin{preface}
%\end{preface}

%% ToC
\tableofcontents

%% Strictly optional!
%\frontquote%
%{Writing in English is the most ingenious torture\\
%   ever devised for sins committed in previous lives.}%
%  {James Joyce}

\end{frontmatter}

%% Start the content body of the thesis
\begin{mainmatter}
  %% Actually, more semantic chapter filenames are better, like "chap-bgtheory.tex"

\newcommand{\FigTheoryMuonDecayCloudChamber}{
\begin{figure}[tb]
\centering 
%\fbox{
\includegraphics[width=0.95\textwidth]{figs/theory/EarlyCloudChamberMuonDecay.pdf}
%}
\caption[One of the earliest cloud chamber photographs of a muon, taken in 1940~\cite{Williams1940102}.]{
One of the earliest cloud chamber photographs of a muon, taken in 1940~\cite{Williams1940102}.
In the left-most image the muon enters the chamber at point A and travels to point F, where it eventually decays to an electron which can be seen faintly leaving the image at point G.
The images to the right are a stereoscopic zoom in on point F, showing the relatively slow and more ionising muon and the faster, less ionising electron.
\figlabel{theory:muonCloudChamber}}
%\footnote{though the author has failed to reproduce the stereoscopic effect with his own eyes}
\end{figure}
}

\newcommand{\FigTheoryHincksPontecorvoMuEGamma}{
\begin{figure}[tb]
\centering 
%\fbox{
\includegraphics[width=0.85\textwidth]{figs/theory/OriginalMuEGammaExperiment.png}
%}
\caption[The setup of the first experiment to look for photons produced during muon decay taken from~\cite{Hincks194802}.]{
The setup of the first experiment to look for photons produced during muon decay taken from~\cite{Hincks194802}.
Cosmic muons arrived from the top, slowing down in the big block of lead, triggering two Geiger-Muller counters (A and B) as they passed and eventually coming to stop in the graphite.
From there, electrons and any potential photons would be detected in the counters above and below the graphite (B and C).  
No photons were seen in coincidence with an electron from muon decay, which lead theorists to hypothesise two distinct neutrino flavours.
\figlabel{theory:originalMEG}}
%\footnote{though the author has failed to reproduce the stereoscopic effect with his own eyes}
\end{figure}
}

\newcommand{\FigTheoryMuEGammViaNeutrino}{
\begin{figure}[t]
\centering 
%\fbox{
%\input{figs/feynman/mu_to_e_gamma_via_SM-Wgamma.tex}
\subfloat[][\figlabel{theory:feyn:muDecay}\ac{SM} Muon Decay]{\includegraphics[width=0.45\textwidth]{figs/feynman/pdfs/mu_decay.pdf}}\hspace{0.03\textwidth}
\subfloat[][\figlabel{theory:feyn:muEGammaViaNu}Neutrinoless Muon Decay with Photon Emission]{\includegraphics[width=0.45\textwidth]{figs/feynman/pdfs/mu_to_e_gamma_via_SM-Wgamma.pdf}}
%}
\caption[Feynman diagrams for Standard Model muon decay and neutrino oscillation-mediated \mutoe gamma]{
%Feynman diagram for the neutrinoless muon decay to a photon and electron mediated by a neutrino oscillation
The outgoing neutrinos produced in the \ac{SM} decay \protect\subref{fig:theory:feyn:muDecay} can be made into an internal neutrino propagator \protect\subref{fig:theory:feyn:muEGammaViaNu} via neutrino oscillations, whilst the addition of a photon ensures conservation the 4-momentum.
Although allowed in the \ac{SM} with neutrino oscillations, the actual rate from such a diagram is well below present experiment sensitivities.
A similar diagram was envisaged to show that the lack of observation of \mueg implied distinict neutrino flavours.
\figlabel{theory:feyn:decay}}
%\footnote{though the author has failed to reproduce the stereoscopic effect with his own eyes}
\end{figure}
}

\newcommand{\FigTheoryMuLFVLimits}{
\begin{figure}[tb]
\centering 
%\fbox{
\includegraphics[width=0.95\textwidth]{figs/theory/MuLFVSearchLimits.pdf}
%}
\caption
[The historic evolution of experimental upper limits on \ac{CLFV} in muon channels]{
The experimental limits on \ac{CLFV} in muon channels has been reduced by some 13 orders of magnitude since the first search in 1947.
Each improvement in the limit is primarily due to the increasing intensity of the muon source:
cosmic muons were used around 1950; accelerator-produced muons followed; surface muon sources further increased the intensity; and nowadays superconducting solenoids will capture and focus the resulting muon beam.
%The improvements in the limit are primarily linked to the intensity of the available muon source: around 1950 cosmic muons were used; accelerator-based muon sources improved this, 
The shaded grey region shows measurements that are currently under preparation.
%, including MEG-II, Mu3e, Mu2e, and COMET \phaseI and \phaseII.
Data until 2015 is taken from the summary tables from~\cite{Bernstein2013hba}.
\figlabel{theory:muModesLimits}}
\end{figure}
}

\newcommand{\FigTheoryMuEConvViaNeutrino}{
\begin{figure}[t]
\centering 
%\fbox{
%\input{figs/feynman/mu_to_e_gamma_via_SM-Wgamma.tex}
\subfloat[][\figlabel{theory:feyn:muecViaNu:W}$\gamma$ Penguin]  {\includegraphics[width=0.43\textwidth]{figs/feynman/pdfs/mu_e_conversion_via_SM-Wgamma.pdf}}\hspace{0.08\textwidth}
\subfloat[][\figlabel{theory:feyn:muecViaNu:Z}$Z$ Penguin]  {\includegraphics[width=0.43\textwidth]{figs/feynman/pdfs/mu_e_conversion_via_SM-nuZ.pdf}}\\
\subfloat[][\figlabel{theory:feyn:muecViaNu:d}$d$-quark Box]{\includegraphics[width=0.43\textwidth]{figs/feynman/pdfs/mu_e_conversion_via_SM-downBox.pdf}}\hspace{0.08\textwidth}
\subfloat[][\figlabel{theory:feyn:muecViaNu:u}$u$-quark Box]{\includegraphics[width=0.43\textwidth]{figs/feynman/pdfs/mu_e_conversion_via_SM-upBox.pdf}}\\
%}
\caption
[Feynman diagram for the neutrinoless muon decay in the presence of an atomic nucleus---\mueconv{}---caused by neutrino oscillations.]{
Feynman diagram for the neutrinoless muon decay in the presence of an atomic nucleus---\mueconv{}---caused by neutrino oscillations.
Compared to \mueg there are now four possible diagrams, so that the rate depends on the nucleus and picks up interference terms.
\figlabel{theory:feyn:muecViaNu}}
%\footnote{though the author has failed to reproduce the stereoscopic effect with his own eyes}
\end{figure}
}


\newcommand{\FigTheoryMuEConvNewPhysics}{
\begin{figure}[tb]
%\vskip1cm
\centering
\subfloat[][\figlabel{theory:feyn:muecNP:HiggsDirect}Exotic Higgs]{ \includegraphics[width=0.16\textheight]{figs/feynman/pdfs/mu_e_conversion_Higgs.pdf}}\hspace{0.045\textwidth}
\subfloat[][\figlabel{theory:feyn:muecNP:Zprime}$Z$-prime]{         \includegraphics[width=0.16\textheight]{figs/feynman/pdfs/mu_e_conversion_Z_prime.pdf}}\hspace{0.045\textwidth}
\subfloat[][\figlabel{theory:feyn:muecNP:Leptoquark}Leptoquarks]{   \includegraphics[width=0.16\textheight]{figs/feynman/pdfs/mu_e_conversion_Leptoquark.pdf}}\\
\subfloat[][\figlabel{theory:feyn:muecNP:HeavyN}Heavy Neutrinos]{   \includegraphics[width=0.2\textheight]{figs/feynman/pdfs/mu_e_conversion_via_heavy_neutrino.pdf}}\hspace{0.02\textwidth}
\subfloat[][\figlabel{theory:feyn:muecNP:HiggsTopLoop}Exotic Higgs]{\includegraphics[width=0.17\textheight]{figs/feynman/pdfs/mu_e_conversion_HiggsTopLoop.pdf}}\hspace{0.02\textwidth}
\subfloat[][\figlabel{theory:feyn:muecNP:SUSY}Supersymmetry]{       \includegraphics[width=0.2\textheight]{figs/feynman/pdfs/mu_e_conversion_via_susy.pdf}}\\
%\subfloat[][\label{fig:FD-Z-h}Z-prime and extended Higgs]{\input{feynman/mu_e_conversion_Z_prime}}
\caption
[Feynman diagrams that produce \mueconv through New Physics models.]{
Feynman diagrams that produce \mueconv through New Physics models.
The upper three diagrams (\protect\subref{fig:theory:feyn:muecNP:HiggsDirect} to \protect\subref{fig:theory:feyn:muecNP:Leptoquark}) all connect to the nucleus via some massive exchange particle, 
whereas the lower three diagrams (\protect\subref{fig:theory:feyn:muecNP:HeavyN} to \protect\subref{fig:theory:feyn:muecNP:SUSY}) all connect via an exchanged photon.
In addition to interactions with the quarks, since \mueconv interacts with the whole nucleus, there are also models where the interaction involves external gluon lines.
\figlabel{theory:feyn:muecNP}}
\end{figure}
}

\newcommand{\FigTheoryMuecVsMueg}{
\begin{figure}[tb]
\centering 
%\fbox{
\includegraphics[width=0.95\textwidth]{figs/theory/MuecVsMueg.pdf}
%}
\caption
[Comparison of \mueconv and \muegamma]{
Searches for \mueconv and \muegamma have relative sensitivities that depend on the underlying physics, making the two channels highly complementary.
As shown on the left, New Physics can produce a signal in both channels, but one channel or the other can be comparatively suppressed due to the need to include extra vertices and loops.
The plot on the right is adapted from~\cite{ProposalPhase1}, based on~\cite{DEGOUVEA2009303}, and shows the relative sensitivity for the toy lagrangian of equation~\eq{theory:toy} as a function of $\kappa$, how non-photonic the New Physics is, and $\Lambda$, the mass scale assuming coupling strengths of unity.
\figlabel{theory:muecVsMueg}}
%\footnote{though the author has failed to reproduce the stereoscopic effect with his own eyes}
\end{figure}
}


\newcommand{\FigMuecCreation}{
\begin{figure}[bt]
\centering 
%\fbox{
%\input{figs/feynman/mu_to_e_gamma_via_SM-Wgamma.tex}
\subfloat[][\figlabel{muec:underlying}Underlying Process]{%
\includegraphics[width=0.32\textwidth,trim=-0.6cm 0 -0.6cm 0,clip]{figs/feynman/pdfs/mu_e_conversion.pdf}}\hspace{0.01\textwidth}%
\subfloat[][\figlabel{muec:atomSketch}Conversion from Ground State]{%
\includegraphics[width=0.30\textwidth]{figs/mueconv/MuEConversion-atom-sketch.pdf}}\hspace{0.01\textwidth}%
\subfloat[][\figlabel{muec:beamOnTgt}Muon Beam Stopped in Target]{%
\includegraphics[width=0.32\textwidth,trim=2.5cm 0.5cm 2.5cm 0.5cm,clip]{figs/mueconv/BeamOnTarget.png}}
%}
\caption{\figlabel{muec:creation}
The underlying \mueconv process \protect\subref{fig:muec:underlying} occurs from the ground state of a muonic atom~\protect\subref{fig:muec:atomSketch}.
To produce the muonic atoms a beam of muons has to be stopped in a target~\protect\subref{fig:muec:beamOnTgt}.
}
%\footnote{though the author has failed to reproduce the stereoscopic effect with his own eyes}
\end{figure}
}


\chapter{\mueconv and the Muonic Atom}
\sectlabel{theory:atomicMuon}
Muon-to-electron conversion is the spontaneous decay of  muon to an electron within the Coulomb potential of an atomic nucleus and without the emission of neutrinos.
It is given by the formula:
\begin{equation}
\mu^{-}+N(A,Z) \rightarrow e^{-}+N(A,Z)
\end{equation}

\FigMuecCreation
In general the nucleus involved can be excited under \mueconv, although all experimental searches to date have additionally required that the nucleus be left unchanged.
This constraint has two effects: firstly, coherent terms in the \mueconv cross section dominate since the interaction will largely be with the whole nucleus.
Being coherent, the rate of \mueconv will in general grow more quickly as a function of the atomic mass or number (which one of these itself is model dependent).
Secondly, the constraint of an unchanged nucleus means that all the free energy of the initial muon has to go into the kinetic energies of the electron and the nuclear recoil.
Since the initial system is at rest, the fact this is a two body decay fixes the energy of the outgoing electron:
\begin{equation}
E_e=M_\mu-E_{\mu,\mathrm{binding}}-E_\mathrm{recoil}
\end{equation}
where $M_\mu=$105.66~MeV/c$^2$ is the muon mass, $E_{\mu,\mathrm{binding}}$ the
binding energy of the muon in the ground state of the muonic atom, and
$E_\mathrm{recoil}$ is the kinetic energy of the recoiling nucleus.
In the aluminium target used for COMET (see section \sect{stop-tgt}) the electron energy is $E_e=104.97$~MeV.
The simplicity and model independence of the signal -- a single, monoenergetic electron -- makes the process experimentally very attractive.

The underlying physics takes place in the interaction between the muon in the ground state atomic orbital and the atomic nucleus, as illustrated in \fig{muec:creation}.
To produce the muonic atoms a beam of negative muons is brought to stop in a target, which would produce electrons that are then detected.
When negative muons in a material reach energies of a few keV of less, they become atomically captured around the nucleus of the target.
From here, on the order of 100~fs, these muons will undergo Auger and radiative transitions to the atomic ground-state.
The X-rays emitted during this electromagnetic cascade have well defined energies and intensities and can be detected as a means to evaluate the number of muons stopped in the target.
\Fig{muec:alXrays} shows the X-ray spectrum for muonic aluminium.
\CHECK{Add muonic Xrays plots}

From the ground state, in addition to the anticipated \mueconv process, there are two \ac{SM} processes that can occur to the bound muon:
\ac{DIO} and nuclear capture.
\ac{DIO} is the normal decay of a muon, during which process two neutrinos are emitted, although the spectrum of the emitted electron is modified compared to the free muon decay due to the presence of the nucleus.
Nuclear capture of the muon is the process where the muon is absorbed into the nucleus in analogue to electron nuclear capture and inverse beta decay.
A single muon-neutrino is emitted as well as various other possible particles, since the daughter nucleus is often unstable.
Both of these are important in \mueconv searches since they impose various experimental constraints and will be discussed more shortly.

These two processes determine the lifetime of the bound muon, which is not the same as the free muon.
In the case of decay, being bound to the nucleus reduces the available energy, therefore reducing the available phase-space for decay. 
In addition, a time-dilation effect occurs because the muon is never truly at rest. 
As a result the lifetime due to muon decay increases in the bound muon system compared to the free muon, and this increase grows with the atomic number, as the muon is bound tighter and tighter to the nucleus.
However, whilst the rate of decay decreases with atomic number the rate of muon capture increases.
This occurs firstly because there are more protons against which to capture, and secondly because the muon wavefunction overlaps more and more with the nucleus.
For atomic numbers larger than $Z=30$ the muon wavefunction is contained almost completely within the nucleus.\CHECK{Calculate this? Is 30 a reasonable value of Z for this statement?}
Whilst for light elements, up to around $Z=12$, the decay process dominates, for the rest of the periodic table the capture process determines the muon lifetime.
For an aluminium target, the two processes are roughly equal, with branching ratios of 61:39 for capture to decay, and a muon lifetime of 864~ns~\cite{Measday2007Comparison}.
\CHECK{Do I want the lifetime plot in here, or is it better left in the detector section?}

Since the muon is 200~times heavier than the electron, the muon wavefunction feels the effect of the nucleus a lot more.
Rather than the full branching ratio, typically \mueconv experiments discuss the conversion rate, which is given by:
\begin{equation}
\mathcal{C.R.}=\frac{\Gamma\left(\mathrm{\mueconv}\right)}{\Gamma\left(\mathrm{nuclear~capture}\right)}
\end{equation}
The key advantage over using the full branching ratio is that by normalising to the number of muons that undergo nuclear capture, as opposed to the total number of stopped muons, the theoretical uncertainty due to the initial muon wave-function is reduced since this is would be needed to predict the decay rate.

From this, one defines the \acf{ses} to be:
\begin{equation}
	\eqlabel{det:ses}
\mathrm{S.E.S}(\muec)=\frac{1}{N_\mu \mathcal{B}_\mathrm{capture} A_{\mu\rightarrow e}}
\end{equation}
where $N_\mu$ is the number of muons stopped, $\mathcal{B}_\mathrm{capture}$ is the branching ratio for muon nuclear capture, and $A_{\mu\rightarrow e}$ is the total acceptance of electrons coming from \mueconv.

\section{Muon Decay in Orbit}
In free muon decay the maximum energy for the outgoing electron occurs when the neutrinos recoil back-to-back with the electron.
In this configuration, exactly half the energy released in the decay is available to the electron, so that the maximum energy of an electron coming from the decay of a free muon at rest is: $\max(E_{e}^\textrm{free})=m_\mu/2=52.5$~MeV.

This situation is changed completely when the muon becomes bound to the nucleus of an atom.
Once bound, the neutrinos can be arranged back-to-back with each one another, carrying away a negligible amount of energy.
Four-momentum can still be conserved however, since the nucleus of the atom recoils against the electron.  
Given the enormous mass of any nucleus compared to the electron, momentum can be conserved with only a small amount of kinetic energy and the maximum electron energy is hugely increased compared to the free decay.
In fact, if the neutrinos take away no energy the kinematic configuration of this decay becomes identical to that of \mueconv but for the mass of the neutrinos: $\max(E_{e}^\textrm{DIO})\simeq{}E_{e}^\textrm{conversion}$.

The spectrum of electrons from \ac{DIO} in aluminium is shown in \fig{muec:DIO}.
It can be seen how the peak electron energy is close to the free muon decay end-point, and in reality about 99\% of \ac{DIO} electrons will be emitted below 55~MeV \CHECK{99\% energy for DIO electrons}.
Whilst the end-point for the spectrum is indeed around 104.97~MeV, it is clear how suppressed this part of the spectrum is -- some twenty orders of magnitude less likely than at the peak energy.
Achieving the end-point energy requires radiative connections between the nucleus and the muon or electron; the low neutrino momentum brings about a helicity suppression; and the small amount of phase-space available to produce low-energy electrons further suppresses things.
\CHECK{Add DIO spectrum plot and label the free muon decay and the \mueconv end point}

For these reasons, \mueconv searches historically described themselves as `background free'.
However, given the projected sensitivities of modern experiments, the \ac{DIO} rate close to the end-point of the spectrum is now at an appreciable level.
Indeed, the next generation of searches (and COMET \phaseI in particular) will be the first to measure the \ac{DIO} spectrum above 90~MeV, which will form an important check for the theoretical prediction of muon \ac{DIO}.

\section{Muon Nuclear Capture}
%\begin{easylist}
%# Inverse beta-decay
%# Prompt process of muon-proton -> neutron and neutrino
%# Nucleon clustering means prompt protons also observed (muon-nucleon cluster -> neutrino-neutron-proton)
%# about 50~MeV excitation of the nucleus
%# Emission of various particles during nuclear de-excitation: protons, neutrons, gammas, deuterons, triton, alphas
%# Difficult to predict products and rates theoretically
%# Interest in this process towards the end of the 70s as a means to test nuclear theory but since then interest has waned
%# Lack of experimental measurements for Al($\mu$,X)Mg so needed to measure things: \alcap
%# See appendix for an overview of \alcap
%# Proton emission around 3.5\% of every muon captured
%\end{easylist}
The nuclear capture of negative muons is governed by the equation:
\begin{equation}
\mu^-+N(A,Z)\rightarrow \nu_\mu+N'(A,Z-1)
\end{equation}
Whilst it is clearly an incoherent process, the direct process can occur directly between a muon and proton, resulting in a prompt neutron, or between a cluster of nucleons, which can cause both prompt neutrons and protons to be produced.
The nuclear excitation after such a process is typically around 50~MeV, with the other half of the total incoming energy lost to the outgoing neutrino.
Whilst both prompt neutrons and protons are possible, the remnant nucleus is often left in an excited, unstable state, such that during de-excitation other particles can also be emitted.
These include neutrons and protons but also gammas, deuterons, triton and alpha particles.

From the perspective of a sensitive \mueconv experiment the emission products following nuclear capture can be dangerous, since, in the case of charged particles, they can swamp the detector if left unchecked.
Similarly, neutrons and gamma rays produced by nuclear capture can become dangerous for electronics systems if left unchecked.
As such it is important to understand the rates of these particles emitted after nuclear capture of the muon.

However, theoretical predictions of the rates and energy distributions of capture products is extremely complex and experimental measurements are necessary.
Unfortunately, in the case of aluminium, the target choice for \COMET, the existing experimental data is not extensive.
\Fig{muec:capture:data} shows a summary of the available information.  
Accordingly it has been necessary to measure this directly.
The \alcap experiment~\cite{AlcapProposal2012} is a joint effort between COMET and Mu2e tasked with measuring the rate and spectra of particles emitted following muon capture in aluminium.
Three runs have been held at \ac{PSI} from 2013 to 2015, and data analysis is on-going, although preliminary neutron and proton spectra and rates have been achieved.
The measured proton rate is low enough that \COMET does not expect to have to take any precautionary measures to reduce it further.
For more information on \alcap, see appendix~\sect{appendix:alcap} and the PhD thesis by Nam Tran~\cite{NamThesis}.
\CHECK{Add figure summarizing the capture data}

\newcommand{\FigBentSolenoidRelativeDrift}{
\begin{figure}[t]
\centering
\includegraphics[width=0.6\textwidth]{figs/detector/BentSolenoids_RelativeDrift}
\caption{
Angular dependence of the magnitude of vertical drift in a bent solenoid field.
The total variation (black) remains below 10\% for pitch angles below 50\degree.
}
\figlabel{detector:bent-solenoids:angularDependence}
\end{figure}
}

\newcommand{\FigPhaseII}{
\begin{figure}[t]
\centering
\includegraphics[width=0.9\textwidth]{figs/detector/PhaseII_schematic}
\caption{
Schematic layout of COMET \phaseII. 
The 8 GeV proton beam enters from the top-left, producing (amongst other things) pions.
Pions and muons travelling backwards with respect to the proton beam are then transported around 180 degrees of bent solenoid, during which time most of the pions decay producing an intense muon beam.
About 40\% of these muons then stop in the stopping target (centre of image).
Any electrons coming from  \mueconv are then transported through another 180 degrees of bent solenoid into the detector system.
}
\figlabel{detector:PhaseII:setup}
\end{figure}
}

\newcommand{\FigPhaseI}{
\begin{figure}[t]
\centering
\includegraphics[height=0.4\textheight]{figs/detector/PhaseI_schematic}
\caption{
Schematic layout of COMET \phaseI. 
}
\figlabel{detector:PhaseI:setup}
\end{figure}
}

\newcommand{\TabBackgroundSummary}{
\begin{table}
\begin{tabular}{lldd}
     \hline
     \hline\\[-1.8ex]
     Type           & Background & \multicolumn{2}{c}{Predicted number of events during run} \\
                    &  & \multicolumn{1}{c}{\phaseI \cite{TDR2014}} & \multicolumn{1}{c}{\phaseII \cite{CDRphase2} } \\
     \hline\\[-1.8ex]
     Intrinsic & Muon Decay-in-Orbit                       & 0.01              & 0.15    \\
               & Radiative Muon Capture                    & 0.00056           & <0.001  \\
               & $\mu^-$ Capture w/ n Emission             & <0.001            & <0.001  \\
               & $\mu^-$ Capture w/ Charged Part. Emission & <0.001            & <0.001  \\
     Prompt    & Radiative Pion Capture                    & 0.00023           & 0.05    \\
               & Beam Electrons                            & 0.00083           & <0.1^*  \\
               & Muon Decay in Flight                      & \le0.0002         & <0.0002 \\
               & Pion Decay in Flight                      & \le0.00023        & <0.0001 \\
               & Neutron Induced                           & -                 & 0.024   \\
               & Other beam induced B.G.                   & <2.8\times10^{-6} & -       \\
     Delayed   & Delayed Radiative Pion Capture            & \sim0             & 0.002   \\
               & Anti-proton Induced                        & 0.007             & 0.007   \\
               & Other delayed B.G.                        & \sim0             & -       \\
     Cosmic    & Cosmic Ray Muons                          & -                 & 0.002   \\
               & Electrons from Cosmic Ray Muons           & <0.0001           & 0.002   \\
     \hline\\[-1.8ex]
     \multicolumn{2}{c}{Total background}                      & 0.019         & 0.34    \\
     \multicolumn{2}{c}{Signal (Assuming $B=1\times10^{-16}$)} & 0.31          & 3.8     \\
     \hline
     \hline
\end{tabular}
\caption{
	\CHECK{UPDATE \phaseI values with TDR 2016}
	Backgrounds for COMET \phaseI \cite{TDR2014} and \phaseII \cite{CDRphase2}.
	Prompt backgrounds arise by protons that occur in between bunches and are therefore suppressed by the extinction factor.
	For \phaseI, the recently measured value of $10^{-12}$ was used for the extinction factor, but for \phaseII the older expectation of $10^{-9}$ was used.
}
\tablabel{detector:backgrounds}
\end{table}
}

\newcommand{\FigMuonNuclearParams}{
\begin{figure}[bp]
\centering
\subfloat[][\figlabel{detector:mu-nucl-params:lifetimes}Lifetimes]{\includegraphics[width=0.9\textwidth]{figs/detector/MuNuclearParams_All_lifetimes.pdf}}\\
\subfloat[][\figlabel{detector:mu-nucl-params:end-point}End-point Shift]{\includegraphics[width=0.49\textwidth]{figs/detector/MuNuclearParams_energies.pdf}}%\hspace{0.5cm}%
\subfloat[][\figlabel{detector:mu-nucl-params:branching-ratio}Branching Fraction]{\includegraphics[width=0.49\textwidth]{figs/detector/MuNuclearParams_branching_fraction.pdf}}
\caption{\figlabel{detector:mu-nucl-params}
The effect of changing the atomic number on the branching ratio, lifetime and electron energy spectrum end-point.
For the branching ratio and lifetime plots, the partial rate for muon nuclear capture and decay-in-orbit are shown separately.
The capture and decay rates are taken from the Geant4 \cite{Geant42003} parametrisation for stopped negative muons.  
Only elements for which at least 1 isotope uses a measured value are plotted.
The values for the end-point energy level are calculated using the Bohr model for the muon ground-state binding energy.
}
\end{figure}
}

\newcommand{\FigICEDUSTOverview}{
\begin{figure}[t]
\centering
%\fbox{
\includegraphics[width=1.00\textwidth,trim=0.85cm 0.5cm 0.5cm 0.5cm,clip=true]{figs/software/ICEDUST_structure}
%}
\caption{
Overview diagram for the ICEDUST framework.
Data produced from simulation or taken in the real experiment are treated identically through the calibration and onwards up to analysis.
}
\figlabel{software:ICEDUSTOverview}
\end{figure}
}

\newcommand{\FigNDTwoEighty}{
\begin{figure}[t]
\centering
\includegraphics[width=0.95\textwidth]{figs/software/ND280SoftwareDiagram}
\caption{
Overview diagram for the ND280 framework.
}
\figlabel{software:ND280}
\end{figure}
}

\newcommand{\FigSimulationOverview}{
\begin{figure}[t]
\centering
%\fbox{
\includegraphics[width=1.00\textwidth]{figs/software/Simulation_structure}
%}
\caption{
Diagram showing the stages used to simulate COMET.
The timing schematics on the right show how a simulated event is built up, firstly by producing many individual proton interactions with the production target,
then by transporting the secondary particles to produce energy deposits in the detector, which are then combined with the truth hits from other proton events to produce a realistic bunch structure.
Finally, these bunch events are processed through the detector response simulation to produce fake waveforms and other detector read-outs.
}
\figlabel{software:SimulationOverview}
\end{figure}
}

\newcommand{\FigGeometryHeirarchy}{
\begin{figure}[t]
\centering
%\fbox{
\includegraphics[width=0.70\textwidth]{figs/software/ComponentHeirarchy}
%}
\caption{
How parameters are shared amongst different components.
Parameters of component 1.a (in red) can access the value of parameters owned by components contained in the larger violet region.
}
\figlabel{software:componentHeirarchy}
\end{figure}
}

\newcommand{\FigGeometryParameters}{
\begin{figure}[t]
\lstinputlisting[style=customc]{figs/software/demo-parameters.mac}
\caption{
An example set of parameter definitions which control the geometry for the Torus2.
Parameter specifications use natural arithmetic notation and can reference other parameters and use standard units.
They can also be formed as sets where each element has a different value, such as the \texttt{Coils:Position} parameter.
}
\figlabel{software:geom:paramAssignments}
\end{figure}
}

\newcommand{\FigGeometryScreenshots}{
\begin{figure}[tb]
        \subfloat[][\figlabel{software:geom:screenshots:phaseI}\phaseI]  {\includegraphics[height=0.25\textheight]{figs/software/Phase-I-UpdateGeom}}\hspace{2ex}%
        \subfloat[][\figlabel{software:geom:screenshots:phaseII}\phaseII]{\includegraphics[height=0.25\textheight,trim=0 0.3cm 0 2.6cm,clip=true]{figs/software/Phase-II-UpdateGeom.png}}
\caption{\figlabel{software:geom:screenshots} %
Two of the possible simulation `worlds' that can be selected at run-time:
        \protect\subref{fig:software:geom:screenshots:phaseI} \phaseI with the CyDet detector installed, and
        \protect\subref{fig:software:geom:screenshots:phaseII} \phaseII.  
	Mutiple \phaseI worlds exist, one for each potential running configuration.
}
\end{figure}
}

\newcommand{\FigPiYieldHadronCodes}{
\begin{figure}[b]
\centering
%\fbox{
\includegraphics[width=0.95\textwidth,trim=2cm 1cm 0.8cm 0.5cm,clip=true]{figs/software/PionYield_AEdmondsThesis}
%}
\caption{
Comparison of various hadron production codes with experimental data from the HARP experiment, taken from the thesis of A. Edmonds~\cite{AEdmondsThesis}.
Points with error bars are the experimental data.  Left: double differential-production cross-section for pion production from 20 to 32\degree with respect to the incoming proton direction; right: from 100 to 112\degree.
The hadron production code that best reproduces the data depends strongly on the angular region under consideration.
}
\figlabel{software:piYield}
\end{figure}
}

\newcommand{\FigSoftwarePhysicsSpectra}{
\begin{figure}[p]
\centering
%\fbox{
\subfloat[][\figlabel{software:customPhysic:DIO}Electrons from $\mu$ Decay-in-Orbit]{
\includegraphics[width=0.45\textwidth,trim=0cm 0cm 0.0cm 1.3cm,clip=true]{figs/software/160822_BoundDecay_Geant4_vs_Czarnecki-lin.pdf}
\includegraphics[width=0.45\textwidth,trim=0cm 0cm 0.0cm 1.3cm,clip=true]{figs/software/160822_BoundDecay_Geant4_vs_Czarnecki-log.pdf}
}\\
\subfloat[][\figlabel{software:customPhysic:ProtMuCap}Protons Emitted Following $\mu$ Nuclear Capture]{
\includegraphics[width=0.45\textwidth,trim=0cm 0cm 1.8cm 1.9cm,clip=true]{figs/software/160822_Geant4VsAlcap-lin.pdf}
\includegraphics[width=0.45\textwidth,trim=0cm 0cm 1.8cm 1.9cm,clip=true]{figs/software/160822_Geant4VsAlcap-log.pdf}
}
%}
\caption{
\figlabel{software:customPhysic}
Comparison of the realistic spectra for \ac{DIO} electrons, \protect\subref{fig:software:customPhysic:DIO} (normalised to agree at 35~MeV), and protons coming from muon nuclear capture, \protect\subref{fig:software:customPhysic:ProtMuCap} (normalised to have the same maximum value), each on a linear scale (left) and a logarithmic scale (right).
The \ac{DIO} spectrum used in default Geant4 has a sharp cut-off slightly above the free muon decay end-point, to be compared with the long but steeply falling tail of the Czarnecki \etal theoretical calculation~\cite{Czarnecki2011}.
The comparison of protons coming from muon capture between the preliminary result from AlCap and default Geant4 shows that the true proton spectrum is much softer than the Geant4 model.
}
\end{figure}
}

\newcommand{\FigSimulationPhysicsClasses}{
\begin{figure}[tb]
\centering
%\fbox{
\includegraphics[width=1.00\textwidth]{figs/software/SimulationMuonPhysicsClasses}
%}
\caption{
The various classes involved in simulating the various processes of stopped negative muons.
%Classes in red have been implemented for COMET and augment the existing Geant4 classes which are shown in blue.
The standard Geant4 model is activated by registering `G4MuonMinusCapture', which instantiates `G4MuonMinusBoundDecay' and `G4CascadeInterface' to run the \ac{DIO} and nuclear capture respectively.
To use the custom COMET muon physics, an instance of `COMETStoppedMuonProcess' should be registered, which sets up `COMETMuonMinusBoundDecayOrConversion' to produce the electron (and possibly neutrinos) from \ac{DIO} or conversion, and `COMETMuonMinusCapture' to do the nuclear capture.
}
\figlabel{software:ExtendedMuonClasses}
\end{figure}
}

\newcommand{\FigSoftwareFieldMap}{
\begin{figure}[b]
\centering
%\fbox{
\subfloat[][\figlabel{software:field:Opera}Opera calculation]{
\includegraphics[width=0.45\textwidth,trim=0cm 0cm 0.0cm 13cm,clip=true]{figs/software/Plot_Opera.png}
}
\subfloat[][\figlabel{software:field:G4Beamline}G4Beamline calculation]{
\includegraphics[width=0.45\textwidth,trim=0cm 0cm 0.0cm 13cm,clip=true]{figs/software/Plot_G4Beamline.png}
}
%}
\caption{
\figlabel{software:field}
Fieldmap produced by \protect\subref{fig:software:field:Opera} Opera and \protect\subref{fig:software:field:G4Beamline} G4Beamline.
Although the fringe field is larger with the G4Beamline calculation, the lack of material effects make this calculation less reliable.
Note that the G4Beamline calculation does not include the detector solenoid.
}
\end{figure}
}

\newcommand{\FigSoftwareFieldMapComparison}{
\begin{figure}[t]
\centering
%\fbox{
\includegraphics[width=0.9\textwidth,trim=0cm 0cm 0.0cm 13cm,clip=true]{figs/software/Plot_ratio_opera-G4Beamline.png}
%}
\caption{
\figlabel{software:field:comparison}
The ratio of the Opera and G4Beamline calculations shown in \fig{software:field}.
For most of the field within the beamline the calculations agree within 10\%, although around the ends of the solenoids the agreement is poorer.
}
\end{figure}
}

\newcommand{\FigSoftwareDipoleField}{
\begin{figure}[t]
\centering
%\fbox{
\includegraphics[width=0.9\textwidth,trim=2cm 8cm 0.0cm 6cm,clip=true]{figs/software/DipoleFields.png}
%}
\caption{
\figlabel{software:field:dipole}
The dipole field calculations used in ICEDUST for \phaseII. 
	The first 180\degree of the muon transport beamline have a positivie dipole field (pointing out of the plane)
	whilst the dipole along the electron spectrometer is negative and with a much larger strength.
	It can also be seen that the calculation for the dipole along the muon transport beamline contains realistic features (fringe fields, non-uniformities, etc)
	whilst the dipole field for the electron spectrometer is artificially uniform.  
}
\end{figure}
}

%\include{muon-beam}
%\include{pile-up}
%\include{alcap}
\newcommand{\TabOptimisationParameters}{%
\begin{table}[tb]%
\begin{tabular}{ll}%
	\hline
	Region for optimisation & Approx.\ No.\  of parameters \\
	\hline
	Production target dimensions and location & $3+3$ \\
	Torus1 dipole field strength & 1 \\
	Torus2 dipole field strength & 1 \\
	Muon beam collimator shapes, position, and material & $3+1+1$ \\
	Stopping target shape and location & $4+3$ \\
	Beam blocker position, form, and material& $3+3+1$ \\
	Electron spectrometer dipole field strength& 1 \\
	DIO blockers in the spectrometer & $4$ \\
	\hline
	Approx. total number of parameters& 32 \\
	\hline
\end{tabular}%
	\caption{\tablabel{optimisation:possible-parameters}%
	Aspects of the experiment that can be optimised and estimates for the number of parameters that define each aspect.
	In the case of the target, beam blocker, and collimator shapes the number of parameters is only approximate; crudely speaking there is at least a width, length and height but in principle one could have a very irregular shape that cannot be parametrised by only three numbers, for example, shapes that change as a function of distance along the beamline.
}%
\end{table}%
}

\newcommand{\FigOptimProdTgtLength}{
\begin{figure}[pt]
\centering
\subfloat[][\figlabel{optimisation:ProdTgtSec:Length:Momentum:Muons}Muons]{
\includegraphics[width=0.48\textwidth,trim=0 0 0 1.5cm,clip]{figs/optimisation/ProdTgtGeom/Length_mu-minus_momentum}}
\subfloat[][\figlabel{optimisation:ProdTgtSec:Length:Momentum:Pions}Pions]{
\includegraphics[width=0.48\textwidth,trim=0 0 0 1.5cm,clip]{figs/optimisation/ProdTgtGeom/Length_pi-minus_momentum}}
\caption{\figlabel{optimisation:ProdTgtSec:Length:Momentum}
Change to momentum distributions at the entrance to the first 90 degrees of the bent muon beam solenoid for different target lengths.
Target length is given as half-length which is the Geant4 convention.  
}
\end{figure}
\begin{figure}[pt]
\centering
\subfloat[][\figlabel{optimisation:ProdTgtSec:Length:Integral:Muons}Muons]{
\includegraphics[width=0.48\textwidth,trim=0 0 2cm 2.8cm,clip]{figs/optimisation/ProdTgtGeom/Length_mu-minus_integral_toZero}}
\subfloat[][\figlabel{optimisation:ProdTgtSec:Length:Integral:Pions}Pions]{
\includegraphics[width=0.48\textwidth,trim=0 0 2cm 2.8cm,clip]{figs/optimisation/ProdTgtGeom/Length_pi-minus_integral_toZero}}
\caption{\figlabel{optimisation:ProdTgtSec:Length:Integral}
Integrated muon and pion yields up to a certain momentum at the entrance to the first 90 degrees of the bent muon beam solenoid as a function of target length.
}
\end{figure}
\begin{figure}[pt]
\centering
\subfloat[][\figlabel{optimisation:ProdTgtSec:Length:IntegralRatio:Muons}Muons]{\includegraphics[width=0.48\textwidth,trim=0 0 2cm 2.8cm,clip]{figs/optimisation/ProdTgtGeom/Length_mu-minus_integral_ratios}}
\subfloat[][\figlabel{optimisation:ProdTgtSec:Length:IntegralRatio:Pions}Pions]{\includegraphics[width=0.48\textwidth,trim=0 0 2cm 2.8cm,clip]{figs/optimisation/ProdTgtGeom/Length_pi-minus_integral_ratios}}
\caption{\figlabel{optimisation:ProdTgtSec:Length:IntegralRatio}
Change in the momentum distribution of muons and pions at the entrance to the first 90 degrees of the bent muon beam solenoid as a function of target length.
}
\end{figure}
}

\newcommand{\FigOptimProdTgtRad}{
\begin{figure}[pt]
\centering
\subfloat[][\figlabel{optimisation:ProdTgtSec:Radius:Momentum:Muons}Muons]{\includegraphics[width=0.48\textwidth,trim=0 0 0 1.5cm,clip]{figs/optimisation/ProdTgtGeom/Radius_mu-minus_momentum}}
\subfloat[][\figlabel{optimisation:ProdTgtSec:Radius:Momentum:Pions}Pions]{\includegraphics[width=0.48\textwidth,trim=0 0 0 1.5cm,clip]{figs/optimisation/ProdTgtGeom/Radius_pi-minus_momentum}}
\caption{
\figlabel{optimisation:ProdTgtSec:Radius:Momentum}
Change to momentum distributions at the entrance to the first 90 degrees of the bent muon beam solenoid for different target radii.
}
\end{figure}
\begin{figure}[pt]
\centering
\subfloat[][\figlabel{optimisation:ProdTgtSec:Radius:Integral:Muons}Muons]{\includegraphics[width=0.48\textwidth,trim=0 0 2cm 2.8cm,clip]{figs/optimisation/ProdTgtGeom/Radius_mu-minus_integral_toZero}}
\subfloat[][\figlabel{optimisation:ProdTgtSec:Radius:Integral:Pions}Pions]{\includegraphics[width=0.48\textwidth,trim=0 0 2cm 2.8cm,clip]{figs/optimisation/ProdTgtGeom/Radius_pi-minus_integral_toZero}}
\caption{\figlabel{optimisation:ProdTgtSec:Radius:Integral}
Integrated muon and pion yields up to a certain momentum at the entrance to the first 90 degrees of the bent muon beam solenoid as a function of target radius.
}
\end{figure}
\begin{figure}[pt]
\centering
\subfloat[][\figlabel{optimisation:ProdTgtSec:Radius:IntegralRatio:Muons}Muons]{\includegraphics[width=0.48\textwidth,trim=0 0 2cm 2.8cm,clip]{figs/optimisation/ProdTgtGeom/Radius_mu-minus_integral_ratios}}
\subfloat[][\figlabel{optimisation:ProdTgtSec:Radius:IntegralRatio:Pions}Pions]{\includegraphics[width=0.48\textwidth,trim=0 0 2cm 2.8cm,clip]{figs/optimisation/ProdTgtGeom/Radius_pi-minus_integral_ratios}}
\caption{\figlabel{optimisation:ProdTgtSec:Radius:IntegralRatio}
Change in the momentum distribution of muons and pions at the entrance to the first 90 degrees of the bent muon beam solenoid as a function of target radius.
}
\end{figure}
}

\newcommand{\FigOptimProdTgtFinal}{
\begin{figure}[pt]
\centering
\subfloat[][\figlabel{optimisation:ProdTgtSec:Final:Integral:Muons}Muons]{\includegraphics[width=0.48\textwidth,trim=0 0 0 1.5cm,clip]{figs/optimisation/ProdTgtGeom/OptimalLengthRadius_mu-minus_integral_toZero.pdf}}
\subfloat[][\figlabel{optimisation:ProdTgtSec:Final:Integral:Pions}Pions]{\includegraphics[width=0.48\textwidth,trim=0 0 0 1.5cm,clip]{figs/optimisation/ProdTgtGeom/OptimalLengthRadius_pi-minus_integral_toZero.pdf}}
\caption{\figlabel{optimisation:ProdTgtSec:Final:Integral}
Variation in muon and pion yields as a function of target radius when the total target length is set to the optimised value of 32~cm.
Despite the longer target length the optimal radius is still 1~cm.
}
\end{figure}
}

\newcommand{\FigOptimProdTgtComparePhases}{
\begin{figure}[pt]
\centering
\includegraphics[width=0.9\textwidth,trim=0 0 0 1.5cm,clip]{figs/optimisation/ProdTgtGeom/Plot_compare_phase_1and2.pdf}
\caption{\figlabel{optimisation:ProdTgtSec:Phase1vs2}
Comparison of the muon and pion yields per POT for \phaseI and \phaseII.
The difference arises from the change of target material between the phases.
}
\end{figure}
}
\newcommand{\FigOptimMuBeamDipoleMuStops}{
\begin{figure}[bt]
\centering
%	\fbox{
\includegraphics[width=0.85\textwidth,trim=0 0.5cm 0 1.0cm,clip]{figs/optimisation/MuonBeamDipoles/Tidied_stopped_muons.pdf}
%}
\caption{\figlabel{optim:muBeamDipole:stoppedMu}
	Muon stopping rate as a function of the two dipole field strengths (given relative to the \phaseI design specification).
	A clear anti-correlation is visible which is discussed in the text.
}
\end{figure}
}

\newcommand{\FigOptimMuBeamDipolePiStops}{
\begin{figure}[t]
\centering
%	\fbox{
\includegraphics[width=0.85\textwidth,trim=0 0.5cm 0 1.0cm,clip]{figs/optimisation/MuonBeamDipoles/Tidied_stopped_pions.pdf}
%}
\caption{\figlabel{optim:muBeamDipole:stoppedPi}
	Pion stopping rate as a function of the two dipole field strengths (given relative to the \phaseI design specification).
At the level of statistics used to generate each point, no clear trend is obvious.
Empty squares are those where no pions stopped in the run.
}
\end{figure}
}

\newcommand{\FigOptimMuBeamDipoleMuDispersion}{
\begin{figure}[t]
\centering 
\subfloat[][\figlabel{optim:MuBeamDipole:MuDispersion:Entry}Torus1 Entry]{\includegraphics[width=0.45\textwidth,trim=0 0.9cm 0 1.9cm,clip]{figs/optimisation/MuonBeamCollimators/Tidied_Muon_dispersion_at_entrance.pdf}}
\subfloat[][\figlabel{optim:MuBeamDipole:MuDispersion:Exit}Torus2 Exit]  {\includegraphics[width=0.45\textwidth,trim=0 0.9cm 0 1.9cm,clip]{figs/optimisation/MuonBeamCollimators/Tidied_Muon_dispersion_at_exit.pdf}}
\caption{\figlabel{optim:MuBeamDipole:MuDispersion}
Dispersive effect of the 180\degree bent transport solenoid and dipole field on muons.
No collimating material is yet included, so the high-energy muons being removed is due purely to the beam-pipe itself.
}
\end{figure}
}

\newcommand{\FigOptimMuBeamCollimMuonPaths}{
\begin{figure}[p]
\centering 
	\subfloat[][\figlabel{optim:MuBeamCollim:Beamline:All}All Muons]          {\includegraphics[width=1\textwidth,trim=18cm 1.0cm 26cm 1cm,clip]{figs/optimisation/MuonBeamCollimators/Tidied_All_muons_wGeom.png}}\\
\subfloat[][\figlabel{optim:MuBeamCollim:Beamline:Stopped}Stopped Muons]          {\includegraphics[width=1\textwidth,trim=18cm 1.0cm 26cm 1cm,clip]{figs/optimisation/MuonBeamCollimators/Tidied_Stopped_muons_wGeom.png}}\\
\subfloat[][\figlabel{optim:MuBeamCollim:Beamline:HighP}Muons with $p>70$ MeV/c around the stopping target]          {\includegraphics[width=1\textwidth,trim=18cm 1.0cm 26cm 1cm,clip]{figs/optimisation/MuonBeamCollimators/Tidied_HighP_muons_wGeom.png}}\\
\subfloat[][\figlabel{optim:MuBeamCollim:Beamline:Diff}Stopped $\mu-$High-$p$ $\mu$]{\includegraphics[width=1\textwidth,trim=18cm 1.0cm 26cm 1cm,clip]{figs/optimisation/MuonBeamCollimators/Tidied_Where_to_collimate_wGeom.png}}
\caption{\figlabel{optim:MuBeamCollim:Beamline}
The heights of muons as they pass along the beamline.  
	\protect\subref{fig:optim:MuBeamCollim:Beamline:All} The path of all muons.
	\protect\subref{fig:optim:MuBeamCollim:Beamline:Stopped}: The paths of muons that stop in the target.
	\protect\subref{fig:optim:MuBeamCollim:Beamline:HighP}: The heights of muons with momentum greater than 70 MeV/c when they enter the region around the stopping target.  These could potentially decay in flight to give electrons with 100 MeV/c or greater.
	\protect\subref{fig:optim:MuBeamCollim:Beamline:Diff}: The difference between plot \protect\subref{fig:optim:MuBeamCollim:Beamline:Stopped} and plot \protect\subref{fig:optim:MuBeamCollim:Beamline:HighP}.
	Regions in dark blue would give the greatest impact in removing high-momentum muons whilst leave the stopping muons untouched.
	These plots should be compared to those of \fig{optim:MuBeamCollim:BeamWColl} once collimators have been introduced.
}
\end{figure}
}

\newcommand{\FigOptimMuBeamCollimTransverseSep}{
\begin{figure}[t]
\centering 
\subfloat[][\figlabel{optim:MuBeamCollim:TransverseSep:TS1Entry}At 0\degree]  {\includegraphics[height=0.35\textheight,trim=0.0cm 0.8cm 1.3cm 1.9cm,clip]{figs/optimisation/MuonBeamCollimators/MuonTransversePos_Torus1Tor1FirstPoint}}
\subfloat[][\figlabel{optim:MuBeamCollim:TransverseSep:TS3}At 90\degree] {\includegraphics[height=0.35\textheight,trim=1.7cm 0.8cm 1.3cm 1.9cm,clip]{figs/optimisation/MuonBeamCollimators/MuonTransversePos_Torus2Monitor_0}}
\subfloat[][\figlabel{optim:MuBeamCollim:TransverseSep:TS2Exit}At 180\degree]{\includegraphics[height=0.35\textheight,trim=1.7cm 0.8cm 1.3cm 1.9cm,clip]{figs/optimisation/MuonBeamCollimators/MuonTransversePos_StopTgtSecMonitor_0}}
\caption{\figlabel{optim:MuBeamCollim:TransverseSep}
The separation between stopping and dangerous muons.
The separation is largest at the exit (180\degree), reasonable at the entrance (0\degree), and smallest around the mid-point (90\degree).
}
\end{figure}
}

\newcommand{\FigOptimMuBeamCollimMuonPathsWColl}{
\begin{figure}[ph]
\centering 
\subfloat[][\figlabel{optim:MuBeamCollim:BeamWColl:All}All Muons]                  {\includegraphics[width=1\textwidth,trim=18cm 1.0cm 26cm 1cm,clip]{figs/optimisation/MuonBeamCollimators/Tidied_WColl_AllMuons_WGeom.png}}\\
\subfloat[][\figlabel{optim:MuBeamCollim:BeamWColl:Stopped}Stopped Muons]          {\includegraphics[width=1\textwidth,trim=18cm 1.0cm 26cm 1cm,clip]{figs/optimisation/MuonBeamCollimators/Tidied_WColl_StoppedMuons_WGeom.png}}\\
\subfloat[][\figlabel{optim:MuBeamCollim:BeamWColl:HighP}Muons with $p>70$ MeV/c around the stopping target]          {\includegraphics[width=1\textwidth,trim=18cm 1.0cm 26cm 1cm,clip]{figs/optimisation/MuonBeamCollimators/Tidied_WColl_HighPMuons_WGeom.png}}\\
\caption{\figlabel{optim:MuBeamCollim:BeamWColl}
The heights of muons as they pass along the beamline.  
	\protect\subref{fig:optim:MuBeamCollim:Beamline:All} The path of all muons.
	\protect\subref{fig:optim:MuBeamCollim:Beamline:Stopped}: The paths of muons that stop in the target.
	\protect\subref{fig:optim:MuBeamCollim:Beamline:HighP}: The heights of muons with momentum greater than 70 MeV/c when they enter the region around the stopping target.  These could potentially decay in flight to give electrons with 100 MeV/c or greater.
	These plots should be compared to those of \fig{optim:MuBeamCollim:Beamline} before collimators were introduced, where it is clear how well the dangerous muons are being suppressed.
}
\end{figure}
}

\newcommand{\FigOptimMuBeamCollimTorusOne}{
\begin{figure}[t]
\centering 
\subfloat[][\figlabel{optim:MuBeamCollim:Torus1:perPOT}Particles Surviving per POT]{\includegraphics[width=0.5\textwidth,trim=0.8cm 0.8cm 0.6cm 1.9cm,clip]{figs/optimisation/MuonBeamCollimators/Survived_Coll1_unNormalised-log.pdf}}
\subfloat[][\figlabel{optim:MuBeamCollim:Torus1:fraction}Fraction Surviving Collimator]{\includegraphics[width=0.5\textwidth,trim=0.8cm 0.8cm 0.6cm 1.9cm,clip]{figs/optimisation/MuonBeamCollimators/Survived_Coll1_Normalised-lin.pdf}}
\caption{\figlabel{optim:MuBeamCollim:Torus1}
The effect of changing the height of the collimator in Torus1 on the particle distributions.
}
\end{figure}
}

\newcommand{\FigOptimMuBeamCollimTorusTwoPerPOT}{
\begin{figure}[t]
\centering 
\includegraphics[width=0.95\textwidth,trim=0.3cm 0.3cm 0.8cm 0.1cm,clip]{figs/optimisation/MuonBeamCollimators/Survived_Coll2_unNormalised-lin.pdf}
\caption{\figlabel{optim:MuBeamCollim:Torus2:perPOT}
The number of particles reaching the end of the Torus2 solenoid per POT for different heights of both collimators in Torus1 and Torus2.
}
\end{figure}
}


\newcommand{\FigOptimMuBeamCollimTorusTwoFraction}{
\begin{figure}[t]
\centering 
\includegraphics[width=0.8\textwidth,trim=0.3cm 0.3cm 0.8cm 0.1cm,clip]{figs/optimisation/MuonBeamCollimators/Survived_Coll2_unNormalised-lin.pdf}
\caption{\figlabel{optim:MuBeamCollim:Torus2:fraction}
	The number of particles reaching the end of the Torus2 solenoid relative to the number that enter the Torus1 solenoid (\ie the survival probability) for different heights of both collimators in Torus1 and Torus2.
}
\end{figure}
}

\newcommand{\FigOptimMuBeamCollimTorusTwoContours}{
\begin{figure}[bt]
\centering 
%\fbox{
\includegraphics[width=0.75\textwidth,trim=0.1cm 0.3cm 2.8cm 1.1cm,clip]{figs/optimisation/MuonBeamCollimators/Survived_Coll2_StoppedVsHighP-Muons.pdf}
%}
\caption{\figlabel{optim:MuBeamCollim:Torus2:contours}
Contours showing 2.5 percentage point changes to the stopping (blue) and dangerous (red) muon flux, as a function of the collimator heights.
100\% acceptance is found in the bottom right corner.
%For example, for collimator heights within the first blue contour towards the bottom-right corner, less than 2.5\% of stopped muons are lost.
}
\end{figure}
}

\newcommand{\FigOptimESTDipoleBeamHeightTwoD}{
\begin{figure}[tp]
\centering 
	\subfloat[][\figlabel{optim:ESTDipole:Beam:0.0}No Dipole]  {\includegraphics[width=1\textwidth,trim=4cm 0.5cm 9.5cm 0.7cm,clip]{figs/optimisation/EST_dipole/Tidied_signal_height-dipole_00}}\\
\subfloat[][\figlabel{optim:ESTDipole:Beam:0.1}0.1~T Dipole]{\includegraphics[width=1\textwidth,trim=4cm 0.5cm 9.5cm 0.7cm,clip]{figs/optimisation/EST_dipole/Tidied_signal_height-dipole_10}}\\
\subfloat[][\figlabel{optim:ESTDipole:Beam:0.2}0.2~T Dipole]{\includegraphics[width=1\textwidth,trim=4cm 0.5cm 9.5cm 0.7cm,clip]{figs/optimisation/EST_dipole/Tidied_signal_height-dipole_20}}\\
\caption{\figlabel{optim:ESTDipole:Beam}
The heights of signal electrons for different dipole field values.
}
\end{figure}
}

\newcommand{\FigOptimESTDipoleBeamHeightMean}{
\begin{figure}[tb]
\centering 
%	\fbox{
\includegraphics[width=0.8\textwidth,trim=1.0cm 1.3cm 3.8cm 0.9cm,clip]{figs/optimisation/EST_dipole/Tidied_NoShift-Height}
%}
\caption{\figlabel{optim:ESTDipole:MeanHeight}
Mean height of signal electrons for different values of the dipole field strength.
}
\end{figure}
}

\newcommand{\FigOptimESTDipoleBeamFluxMean}{
\begin{figure}[tb]
\centering 
\includegraphics[width=0.8\textwidth,trim=1.0cm 0.3cm 3.8cm 1.8cm,clip]{figs/optimisation/EST_dipole/Tidied_NoShift-Flux}
\caption{\figlabel{optim:ESTDipole:MeanFlux}
Survival probability for signal electrons as a function of the distance along the beamline for different values of the electron spectrometer's dipole field strengths.
}
\end{figure}
}

\newcommand{\FigOptimESTDipoleAcceptanceVsDipole}{
\begin{figure}[tb]
\centering 
\includegraphics[width=0.6\textwidth,trim=0.3cm 0.3cm 1.9cm 1.1cm,clip]{figs/optimisation/EST_dipole/Tidied_acceptance}
\caption{\figlabel{optim:ESTDipole:acceptance}
Geometric acceptance into the StrECAL detector as a function of the dipole field strength over the electron spectrometer.
}
\end{figure}
}

\newcommand{\FigOptimStopTgtPosMuStops}{
\begin{figure}[bt]
\centering 
%	\fbox{
\includegraphics[width=0.7\textwidth,trim=2.0cm 0.05cm 0.9cm 0.3cm,clip]{figs/optimisation/StopTgtPosition/Tidied_MuonStoppingRate.pdf}
%}
\caption{\figlabel{optim:StopTgtPos:MuStops}
Muon stopping rate per \ac{POT} for different target positions.
The linear behaviour arises from the reduced field strength and fixed target radius such that fewer muons impact the target as it is moved downstream.
}
\end{figure}
}

\newcommand{\FigOptimStopTgtPosSensitivitySpect}{
\begin{figure}[tb]
\centering 
%	\fbox{
\includegraphics[width=0.9\textwidth,,trim=1.3cm 0.2cm 0.7cm 0.6cm,clip]{figs/optimisation/StopTgtPosition/Tidied_-AcceptedMomentum.pdf}
%}
\caption{\figlabel{optim:StopTgtPos:AcceptedMomSpect}
The momentum dependence of the electron acceptance into the detector for different target positions.
The spectrum for each target position is normalised to the muon stopping rate for that position, such that each curve shows the sensitivity to electrons of that momentum.
}
\end{figure}
}

\newcommand{\FigOptimStopTgtPosSensitivityNoBeamBlock}{
\begin{figure}[tb]
\centering 
%	\fbox{
\includegraphics[width=0.9\textwidth,trim=1.3cm 0.2cm 0.7cm 0.6cm,clip]{figs/optimisation/StopTgtPosition/Tidied_NoBeamBlocker-AcceptedMomentum.pdf}
%}
\caption{\figlabel{optim:StopTgtPos:AcceptedMomSpectNoBeamBlock}
The momentum dependence of the electron acceptance into the detector for different target positions when the beam blocker is removed.
}
\end{figure}
}

\newcommand{\FigOptimStopTgtPosSensitivityIntegral}{
\begin{figure}[tb]
\centering 
%	\fbox{
\subfloat[][\figlabel{optim:StopTgtPos:AcceptIntegral:Adjacent}Sensitivity]{%
\includegraphics[width=0.49\textwidth,trim=0.5cm 0.5cm 0.3cm 1.9cm,clip]{figs/optimisation/StopTgtPosition/Tidied_-AcceptedMomentum-Integrated_adjacent.pdf}}
\subfloat[][\figlabel{optim:StopTgtPos:AcceptIntegral:Ratio}Change in Shape]{%
\includegraphics[width=0.49\textwidth,trim=0.5cm 0.5cm 0.3cm 1.9cm,clip]{figs/optimisation/StopTgtPosition/Tidied_-AcceptedMomentum-Integrated_ratios.pdf}}
%}
\caption{\figlabel{optim:StopTgtPos:AcceptIntegral}
\protect\subref{fig:optim:StopTgtPos:AcceptIntegral:Adjacent} 
The variation in sensitivity (acceptance $\times$ stopping rate) to electrons with different momenta as a function of the target position with respect to the nominal location. 
The darkest red line towards the top of the plot represents the sensitivity to signal, and it is that line that should therefore be maximised.
\protect\subref{fig:optim:StopTgtPos:AcceptIntegral:Ratio} 
The change in the shape of the acceptance vs. momentum spectrum as a function of the stopping target location.
}
\end{figure}
}

\newcommand{\FigOptimStopTgtPosHeights}{
\begin{figure}[p]
\centering 
\subfloat[][\figlabel{optim:StopTgtPos:Height:42.5}Electrons from 40 to 45 MeV/c]{%
\includegraphics[width=0.95\textwidth,trim=1.15cm 0.05cm 0.3cm 0.180cm,clip]{figs/optimisation/StopTgtPosition/WithBeamBlocker-Height-VaryShifts-Momentum_42-5.pdf}}%
\\\subfloat[][\figlabel{optim:StopTgtPos:Height:62.5}Electrons from 60 to 65 MeV/c]{%
\includegraphics[width=0.95\textwidth,trim=1.15cm 0.05cm 0.3cm 0.180cm,clip]{figs/optimisation/StopTgtPosition/WithBeamBlocker-Height-VaryShifts-Momentum_62-5.pdf}}%
\\\subfloat[][\figlabel{optim:StopTgtPos:Height:82.5}Electrons from 80 to 85 MeV/c]{%
\includegraphics[width=0.95\textwidth,trim=1.15cm 0.05cm 0.3cm 0.180cm,clip]{figs/optimisation/StopTgtPosition/WithBeamBlocker-Height-VaryShifts-Momentum_82-5.pdf}}%
\\\subfloat[][\figlabel{optim:StopTgtPos:Height:102.5}Electrons from 100 to 105 MeV/c]{%
\includegraphics[width=0.95\textwidth,trim=1.15cm 0.05cm 0.3cm 0.180cm,clip]{figs/optimisation/StopTgtPosition/WithBeamBlocker-Height-VaryShifts-Momentum_102-5.pdf}}%
\\\subfloat[][\figlabel{optim:StopTgtPos:Height:102.5-zoom}Electrons from 100 to 105 MeV/c (Zoomed)]{%
\includegraphics[width=0.95\textwidth,trim=1.15cm 0.05cm 0.3cm 0.180cm,clip]{figs/optimisation/StopTgtPosition/WithBeamBlocker-Height-VaryShifts-Momentum_102-5-Zoom.pdf}}%
\caption{\figlabel{optim:StopTgtPos:Height}
The effect of stopping target position on the height of electrons with a fixed momentum as they pass through the Electron Spectrometer.
The size of the variation indicates the stability of the dipole tune; the two parameters are clearly correlated.
Also striking---particularly in \protect\subref{fig:optim:StopTgtPos:Height:102.5-zoom}---is the way the dependence on the helical pitch angles is affected by the stopping target position.
}
\end{figure}
}

%\newcommand{\FigOptimStopTgtPosSensitivityIntegral}{
%\begin{figure}[tb]
%\centering 
%%	\fbox{
%\includegraphics[width=0.9\textwidth,trim=1.2cm 0.05cm 0.9cm 0.6cm,clip]{figs/optimisation/StopTgtPosition/Tidied_-AcceptedMomentum-Integrated_adjacent.pdf}
%%}
%\caption{\figlabel{optim:StopTgtPos:AcceptedIntegrated}
%The variation in sensitivity to different momentum electrons as a function of the target position, with respect to the nominal location.
%The darkest red line towards the top of the plot represents the signal sensitivity, and it is that line that should therefore be maximised.
%}
%\end{figure}
%}
%
%\newcommand{\FigOptimStopTgtPosSensitivityIntegralRatio}{
%\begin{figure}[tb]
%\centering 
%%	\fbox{
%\includegraphics[width=0.9\textwidth,trim=1.2cm 0.05cm 0.9cm 0.6cm,clip]{figs/optimisation/StopTgtPosition/Tidied_-AcceptedMomentum-Integrated_ratios.pdf}
%%}
%\caption{\figlabel{optim:StopTgtPos:AcceptedIntegratedRatio}
%The relative avvep
%}
%\end{figure}
%}

\newcommand{\FigOptimDIOBeamBlockGeometry}{
\begin{figure}[tb]
\centering 
%\fbox{
\includegraphics[width=0.8\textwidth]{figs/optimisation/BeamAndDIOBlocker/StopTgt_to_Detector-cropped.png}
%}
\caption{\figlabel{optim:DIOBeamBlock:Geometry}
Location of the beam blocker and one possible geometry for the \ac{DIO} blockers, both highlighted in green, shown here before optimisation.
}
\end{figure}
}

\newcommand{\FigOptimDIOBeamBlockESTDispersion}{
\begin{figure}[tb]
\centering 
%\fbox{
\includegraphics[width=0.8\textwidth,trim=0.3cm 0 1.7cm 0.6cm,clip]{figs/optimisation/BeamAndDIOBlocker/Tidied_ElectronDispersion-EST.pdf}
%}\\\fbox{
\includegraphics[width=0.8\textwidth,trim=0.3cm 0 1.7cm 0.6cm,clip]{figs/optimisation/BeamAndDIOBlocker/HandTidied_ElectronDispersion-EST-envelope.pdf}
%\includegraphics[width=0.8\textwidth,trim=1.5cm 0 8.5cm 3.0cm,clip]{figs/optimisation/BeamAndDIOBlocker/Tidied_ElectronDispersion-EST-envelope.png}
%}
\caption{\figlabel{optim:DIOBeamBlock:ESTDispersion}
Momentum-dependent dispersion of electrons passing through the spectrometer.
Top plot: the mean height of different momenta electrons as a function beamline distance, showing how the drift of the centre of gyration is truly proportional to the momentum.
Bottom plot: single standard deviation bands for electrons at different momentum, which shows how the envelope for different momenta overlap considerably, reducing the effectiveness of any collimators.
%Whilst the drift of the centre of gyration is proportional to momentum (top plot), there is reasonable overlap in the overall enevelope
%Mean height of electrons with different momentum from the stopping target to the detector before adding any DIO blockers and before optimised of the beam blocker.
}
\end{figure}
}

\newcommand{\FigOptimDIOBeamBlockAcceptances}{
\begin{figure}[tb]
\centering 
%\fbox{
\subfloat[\figlabel{optim:DIOBeamBlock:Acceptance:40}40 MeV/c]  {\includegraphics[width=0.35\textwidth,trim=0.0cm 0.3cm 0.5cm 1.5cm,clip]{figs/optimisation/BeamAndDIOBlocker/Tidied_Acceptance2D_40.pdf}}\hspace{1cm}
\subfloat[\figlabel{optim:DIOBeamBlock:Acceptance:60}60 MeV/c]  {\includegraphics[width=0.35\textwidth,trim=0.0cm 0.3cm 0.5cm 1.5cm,clip]{figs/optimisation/BeamAndDIOBlocker/Tidied_Acceptance2D_60.pdf}}\\
\subfloat[\figlabel{optim:DIOBeamBlock:Acceptance:80}80 MeV/c]  {\includegraphics[width=0.35\textwidth,trim=0.0cm 0.3cm 0.5cm 1.5cm,clip]{figs/optimisation/BeamAndDIOBlocker/Tidied_Acceptance2D_80.pdf}}\hspace{1cm}
\subfloat[\figlabel{optim:DIOBeamBlock:Acceptance:100}100 MeV/c]{\includegraphics[width=0.35\textwidth,trim=0.0cm 0.3cm 0.5cm 1.5cm,clip]{figs/optimisation/BeamAndDIOBlocker/Tidied_Acceptance2D_100.pdf}}
%}
\caption{\figlabel{optim:DIOBeamBlock:Acceptance}
Acceptance into the straw tracker for electrons with different momentum at the stopping target as a function of the beam and DIO blocker dimensions.
Note the logarthmic scale for the colour bar.
}
\end{figure}
}

\newcommand{\FigOptimDIOBeamBlockHitRate}{
\begin{figure}[tb]
\centering 
%\fbox{
%\hspace{-1.3cm}
\subfloat[\figlabel{optim:DIOBeamBlock:HitRate}Hit Rate from DIO]               {\includegraphics[width=0.45\textwidth,trim=0.8cm 0.3cm 0.3cm 1.05cm,clip]{figs/optimisation/BeamAndDIOBlocker/HitRate2d.pdf}}\hspace{0.04\textwidth}
\subfloat[\figlabel{optim:DIOBeamBlock:HitVAccept}Signal Acceptance vs Hit Rate]{\includegraphics[width=0.45\textwidth,trim=0.8cm 0.3cm 0.3cm 1.05cm,clip]{figs/optimisation/BeamAndDIOBlocker/HitRateVsAcceptance2dRatio.pdf}\vspace{1cm}{}}
%\subfloat[\figlabel{optim:DIOBeamBlock:HitVAccept}Signal Acceptance vs Hit Rate]{\includegraphics[width=0.45\textwidth,trim=0.8cm 0.3cm 2.5cm 1.05cm,clip]{figs/optimisation/BeamAndDIOBlocker/HitRateVsAcceptance2d.pdf}\vspace{1cm}{}}
%\hspace{0.12\textwidth}
%}
\caption{\figlabel{optim:DIOBeamBlock:HitRateAcceptance}
\protect\subref{fig:optim:DIOBeamBlock:HitRate} Number of straw tracker hits per DIO electron.
\protect\subref{fig:optim:DIOBeamBlock:HitVAccept} Ratio between the high-momentum electron ($p>100$~MeV/c) acceptance to the number of hits per DIO electron.  Colour is on a logarithmic scale.
}
\end{figure}
}

\chapter{COMET Phase-II: Optimisation}
\sectlabel{phaseII-optimisation}
%\begin{easylist}
%# Before a substantial sensitivity estimate can be made, need a solidly optimised design
%# Aiming for \sensePII within a single year of running
%# Designs previously optimised \cite{CDR}, and these results are used as nominal design / starting point
%# Fresh optimisation using new software / simulation, updated fieldmaps, physics lists and geometry
%# Some aspects fixed already since \phaseI under construction: Experiment hall, Torus1, detector solenoid, fieldmap and coil parameters?
%# Key areas for optimising:
%## Production target dimensions
%## Torus1 dipole field strength
%## Torus2 dipole field strength
%## Electron spectrometer dipole field strength
%## collimator shapes and locations
%## stopping target and beam blocker position and form
%## DIO blockers on spectrometer
%\end{easylist}
The last study into the sensitivity of \COMET \phaseII was performed in 2009~\cite{CDRphase2}, before the staged approach and \phaseI design had even been considered.
That study found that a \ac{ses} of $2.6\times10^{-17}$ could be achieved in $2\times10^{7}$ seconds of running, with a total expected background count rate of fewer than 0.34 events over the entire run period.
Since then, the collaboration's focus has shifted to design and R\&D for \phaseI and no further studies have been made of the \phaseII design.

The purpose of this chapter is to make use of the updates in the fieldmap calculation, the geometry handling, and physics modelling to revisit the design of \phaseII and demonstrate that it can do at least as well as the previous 2009 design.
In addition to updates in the simulation, some aspects of the actual design have been refined and fixed alongside \phaseI preparation, such as the experiment hall and the coils and cold-mass of the first stages of the muon beamline.
The aspects of the experiment that remain open for optimisation are shown in \tab{optimisation:possible-parameters}.

As an initial configuration, much of the design from the 2009 study will be used, with updates included for the areas of the experiment that have since been refined by the \phaseI design.
\TabOptimisationParameters%

\section{Optimisation Strategy}
%\begin{easylist}
%# Take some aspects as fixed
%# Limit scope and approach:
%## Ideally, each aspect optimised in combination to maximise signal acceptance and reduce background
%## How decoupled are each section?
%## In practise such an optimisation is not easy to do, instead aim to produce a baseline optimisation so that all backgrounds / issues can be identified
%## This can then form basis for further optimisation, with perhaps a smarter more integrated approach
%# Method:
%## Production target optimisation
%### Maximise muon and pion yield between 0 and 80 MeV at entrance to muon beamline
%### Parameters to vary: target length, target radius
%## Muon beam optimisation
%### Maximise muon stopping rate in stopping target
%### Minimise pion stopping rate
%### vary dipole along TS2 and TS4
%### vary Collimators: TS2 and at TS3
%## Electron spectrometer optimisation
%### Optimise dipole to increase signal acceptance
%### Optimse DIO blockers so DIO rate per straw is less than 1 kHz
%### Vary solenoidal field to increase separation?
%## Stopping target / beam blocker optimisation
%### Maximise reflection of signal electrons from upstream by tuning target position
%## Detector optimisation
%\end{easylist}
To perform a comprehensive optimisation of \phaseII, there are many parameters to be changed, corresponding to the shapes, positions, field strengths and so on of the different regions along the beamline.
\Tab{optimisation:possible-parameters} lists some 32 parameters, each one of which could be studied.
And yet, this number alone does not describe the full challenge of optimising the \phaseII design, since many of these parameters will correlate to one another.
For example, a correlation likely exists between the position and shapes of the muon beam collimator(s), the stopping target and the beam blocker, since all involve removing or stopping muons and other particles in the beam.
Other less intuitive correlations might also exist and a complete optimisation should be able to include the impact of these as well.

A full optimisation study then would involve a scan through a parameter space with at least 32 dimensions.
A brute force search of such a space would be nightmarishly slow and require an enormous amount of computing power.
Machine learning algorithms or intelligent scanning techniques might be able to tackle such a problem, and perhaps in the future these methods will be used.
In the meantime, however, we make use of the technique known as `physicists intuition' to approach the problem, whereby some parameters are assumed uncorrelated whilst others are disregarded on the expectation that their impact be small.
%This should be sufficient to establish a reasoned baseline geometry, whilst allowing for future work to refine the design further.
The goal of this chapter therefore is not only to optimise the experiment but also to evaluate the correlations.

This challenge is further complicated because it is not just the signal sensitivity that must be maximised, but the background rate that must be kept low.
When this work began, no start-to-finish simulation existed to estimate the signal and background rates.
As such, this study will progress by focussing on optimising the signal acceptance in the first instance, and then move to background estimation described in chapter~\sect{bagrounds}.

%Ideally this will also identify which parameters are the most strongly correlated, which would help make future optimisation studies more efficient.
The outputs of this optimisation should not be considered as final but should instead be treated as a baseline from which more intelligent approaches or physicists can iterate and improve.

%\section{Optimisation Goals}
%\begin{easylist}
%# Set sensitivity goal and optimise to reach this
%# Identify which parameters are strongly correlated to help focus and improve the efficiency of future optimisation studies.
%# Single event sensitivity only considers signal acceptance, but also need to understand backgrounds in terms of final confidence limit that can be set
%\end{easylist}

\section{Production Target Optimisation}
In the \phaseII \ac{CDR}, the production target is given as being 16~cm in
length and 4~mm in radius~\cite{CDRphase2}.  Since then, there have been changes to
the magnetic field in this region as well as the lengths and locations of
solenoids, shielding and beam-pipe, and the proton beam.  Previous studies have
looked at comparing the tungsten target proposed for \phaseII to other
materials~\cite{AEdmondsThesis}, and also drawn a comparison between
MARS~\cite{MARS1995}, Geant4~\cite{Geant42003} and the limited data available.

\subsection{Production Target Simulations} 
The goal of this study is to maximise the total muon and pion
yield below 80~MeV at the entrance to the Torus1 bent solenoid, by varying the
radius and length of the production target.
Whilst it is really the rate that negative muons stop in the Stopping Target that must be optimised, 
since it is known that muons above about 40 to 50~MeV/c will not stop, and since the bent transport solenoids themselves need optimising
we assume here that by maximising the number of muons and pions below 80~MeV/c will also maximise the stopping rate.
In addition this had the advantage of increasing the speed of the simulation.
The production target is one of the more computationally intensive aspects
of the experiment, given the relatively high energy protons and hadronic nature of their interactions with the target, which together result in a large number of secondaries.

The production target is formed from a solid cylinder of tungsten.
To find the optimal target geometry the radius and lengths were both varied individually, and then a combined scan performed where the length was held at the observed optimum. 
The location and orientation of the target were held fixed, since the axis of the proton beamline is fixed to intercept the muon beam axis
at a given point.  Once a more realistic proton beam becomes available, the location and direction of the production target
would also benefit from optimisation, however.  During the scan over
length, the back face of the target was kept 8~cm away from the muon beam since
the radiation shielding has previously been optimised.
%, and beyond this
%point the magnetic field will no longer be able to capture the pions and muons
%produced.

Primary protons were generated according to the description in the 2014 \phaseI TDR~\cite{TDR2014}.
This provides a double Gaussian position profile as well as the spread in momenta.
It does not discuss any dispersion or divergence in the beam, however.
The proton beam is currently under study, which is why the description at this stage is fairly simplistic.
The impact of a more realistic proton beam description will need to be studied in future work.

Protons in the simulation originated from a plane but, since there is some scope to tune the proton
beam's position, the input particle plane was moved to remain 1~cm away from
the front surface of the target.  Since the aim is to maximise the muon and
pion yield by varying only the length and radius, shifting the proton beam
input plane in this way removes any variation of target acceptance due to
divergences of the proton beam in the field before the target.

\FigOptimProdTgtLength
\subsection{Length Scan}
Different length targets were simulated with \num{3e5}~\ac{POT} per length.
The target length was varied in steps of 4~cm from 4 to 32~cm, whilst the target radius was held fixed at the CDR value of 4~mm.

\Fig{optimisation:ProdTgtSec:Length:Momentum} shows the momentum distributions of pions and muons for different target lengths, which are 
given as half-length following the Geant4 convention.  
\Fig{optimisation:ProdTgtSec:Length:Integral} then shows these distributions integrated up to different momenta.
From these plots it can be seen that for both muons and pions, the optimum target length occurs around a total length of 32~cm.

Additionally it can be seen from
\fig{optimisation:ProdTgtSec:Length:IntegralRatio} that the shape of the
momentum distributions changes only weakly as a function of the target length.
These plots were produced by normalising the integrated momentum contours of
\fig{optimisation:ProdTgtSec:Length:Integral} to the total integral below
400~MeV.  As a result, it is possible that the actual shape variation is even
weaker than apparent here, since in the present sample size, the high momentum
tail is not well sampled at small target lengths, such that a skew in the
normalisation might occur.

\FigOptimProdTgtRad
\FigOptimProdTgtFinal
\subsection{Radius scan}
In parallel to the length optimisation scan, different radii targets were also simulated.
Targets with radii of 2, 4, 6, 8, 10, 14, 18, 22, 26, and 30~mm were tested.
The target length was held at the \ac{CDR} value of 16~cm in total.

The results of these scans can be seen in \fig{optimisation:ProdTgtSec:Radius:Momentum} and \fig{optimisation:ProdTgtSec:Radius:Integral},
where it can be seen that a maximum in both the muon and pion yields at the entrance to the Torus1 section is achieved at a radius of about 10~mm.
As in the length scan, the varation in the shape of the momentum distributions is rather weak as a function of target radius.

\subsection{Final Result}
\FigOptimProdTgtComparePhases
Since the length and radius scan were performed in parallel, a final cross check was performed where the optimal radius was confirmed at the optimised target length.
The integrated spectrum is shown in \fig{optimisation:ProdTgtSec:Final:Integral} where it can be seen that the optimum radius once the target length is increased to 32~cm is still 10~mm.

\Fig{optimisation:ProdTgtSec:Phase1vs2} shows the total muon and pion yields at the entrance to the Torus1 for the final optimised \phaseII target and compares this to the optimised \phaseI target design.
The increased yield in the low momentum range is due to the heavier target nucleus which produces more low-momentum pions in the backwards direction.
Since only muons below around 50~MeV/c tend to stop in the target, this increase in the low momentum yield amounts to a factor 2 or so gain in the stopping muon rate per POT.

\section{Dipole Strengths of the Muon Beamline}
%\begin{easylist}
%# Production target simulation
%	## Plots of input distribtion (possibly shown in previous section)
%	## Mention gas pressure at exit of ProdTgtSec but show why this shouldn't have been an issue
%# Optimisation technique
%	## Vary dipole field strength in Torus1 and Torus2 simultaneously
%	## Mesh scan through parameter space
%	## Look at muon stopping rate
%	## Stopping target as in the CDR 
%	## No collimator material included
%# Results
%	## Muon stopping rate vs dipole field plot
%	## Pion stopping rate vs dipole field plot
%	## Discussion of corraltion
%# Selected optimal value
%\end{easylist}

The full 180\degree bent solenoid that makes up the bulk of the muon transport beam line is actually broken down into two
90\degree pieces -- Torus1 and Torus2 (also known as TS2 and TS4 by the magnet group).
Each of these sections has its own dipole field which need not provide the same dipole field strength.
The Torus1 section has been built already and was previously optimised for both \phaseI~\cite{TDR2016} and \phaseII~\cite{CDRphase2}. 
The dipole coils that it contains are designed to produce a dipole field of 0.055~T.
By running a lower current through these coils one could reduce the Torus1 dipole field without too much effort. 
However if \phaseII should require a greater dipole field strength for this region
that might be trickier since it could require additional windings to be inserted -- this would not be impossible but might costly.

\subsection{Large-sample Production Target Simulation}
Before the muon beam section could be optimised, a large set of \ac{POT} events were produced transporting all particles up to the entrance of the muon beam section.
In total, $2.3\times10^{8}$~POT were simulated through this stage, equivalent to about 1.5 bunches at \phaseII.
The production target used the optimised geometry from the previous section.
All particles that hit the surface of the Torus1 container volume were read out for later re-use in a way that ensured double-counting of particles could not occur.
In addition, particles entering the proton beam dump were also saved if later simulations wished to study their impacts.

%\CHECK{plot that shows the momentum distributions of particles in this input}
\begin{table}[t]
\CHECK{Fill this table in properly}
	\CHECK{Better yet make this and \tab{optimisation:ProdTgtSec:} into an appendix / list of simulations and configurations}
\centering
\begin{tabular}{lp{0.7\textwidth}}
	\hline
	\multicolumn{2}{l}{\bf Pions and muons from Prod.\ Tgt.\ simulation}\\
	%Energy spread, $\sigma_E$ & 0.135 MeV \\
	\hline
	\multicolumn{2}{l}{\bf Software configuration}\\
	Packages & \ttfamily{---}  \\
	Externals & \ttfamily{---}  \\
	Fieldmap & 160104 \CHECK{} with additional scale factors applied to Torus dipole field\\
	\hline
	\multicolumn{2}{l}{\bf Sample Sizes}\\
	Final scan &   \\
\hline
\end{tabular}
\caption{Key parameters in the configuration of the muon beamline dipole field optimisation.}
\tablabel{optimisation:ProdTgtSec:configuration}
\end{table}

\subsection{The Optimised Dipole Field Strengths}
The figure-of-merit for this optimisation is the muon stopping rate, which will be maximised at the optimal field configuration.
%Toshiba have previously performed a realistic field calculation using OPERA for a single coil whose on-axis field has a mean of about 0.055~T.
To identify such a configuration, a 2D-grid scan was done where the Torus1 and Torus2 dipole fields were simultaneously varied.
Varying the field strength was done by applying a scale factor to the fieldmaps for each 90\degree section described in \sect{sw:fieldmap}.
These scale factors were varied in steps of 0.125 (equivalent to 6.875~mT) and for each point muons and pions from $8\times10^6$~POT were transported to the stopping target.
The stopping rate for each combination of scale factors was then assessed.

\FigOptimMuBeamDipoleMuStops
\FigOptimMuBeamDipolePiStops

The muon stopping rate as a function of the two dipole field strengths are shon in \fig{optim:muBeamDipole:stoppedMu}.
A scale factor of 1 on the x-axis means no change to the current \phaseI design and since there is little improvement by moving to larger values for Torus1 the optimal scale factors are chosen to be 1 and 0.5 for the Torus1 and TTorus2 respectively.  
This translates to dipole field strengths of 0.055 and 0.0275~T respectively.
%\CHECK{Are these the final selected dipole field strengths?}

Also striking from \fig{optim:muBeamDipole:stoppedMu} is the anti-correlation between the two dipole field strengths.  
Roughly speaking the sum of the optimal dipole field values is constant, \ie $B_1^\text{optimal} + B_2^\text{optimal} = \text{const}$.
Although this was not previously foreseen, such a correlation can be understood accordingly.
The stopping target geometry and position is held fixed in this parameter scan, which essentially fixes the upper momentum of muons that can stop in the target.
%The initial momentum distribution grows roughly linearly with momentum for the momenta of muons that will stop in the target (see the later sections for more on this), which would suggest tuning to the higher momentum would increase the number of muons reaching the target.
%That effect however is greatly reduced by the fact that the initial muon beam has a very large transverse position profile
The drift of a particle due to the dipole field is proportional only to the distance travelled in that dipole field and the dipole field strength, and does not depend on the particle's momentum.
This means that the total drift due to the Torus1 and Torus2 dipole components is proportional to the weighted sum of the two dipole field strengths, where the weight is the distance travelled through each of the two sections.
Since each section is the same length, the time for a particle to pass through Torus1 will be roughly the same as the time in Torus2 so that the weights would be roughly equal.
This leads to the total drift being proportional to the sum of the two dipole fields.
Since the production target (as the source of muons) and stopping target (as the muons' destination) are at a fixed height with respect to one another the optimal vertical drift is also fixed so that the sum of the dipole field strengths should also be roughly fixed.

With that said, the Torus1 section has a higher pion flux which causes some assymetry between the two sections.
Keeping more pions on-axis in the Torus1 section means that more muons will enter the Torus2 section from those pions that have decayed.
But since the pion momentum distribution is slightly higher than the muon distribution keeping pions on axis requires a larger dipole field strength.
This could explain the slight assymetry where the muon stopping rate appears slightly larger if Torus1's field is larger than that of Torus2.

It is also interesting to consider the pion stopping rate as a function of the dipole field strengths.
However, as can been from \fig{optim:muBeamDipole:stoppedPi} the stopping rate is close to the level of POT events used in the simulation so that the plot is dominated by statistical fluctuations.

\section{Electron Spectrometer's Dipole}
%\begin{easylist}
%# Method
%	## Realistic muon stopping distribution
%	## Signal electrons at the target 
%	## vary the dipole field strength and study the acceptance
%	## Uniform dipole field with sharp turn-on
%# Signal height vs. dipole field strength
%	## 2D plots for no dipole, 0.1 T and 0.2 T
%	## Stack of 1D mean heights
%# Signal acceptance vs. Dipole field
%	## Survival probability vs beamline pos
%	## Overall acceptance at the detector
%# Mean height for other momenta
%\end{easylist}
The next element in the beamline after the muon transport solenoids will be the stopping target.
However, in order to study the impact of changing the stopping target parameters one will need to look at the impact on the signal acceptance into the detector.
To study that requires the components of the beamline intermediate to the target and the detector be optimised, namely the electron spectrometer.
The key free parameter in this section is the dipole field strength along the spectrometer.
The solenoidal field and solenoid aperture could also be optimised in principle, but these are considered beyond the scope of this study at this point and so the CDR values are held fixed here.
The point of this section is to establish the optimal dipole field strength given fixed target parameters which will then be studied separately together with the stability of the dipole field tune checked.

\subsection{Method and Potential Short-comings}
To study the effect of the dipole field on signal acceptance, a realistic muon stopping distribution in the target was produced by transporting muons from the production target simulation through to the stopping target.
Signal electrons were then injected at the target with the realistic stopping distribution and propagated through the beamline to the detector with different dipole field strengths.

A non-trivial short-coming of the current study is that the dipole field along the spectrometer is poorly modelled -- no realistic coil simulation exists, unlike for the bent muon transport beamlines.
As a result, a perfectly uniform dipole field is assumed with a sharp switch on and off at the entrance and exit of the spectrometer.
The impact that this has on the final result is not clear: one might expect it to be small given the relative strengths of the dipole and solenoidal fields and overall it is the integrated field that tends to matter.
However, the sharp switch-on of the field at the entrance and exit of the spectrometer is clearly not physical.
Given that the gradient introduced in the field by bending is present before the actual entrance and exit of the spectrometer (as a fringe field), some drift can be anticipated in this region.
A realistic dipole field with a realistic fringe field might overcome some of this drift however, so that the uniform field used here will most likely not capture this effect.
Nevertheless, given the absence of a realistic dipole field and scope of this study, the uniform one is the only real option at this point.

\subsection{Results}
\FigOptimESTDipoleBeamHeightTwoD
\FigOptimESTDipoleBeamHeightMean
\Fig{optim:ESTDipole:Beam} shows the projection of electron trajectories to the beam axis coordinate system for three different dipole field values.
The potency of this approach is clear from these plots; the tuneable dipole fields allow the momentum of electrons which remain on-axis to be accurately controlled (during run-time), which will benefit systematic and calibration studies that wish to study the \ac{DIO} spectrum at a lower energy.
\Fig{optim:ESTDipole:MeanHeight} then collects these plots with other dipole field strengths, plotting the mean height for all simulated electrons against the distance along the beam axis.
From this plot it can be seen how a dipole field of about 0.18~T appears optimal to keep the signal electrons on axis.

The probability for electrons to reach a given point in the beamline is shown in \fig{optim:ESTDipole:MeanFlux}, and indeed from this it can be seen that to maximize the probability of an electron reaching the detector
a dipole of around 0.18~T is desirable.
The behaviour of the low dipole field values (0 to 0.08~T) in this plot was not expected, but it is believed this is an artifact of the way this plot is made, coupled with a degree of mirroring at the entrance to the spectrometer which is enhanced as the dipole field strength increases.
If correct, a realistic dipole field calculation would be important to quantify and confirm this behaviour.
\FigOptimESTDipoleBeamFluxMean

\FigOptimESTDipoleAcceptanceVsDipole
Finally, to confirm the optimal dipole field strength the true geometric acceptance of the detector system is checked as a function of the dipole field strength, which is shown in \fig{optim:ESTDipole:acceptance}.
An electron is considered to have been geometrically accepted by the detector in this simulation if it produces at least one hit in the detector system.
In principle this could be in any straw plane, but in practice this is almost always in the first layer of straws.
Since this is a different way to analyze the acceptance compared to the survival probability it would not suffer to the artifact seen for low dipole field values in \fig{optim:ESTDipole:MeanFlux}.
Nonetheless, \fig{optim:ESTDipole:acceptance} confirms that the optimal dipole field strength is very close to 0.18~T.

A second important conclusion can be drawn from the fact that the dependence on the dipole field strength is relatively weak around the optimal value of 0.18~T.
A change of about 10\% in the dipole field strength only reduces the signal acceptance by about 3\% whilst a change of about 5\% would see a reduction of only about 0.7\%.


\section{Stopping Target Position}
%\begin{easylist}
%# Method
%	## Shift target position up or downstream by 50 cm 
%	## Muon stopping distribution with realistic beam sim. only include muons for speed
%	## Electron acceptance using realistic stopping distribution
%	## Don't change anything else about the target (size, length, number of disks, beam blocker distance, etc)
%	### Too many parameters, want only to achieve a complete baseline here.
%# Muon stopping rate vs, position
%# Electron acceptance times muon stopping rate = signal sensitivity vs. position
%# Stability of dipole field tune 
%# Impact on low energy electrons with regards to the subsequent DIO blocker tune stability
%# Impact of beam blocker on acceptance
%## See appendix for further discussion and further optimisation in this region
%\end{easylist}
The final aspect to be studied with regards to maximizing the signal sensitivity is the stopping target.
In principle there are many parameters related to the stopping target such as the location, disk shape (profile and thickness), and disk spacing
In addition the beam blocker should be considered in parallel with the stopping target since it sits so close to the target itself and can be expected to have a big impact on the signal acceptance.
However, this leaves far too many parameters to be considered all at once.

Since the field around the target tapers sharply various competing factors must be considered.
For example, prior to the stopping target region the muon beam is transported through a 3~T solenoidal field.
The magnetic field in the target region, however, tapers to about 1~T, which would cause the envelope of the muon beam to grow.
Moving the stopping target downstream would mean that the muon beam arrives with a larger aperture, and would therefore prefer a larger stopping target or else fewer muons will actually hit the target.
On the other hand, from the perspective of signal acceptance, the tapered field can be used to mirror signal electrons that are initially produced heading upstream therefore increasing the signal acceptance.
Moving the target further upstream then will reduce this effect as the difference between the magnetic field strengths at the exit of the bent muon transport solenoid and at the stopping target is reduced.

Given this and the need to reduce the number of parameters inspected, the target and beam blocker design was held fixed in this study and only the position was changed by moving the target upstream and downstream by $\pm50$~cm with respect to the nominal target location as given in the CDR~\cite{CDRphase2}.
Given that the target disks will occupy in total about 1~m, a shift of 50~cm corresponds to half the target length.
In each different position of the stopping target, as for the spectrometer dipole optimisation, a realistic stopping distribution was produced by running muons from the large-scale production target simulation through to the target.
This stopping distribution was then re-used to introduce signal electrons accordingly.
Additionally, low momentum electrons were also studied in order to study the impact of target position on the height of both signal and low-energy electrons as they pass through the spectrometer.
This is important both to check the correlation of the dipole field tune with the stopping target position, but also how the subsequent DIO blocker height optimisation will correlate the target position.
\FigOptimStopTgtPosMuStops

\Fig{optim:StopTgtPos:MuStops} shows how the rate of muon stops per \ac{POT} is affected by changing the position of the stopping target.
The relationship is roughly linear, dropping from around $2.2\times10{-3}$ muon stops per POT when the target is shifted upstream by 50~cm to about $1.3\times10^{-3}$ muon stops per POT if the target is shifted 50~cm in the other direction.
This relationship is as expected given the fixed radius of the target and the growth of the muon beam aperture arising from the reduction in the field strength.

\FigOptimStopTgtPosSensitivitySpect
In \fig{optim:StopTgtPos:AcceptedMomSpect} one can see the way the electron acceptance changes for different target positions.
Acceptance here is defined as producing at least one hit in the detector and the momentum shown is the momentum at the target, which is not necessarily the same as the momentum at which the electron is observed.
Each histogram in \fig{optim:StopTgtPos:AcceptedMomSpect} is normalised to the number of primary electrons introduced at the target per MeV/c and then scaled to the muon stopping rate.
This normalisation makes the value of each curve proportional to the sensitivity of the experiment to different momentum electrons, up to factors such as analysis cuts like timing and reconstruction quality.

\FigOptimStopTgtPosSensitivityIntegral
Since the parameter we wish to optimise here is the location of the stopping target, \fig{optim:StopTgtPos:AcceptIntegral:Adjacent} represents the same data as in \fig{optim:StopTgtPos:AcceptedMomSpect} but with each line representing the content of a different 5~MeV/c bin as a function of the target position.
For signal, it is the 105~MeV/c line (dark burgundy) that is most important and it can be seen that this is optimised for shifts upstream of the nominal position from between 10 and 20~cm.
It is also interesting to note that the acceptance of lower energy electrons is relatively decreased as the target is moved upstream, as can be seen in \fig{optim:StopTgtPos:AcceptIntegral:Ratio}.
This could be useful as a way to suppress hit rate from DIO electrons later.

Whilst we do not intend to optimise the beam blocker at this point, to check the impact that it has on the optimisation of the target position simulations were performed where the blocker was completely removed.
\Fig{optim:StopTgtPos:AcceptedMomSpectNoBeamBlock} shows the product of the stopping rate and electron acceptance when the blocker is removed.
From this the trend is much cleaner, electrons below 70~MeV/c are suppressed as you move the target upstream whereas the high energy electron acceptance increased.
\FigOptimStopTgtPosSensitivityNoBeamBlock


\FigOptimStopTgtPosHeights
Finally the relationship between the stopping target position and the mean height of electrons through the spectrometer is demonstrated in \fig{optim:StopTgtPos:Height}.
Each plot shows the mean heights for electrons with a given momentum for different stopping target positions and it is clear that there is some correlation between the mean height and the position of the stopping target.
For this reason, a more complete optimisation should consider optimising the two parameters simultaneously.
Nonetheless the change is not particularly large, only about a few cm difference by te end of the spectrometer for signal electrons.
This correlation is also likely related to the way the acceptance changes for different target positions when the beam blocker is removed.

The second striking feature from the plots in \fig{optim:StopTgtPos:Height} is the way the mean height aquires a strong sinusoidal component for large target shifts upstream.
This suggests that when the target is shifted upstream the electrons passing the entrance to the spectrometer tend to have a particular value for the pitch and phase angles of the trajectories.
Several separate mechanisms could produce this effect.
Firstly the acceptance around the target itself could aquire a stronger preference for certain pitch and phase angles when the target is moved upstream.
Secondly, since the stopping target disks will see more of the muon beam upstream, the muon stopping distribution could become less homogenous.  Whilst electrons are produced isotropically path with less target material along them will tend to accept outgoing electrons more readily such that if more muons stop at one side of the target than the other a dependence on pitch and phase could arise.

Clearly then the stopping target region is a very complicated area; even though this study has focussed on a single parameter -- the position of the target itself -- it is clear that this correlates to many other variables.
This is a region in the experiment most ripe for further optimisation then, and indeed appendix~\ref{sec:appendix:stopTgtImprove} shows some of the first steps that have made in this direction since the optimisation described here was completed.
Unfortunately, at the time this work was carried out an error in the normalisation of these plots lead to the conclusion that the optimal shift was between 0 and 10~cm upstream.
Given the complexity of the optimisation in this region, it was therefore decided to keep the stopping target at the nominal location for the subsequent steps.
Having corrected the normalisation of the plots, the conclusion now is that the optimal location is between 10 and 20~cm upstream, so perhaps the target should have been shifted back.
However the improvement to the sensitivity would have been small: shifting the target back about 20~cm improves the sensitivity by around 2\% compared to the signal acceptance at nominal position.
%Given this small improvement and the complexity of this region it was decided that the stopping target was kept to the nominal position.

\section{Collimators in the Muon Beamline}
\FigOptimMuBeamDipoleMuDispersion
%\begin{easylist}
%# Muon beam without collimators
%	## Height vs momentum before and after 180 degrees
%	## Beamline projection for stopped and high-P muons
%	## Subtraction of the good and bad muons to suggest the region of interest for collimators
%	## Why the assymetry upstream?
%# Collimator optimisation
%	## Analytic collimator study to reduce simulation time
%	### Cannot study secondary particles
%	### Assumes all particles entering the collimators are killed
%	## Current collimator from Phase-I are in a different location to those suggested here
%# Results of study, recommended collimator design
%\end{easylist}
With the beam line optimised for high signal efficiency, one can look at improving the background by adding collimators into the muon beam line to reduce the flux of high-momentum muons and pions.

\Fig{optim:MuBeamDipole:MuDispersion} shows the dispersive effect of the bent solenoid and dipole field on muons passing through the beamline.
Thanks to this dispersion low momentum muons that can stop in the target and high momentum muons which could produce backgrounds are separated sufficiently so that material in the beam pipe can selectively remove the dangerous, high-momentum muons with only a small impact on the muon stopping rate.

\subsection{Collimator Placement}
\FigOptimMuBeamCollimMuonPaths
The plots in \fig{optim:MuBeamCollim:Beamline} give a sense of where best to locate the collimating material.
In \fig{optim:MuBeamCollim:Beamline:All} the paths of all muons along the beamline is shown.
\Fig{optim:MuBeamCollim:Beamline:Stopped} then separates out the muons that stop in the target which should be compared to 
\fig{optim:MuBeamCollim:Beamline:HighP} showing the paths of muons that reach the stopping target region with momentum greater than 70~MeV/c (the threshold for a muon to decay to an electron with $p>100$~MeV/c).

It is interesting to note the apparent asymmetry in the high-momentum muons at the entrance to the Torus1 that can be seen in \fig{optim:MuBeamCollim:Beamline:HighP}.  
This is not because of some momentum-dependence in the acceptance of the preceding beamline, but due to the fact that muons in that plot are only included if they are `dangerous' in the region around the stopping target.
Even without additional collimators, the beam pipe itself removes high momentum muons that enter in the lower half of the beamline.
The validity of only tagging high momentum muons around the stopping target region comes from the assumption that the products of high-momentum muons that decay before this region can be reliably removed.
It is important then that this assumption be checked, but for this work this is left as a task for the future.

Finally then in \fig{optim:MuBeamCollim:Beamline:Diff} the difference between the high-momentum muons and the paths of those that stop is shown. 
Regions in blue on this plot show where many more high-momentum muons pass than stopping muons; it is in these locations that collimator material will be most effective.
This approach could also be improved, since taking the straight difference between the two plots implies equal weighting for stopping muons and high-momentum ones.
In reality, whilst a muon tagged as stopping is definitely going to stop, a high-momentum muon should be weighted by the probability to produce a signal-like electron and the probability that this electron survives to be accepted into the detector, passing all analysis cuts.
This would make the weighting for the high-momentum muons be a function of the beamline distance itself which again requires a study into how high-momentum electrons are accepted.
However for the purposes of obtaining a qualitative sense of where to collimate the unweighted difference should be sufficient.

Two regions of interest appear: in the upper half of the entrance to Torus1 and the lower half of the exit of Torus2.  
Collimating at the Torus1 entrance is justified by the basis that high momentum muons will tend to have larger gyroradii compared to the muons that stop, and that at this point the beam is largely on-axis and has not yet been dispersed.
Collimating at the exit of Torus2 is readily understood on the grounds that the high momentum muons will have all drifted downwards by this point, compared to the low momentum stopping muons which are kept on-axis by the dipole fields.

\FigOptimMuBeamCollimTransverseSep
These conclusions are backed by the plots shown in \fig{optim:MuBeamCollim:TransverseSep} which show transverse slices through the beamline at the entrance to Torus1 (0\degree of bent solenoid), the midpoint between Torus1 and Torus2 (after 90\degree of bent solenoid) and at the exit of Torus2 (after 180\degree).
From these plots one can also see how in the middle of the bent solenoids (at 90\degree) the separation between muons that will stop and those that will have momentum greater than or equal to 70~MeV/c in the stopping target region is weakest, and hence collimators in this location will not be so effective.

\FigOptimMuBeamCollimTorusOne
\FigOptimMuBeamCollimTorusTwoFraction
\subsection{Collimator Height Optimisation}
To identify the optimum height for the collimators in a computationally efficient way, events were generated without any collimators included.
The full three dimensional trajectory of all particles as well as the decay tree were persisted.
This allows for a `virtual' collimator to be used, where particles that enter a defined region and their secondaries are removed from the downstream plots.
Whilst this method allows for many collimator shapes and heights to be tested quickly, it comes with a couple of limitations.
Firstly, the accuracy depends on the trajectory sampling density, which should be fine but this results in large data sizes.
Secondly, realistic material effects of the collimator cannot be captured, such as the probability a particle is simply scattered rather than stopping completely, or the result of secondary particles produced in the collimator itself.

\Fig{optim:MuBeamCollim:Torus1} shows the results of lowering the bottom edge of the collimator material in Torus1.
\Fig{optim:MuBeamCollim:Torus1:perPOT} shows the probability per \ac{POT} that different types of particle (stopped muons and high momentum muons, pions, and electrons) pass the collimator as a function of the collimator height.
On the other hand, \fig{optim:MuBeamCollim:Torus1:fraction} shows the same plots but normalised to the total number of each particle type that reaches the collimator in the first place, therefore showing the survival probability along just the collimator region.
Based on these plots, for a collimator that starts at 120~mm above the beamline axis, 14\% of high momentum pions are removed, high momentum muons are suppressed by 14\% whilst the muon stopping rate is reduced only by 3\%.

\FigOptimMuBeamCollimTorusTwoContours

For the second collimator at the exit to Torus2, the situation is slightly more complicated since in principle the optimum height could be correlated to the height of the upstream Torus1 collimator.
To account for this, the virtual collimator technique was applied for both the Torus1 and Torus2 collimator sections simultaneously.
As a result, the 1-dimensional plots of \fig{optim:MuBeamCollim:Torus1} become 2D as can be seen in 
%\fig{optim:MuBeamCollim:Torus2:perPOT}, which is normalised to the rate per POT, and 
\fig{optim:MuBeamCollim:Torus2:fraction}, which is normalised to the particle flux just before the collimator similar to \fig{optim:MuBeamCollim:Torus1:fraction}.
%\FigOptimMuBeamCollimTorusTwoPerPOT

\FigOptimMuBeamCollimMuonPathsWColl
\Fig{optim:MuBeamCollim:Torus2:contours} represents the stopping and high-momentum plots in a way that is easier to compare the two directly.
Each line in that plot is a contour showing a change of 2.5 percentage points to the yield.  Total acceptance, or 100\% is in the bottom right corner.
From this plot it can be seen that whilst keeping more than 97.5\% of the muon stopping rate (the bottom-right most blue contour), the maximum high momentum muon suppression is achieved when the Torus1 collimator sits about 140~mm above the beam axis, and the Torus2 collimator sits about 120~mm below it.
To be precise, at these collimator values  the muon stopping rate is kept at 99\% of the no-collimator rate, whilst high momentum pions, muons and electrons drop to 27.6\%, 20.9\% and 11\% (although this last value is very statistically limited) respectively, compared to their no-collimator rates.
At this point the selected values for the collimator heights give conservative background suppressions and tighter values could be chosen.
Given that backgrounds from high-momentum muons are suppressed compared to the actual rate of `dangerous' muons by the geometric acceptance of the remaining beamline and the timing and momentum cuts this seems reasonable at this stage although this will of course be investigated in the next chapter.

Finally, for comparison to the original plots, \fig{optim:MuBeamCollim:BeamWColl} shows the impact the new collimators have on the muon components of the beam.
It is clear how greatly reduced the number of muons passing the stopping target region has now become.
\clearpage

\section{The Beam and Decay-in-Orbit Blockers}
\sectlabel{optim:BeamAndDIOBlocker}
%\begin{easylist}
%# Dispersion of electrons through the spectrometer
%# Width of beam plot (1 sigma band around mean height for different momentum?)
%# Method for optimisation
%# Goal: Hit rate suppression, not backgrounds
%# Signal, DIO acceptance
%# Hit rate
%\end{easylist}
\FigOptimDIOBeamBlockGeometry
\FigOptimDIOBeamBlockESTDispersion
For every muonic atom formed in the stopping target some 39\% will undergo \acf{DIO}.
With some $1.4\times10^{8}$~\ac{POT} per bunch and a muon stopping efficiency around $1.7\times10^{-3}$ stops per \ac{POT}, one can expect about $9\times10^{4}$ DIO events per bunch.
If the detector system were exposed to this level of hits it would be impossible to resolve any signal electrons.
However, above the end-point energy for electrons coming from free muon decay, the DIO spectrum falls extremely quickly, with around 99\% of DIO electrons being produced with less than 59~MeV/c.
To this end the electron spectrometer's primary purpose is to disperse away the electrons below 60~MeV/c whilst keep signal electrons on-axis, such that material in the beamline can be tuned to remove the low energy DIO electrons.
The beam and DIO blockers are highlighted in green in \fig{optim:DIOBeamBlock:Geometry}.

\Fig{optim:DIOBeamBlock:ESTDispersion} demonstrates the dispersion that appears at the end of the spectrometer with the nominal beam blocker design and no DIO blockers in place.
The DIO blockers will be inserted along the bottom of the spectrometer and tuned to scrape away a sufficient number of DIO electrons, which will tend to travel towards the bottom of the beam pipe.
However whilst the centre of gyration drifts with a well-behaved proportionality in a bent solenoid, the actual separation between signal and \ac{DIO} electrons is in reality less clear as can be seen by the lower plot in \fig{optim:DIOBeamBlock:ESTDispersion}.

In addition to the DIO blockers along the spectrometer, the material of the beam blocker immediately after the stopping target also plays a role in suppressing the \ac{DIO} rate since low energy electrons remain closer to the beam axis, as can be seen in the lower plot of \fig{optim:DIOBeamBlock:ESTDispersion} around the stopping target.

This section therefore describes a simultaneous optimisation of the beam blocker radius and the DIO blocker height.
The overall goal is to suppress the DIO rate to less than a single DIO electron per bunch whilst maintain maximal signal acceptance.
As for the muon beam collimators, we use the analysis-based collimator approach, where no blocking material is included during simulation.
Instead a high-sampling density for particle trajectories is used such that during analysis particles and their secondaries can be `killed' if they enter a region that would contain material of either the DIO or beam blocker.

\FigOptimDIOBeamBlockAcceptances
\FigOptimDIOBeamBlockHitRate
The results of this study are shown in \fig{optim:DIOBeamBlock:Acceptance}, where the geometric acceptance into the detector is shown for four different electron momenta as a function of the beam blocker radius and DIO blocker height.
It is clear that for all values of the blockers' dimensions, electrons above 100~MeV/c have a much better acceptance.

The mean hit rate per DIO event is shown in \fig{optim:DIOBeamBlock:HitRate}.  
It is formed by, for each combination of DIO and beam blocker dimension, the weighted integral of the acceptance of electrons as a function of momentum with the mean DIO rate in each momentum bin.
To compare the hit rate to the acceptance of signal, the ratio between the hit rate and the high-momentum electron acceptance is shown in \fig{optim:DIOBeamBlock:HitVAccept}.
It is clear from this figure how quickly (note the logarithmic colour scale) the DIO hit rate is suppressed by increasing the DIO and beam blocker dimensions compared to the signal acceptance.

With a beam blocker of 24~cm radius and DIO blockers set to 35~cm below the beam axis, the DIO hit rate is about $2.2\times10^{-5}$ per DIO event, or about 2 DIO hits per bunch.
For the same blocker dimensions, the geometric signal acceptance is about 0.22\%.
Given the steep drop-off in hit rate versus signal acceptance around these values a finer scan in this region is an important check for the future.
What hit rate in the straw tracker is tolerable is a number for future studies.  In \phaseI, the straw tracker is expected to operate with a hit rate around 1~kHz per straw.
\phaseII will likely use finer straws, but with the current \phaseI design, some 133 straws occupy a layer, so that a total hit rate into the Straw Tracker of around 200~kHz should be acceptable.

Given the scope for future improvements in the granularity of the \phaseII straw tracker and the sensitivity to the beam and DIO blocker dimensions, the somewhat non-conservative values described above of 24~cm and 35~cm for beam blocker radius and DIO~blocker depth below the beam axis are selected here.

%\section{Revisiting the Torus2 Collimator after Beam Background Studies}
\section{Summary of optimised parameters}
The complete set of optimised parameters is shown in \tab{optim:AllParameters}.

\begin{table}[btp]
	\centering
	%\begin{tabular}{r|lp{6cm}}
		\begin{tabular}{lL{0.2\textwidth}L{0.35\textwidth}}
		Parameter & Value & Comments \\
		\hline
		\hline
		Production Target Length & 32~cm & \multirow{2}{4cm}{Placed asymmetrically about muon beamline axis} \\
		Production Target Radius & 1~cm & \\[1ex]
		\hline
		Torus1 (TS2) Dipole Strength & 0.055~T & same as \phaseI \\
		Torus2 (TS4) Dipole Strength & 0.0275~T & same as \phaseI \\
		\hline
		Torus1 Entrance Collimator Height& 14~cm above beam axis & From top of beam pipe downwards  \\
		Torus2 Exit Collimator Height& 12~cm below beam axis & From bottom of beam pipe upwards\\
		\hline
		Stopping Target Shift& 0~cm& Unchanged from CDR value (2.8~m from centre of final coil in Torus2 to front of target )  \\
		\hline
		Electron Spectrometer Dipole & -0.18~T& Negative compared to Torus1 and Torus2 dipole fields \\
		\hline
		Beam Blocker Radius & 24~cm &  \\
		DIO Blocker Height & 35~cm below beam axis& From bottom of spectrometer upwards  \\
		\hline
		\hline
	\end{tabular}
\caption{\tablabel{optim:AllParameters}
Optimised values for the parameters studied in this chapter.
Many more parameters remain to be optimised that were considered beyond the scope of the present work.
}
\end{table}

\section{Future optimisations}
The primary goal of this work is to update the previous optimisation from the 2009 CDR and provide a new baseline design.
However whilst touching every aspect of the layout of \phaseII, the optimisation developed here is not exhaustive and there is much scope for further work.

The following list is a short a summary of some of the areas that could be developed further:
\begin{description}
\item[ Refine the optimisation criteria]
	In many of these studies only the signal efficiency, or even some proxy, is used to identify the optimal value of the parameter under question.
	In reality one ought also to study simultaneously the background rates.  
	These two quantities should be combined into the expected confidence limits given a null or background-only measurement.
	At the same time other quantities such as cost and run-time will also need considering, although the latter is likely reduced simultaneously with maximising the overall sensitivity.
\item[ High energy electron acceptance vs. beamline distance ]
	When tuning the muon beam collimators, the key goal is to remove high-energy particles that can produce high energy electrons. 
	A particularly useful study would be to evaluate the acceptance to signal-like electrons which originate along the beamline.
	This should include those electrons which originate with momentum greater than 105~MeV/c but which arrive at the detector with signal-like energies.
	With this information, it becomes easier to identify how soon along the beamline one must collimate away the high-energy particles which may lead to improvements in the muon stopping rate or the signal acceptance.
\item[ Stopping target shape, thickness, and disk spacing ]
	This study has focussed only on the position of the stopping target.
	Clearly the actual shape should be studied as well. In particular, given the dispersion in the muon beam, and the changing solenoidal field strength around the target, a target design that uses disks of varying profile or even thickness has potential to improve the experimental sensitivity.  
\item[ Correlation between dipole fields and stopping target shape ]
	Since the bent solenoids introduce dispersion into the beam and the dipole fields compensate for this, the height and momentum of the muon beam at the stopping target can be controlled to a degree.
	Whilst in the optimisation of the muon beam dipoles of this thesis the target shape was held fixed, in principle some of the identified correlations might be different if the target design was allowed to vary simultaneously.  
	Such a study would clearly involve an enormous parameter space, so perhaps this is a study that could be best performed with a smarter machine-learning based approach.
\item[ Spectrometer's solenoidal field strength and aperture]
	Since the dispersive effect of a bent solenoid is inversely proportional to the solenoidal field strength, reducing the field in the spectrometer will likely improve the DIO--signal separation, reducing the hit rate.
	This could however also affect the signal acceptance since the trajectories will acquire larger envelopes.
	Increasing the spectrometer aperture size could compensate this, but then one faces an increase in the cost of the spectromter and the detector solenoid which would need to be considered.
\end{description}

\newcommand{\FigSensMuMomentum}{%
\begin{figure}[bt]
\centering 
%\fbox{
\includegraphics[width=0.8\textwidth,trim=0.5cm 0.1cm 1.5cm 0.8cm,clip]{figs/sensitivity/Muon_momentum.pdf}
%}
\caption{\figlabel{sense:muMomenta}
The Momentum and rates of muons reaching the final beam collimator, the stopping target, and actually stopping in the target.
It can be seen how the present target geometry is unable to stop muons of greater than around 50~MeV/c.
}
\end{figure}\xspace}

\newcommand{\FigSensMuStopsTwoD}{%
\begin{figure}[tb]
\centering 
\subfloat[\figlabel{sense:stops:XZ}X-Z plane]{
        \includegraphics[width=0.33\textwidth,trim=0.3cm 1.5cm 0.7cm 1.2cm,clip]{figs/sensitivity/MuStops_2D_XZ.pdf}}
\subfloat[\figlabel{sense:stops:XY}X-Y plane]{
	\includegraphics[width=0.33\textwidth,trim=0.3cm 1.5cm 0.7cm 1.2cm,clip]{figs/sensitivity/MuStops_2D_XY.pdf}}
\subfloat[\figlabel{sense:stops:ZY}Z-Y plane]{
	\includegraphics[width=0.33\textwidth,trim=0.3cm 1.5cm 0.7cm 1.2cm,clip]{figs/sensitivity/MuStops_2D_ZY.pdf}}
\\
\subfloat[\figlabel{sense:stops:X}X axis]{
	\includegraphics[width=0.33\textwidth,trim=0.3cm 1.5cm 0.7cm 1.2cm,clip]{figs/sensitivity/MuStops_1D_X.pdf}}
\subfloat[\figlabel{sense:stops:Y}Y axis]{
	\includegraphics[width=0.33\textwidth,trim=0.3cm 1.5cm 0.7cm 1.2cm,clip]{figs/sensitivity/MuStops_1D_Y.pdf}}
\subfloat[\figlabel{sense:stops:Z}Z axis]{
        \includegraphics[width=0.33\textwidth,trim=0.3cm 1.5cm 0.7cm 1.2cm,clip]{figs/sensitivity/MuStops_1D_Z.pdf}}
\caption{\figlabel{sense:stops}
Projections of the final position of stopped muons in the stopping target.
	Axes are from the SimG4 global coordinate system, so that $+X$ points away from the production target, $+Y$ is vertically upwards, and $+Z$ is the direction of the muon beam at the production target.
	The muon beam in these plots is therefore travelling in the negative-$Z$ direction having passed around through 180\degree of the bent solenoid.
}
\end{figure}\xspace}

\newcommand{\FigSensGeomAccept}{%
\begin{figure}[bt]
\centering 
\subfloat[\figlabel{sense:accept:height}Path of Signal Electrons in the Beamline Projection]{
	\includegraphics[width=0.99\textwidth,trim=8.3cm 1.5cm 13.2cm 2.2cm,clip]{figs/sensitivity/Tidied_SignalHeight2DVsBeamline.png}}\\
\subfloat[\figlabel{sense:accept:flux}Survival Probability for Signal Electrons]{
        \includegraphics[width=0.99\textwidth,trim=1.0cm 0.2cm 1.7cm 0.4cm,clip]{figs/sensitivity/Tidied_SignalSurivivalVsBeamline.pdf}}
\caption{\figlabel{sense:accept}
Geometric acceptance of signal events.
\protect\subref{fig:sense:accept:height} Projection of the trajectories of signal electrons to the surface formed by the beamline axis and the vertical direction.
\protect\subref{fig:sense:accept:flux} Survival probability of signal electrons as a function of the distance along the beamline axis.
The x-axis range is the same in the two plots.
From these, it is clear how the acceptance is diminished by the DIO blocker in the spectrometer, although from that point on the rate of signal loss reduces.
}
\end{figure}\xspace}

\newcommand{\FigSensMomTransfer}{%
\begin{figure}[tb]
\centering 
%\fbox{
\includegraphics[width=0.5\textwidth,trim=0.0cm 0.0cm 0.0cm 1.6cm,clip]{figs/sensitivity/MomentumTransfer.pdf}
%}
\caption{\figlabel{sense:momTransfer}
The transfer matrix for electrons originating at the target, including the geometric acceptance and energy loss.
}
\end{figure}\xspace}

\newcommand{\FigSensMomSpectra}{%
\begin{figure}[tb]
\centering 
%\fbox{
\subfloat[Incl.~energy loss\figlabel{sense:spectra:ELoss}]             {\includegraphics[width=0.48\textwidth,trim=0.0cm 0.0cm 1.0cm 1.0cm,clip]{figs/sensitivity/ConversionVsDio_Spectra.pdf}}
\subfloat[Energy loss \& resolution\figlabel{sense:spectra:resolution}]{\includegraphics[width=0.48\textwidth,trim=0.0cm 0.0cm 1.0cm 1.0cm,clip]{figs/sensitivity/ConversionVsDio_Spectra-wResolution.pdf}}
\caption{
The spectrum of electrons coming from \ac{DIO} and \mueconv assuming a conversion rate of $\mathcal{R}=\sci{3}{-16}$.
Black dashed lines indicate the total electron distribution that would be seen (the sum of $S$ and $B$) if the signal has the assumed conversion rate of \num{3e-16}, half that rate, and one tenth that rate.
\protect\subref{fig:sense:spectra:ELoss} Includes energy losses in the target, beamline, and detector;
\protect\subref{fig:sense:spectra:resolution} also includes resolution effects (for a Gaussian resolution function with a standard deviation of $\sigma=200$~keV/c).
\figlabel{sense:spectra}}
\end{figure}\xspace}

\newcommand{\FigSensMomIntegral}{%
\begin{figure}[tb]
\centering 
%\fbox{
\includegraphics[width=0.99\textwidth,trim=0.5cm 0.0cm 0.5cm 0.3cm,clip]{figs/sensitivity/ConversionVsDio_Integrated.pdf}
%}
\caption{\figlabel{sense:integral}
Relative signal versus \ac{DIO} background as a function of the low-momentum cut value assuming a conversion rate of $\mathcal{R}=\sci{3}{-16}$.
The magenta line is the signal over square root of signal plus background for this conversion rate shown as an indicator of the optimum cut value.
}
\end{figure}\xspace}

\newcommand{\FigSensTiming}{%
\begin{figure}[b]
\centering 
%\fbox{
\subfloat[\figlabel{sense:timing:signal}Signal Arrival Time]      {\includegraphics[width=0.9\textwidth,trim=0.9cm 0.1cm 1.5cm 0.20cm,clip]{figs/sensitivity/160823_MuonLifetime.pdf}}\\
\subfloat[\figlabel{sense:timing:efficiency}Timing Cut Efficiency]{\includegraphics[width=0.9\textwidth,trim=0.9cm 0.1cm 1.5cm 0.4cm,clip]{figs/sensitivity/160823_TimingCutEfficiency.pdf}}
\caption{\figlabel{sense:timing}
Timing of signal electrons.
\protect\subref{fig:sense:timing:signal} The arrival time of signal electrons at the detector, including the effect of the proton pulse width, particle transportation, and the muon lifetime.
\protect\subref{fig:sense:timing:efficiency} the efficiency of the timing window as a function of the switch-on time.  Assumes a pulse separation of 1.17~$\mu$s.
}
\end{figure}\xspace}

\newcommand{\TabSensParams}{%
\begin{table}[b]
\centering
\begin{tabular}{lll}
	\hline
         $I_p$                          & 7~$\mu$A & Proton beam current                            \\ 
	 ${R}_{\mu/p}$      & \VarMuStopsPerPOT & Muon stopping rate per POT                     \\ 
%	 $f_\mathrm{1s}$                & 90\%    & Probability of reaching the ground state \\ 
	 $\mathcal{B}_\mathrm{capture}$ & 61\%    & Branching ratio for muon nuclear capture in Al \\ 
	 $A_{\mu\rightarrow e}$         & \VarTotalSignalAcceptance   & Total signal acceptance of \phaseII            \\ 
	\hline
\end{tabular}
        \caption{\tablabel{sense:ses}
        Parameters that determine the run time and single event sensitivity for COMET \phaseII based on this study.
        }
\end{table}\xspace}

\newcommand{\TabSensEstimates}{%
\begin{table}[t]
\centering
	\begin{tabular}{L{0.25\textwidth}p{0.14\textwidth}S[table-format=2.1]p{0.12\textwidth}p{0.23\textwidth}}
	\hline
							& Single event sensitivity & \multicolumn{1}{p{0.13\textwidth}}{Total \acp{POT} ($\times10^{19}$)} &Beam time $t_\textrm{run}$ (s) & SES in one year of continuous beam \\ 
  \hline
  COMET \phaseII\hspace{1ex}(this study)                & \VarPredictedSES         & 68.3 & \VarRunTime[3]                  & \VarPredictedSESPerYear                         \\ 
  COMET \phaseII\hspace{1ex}(CDR 2009~\cite{CDRphase2}) & \sci{2.6}{-17}           & 85   & \sci{2.00}{7}                  & \sci{1.65}{-17}                         \\ 
  Mu2e~\cite{Mu2e2014}                                  & \sci{2.4}{-17}           & 36   & \sci{6.00}{7}                  & \sci{4.57}{-17}                         \\ 
  COMET \phaseI~\cite{TDR2016}                          & \sci{3.0}{-15}           & 3.2 & \sci{1.26}{7}                  & \sci{1.19}{-15}                         \\ 
	\hline
\end{tabular}
        \caption{\tablabel{sense:comparisons}
        Comparison between the run time and single-event sensitivity from this study and from the 2009 CDR, the \phaseI TDR, and  the Mu2e experiment's TDR.
	The \ac{ses} in one year of continuous beam is the single-event-sensitivity that can be achieved in \num{3.15e7}~seconds of running, assuming no beam shutdown periods.
        }
\end{table}\xspace}

\newcommand{\FigDIOBackground}{
\begin{figure}[tbp]
\centering
%\fbox{
\includegraphics[width=1.0\textwidth,trim=0 0 1cm 0.93cm,clip]{figs/backgrounds/Dio_BackgroundRateVsRuntime.pdf}
%}
\caption{
The DIO background rate as a function of momentum threshold for different total running times.
Given a fixed running time, the total number of stopped muons is also fixed, which in turn sets the signal sensitivity and the DIO background rate.
All signal acceptance parameters were held fixed, except for the efficiency of the momentum threshold, which, when combined with the number of stopped muons, determines the \ac{ses}.
The \ac{ses} is indicated in the number along the lines in units of \num{1e-17}.
\figlabel{bg:dio:rates}}
\end{figure}
}

\newcommand{\FigDIOEndPointComparison}{
\begin{figure}[tbp]
\centering
\includegraphics[width=0.8\textwidth,trim=0 0 0 0,clip]{figs/backgrounds/CompareDIOEndpoints.pdf}
\caption{
Comparison of the various available end-point expansions.
The red and blue lines show the parametrisations reported in the literature, whilst the black shows the digitisation of the spectrum used in SimG4.
For this study, the more conservative parametrisation from the 2011 Czarnecki paper~\cite{Czarnecki2011} has been used.
\figlabel{bg:dio:spectra}}
\end{figure}
}

\newcommand{\FigRMCExperiments}{
\begin{table}[tbp]
\centering
%\fbox{
\includegraphics[width=0.5\textwidth]{figs/backgrounds/RMC_Gorringe_ExperimentSummary.pdf}
%}
\caption{
Summary of experimental values of the rate of \ac{RMC} producing photons with energy greater than 57~MeV, $R_\gamma$, and the observed end-point, $k_\textrm{max}$, redacted from~\cite{RevModPhys.76.31}.
The column lablled `$\alpha$' is the neutron excess for the element, determined by: $\alpha=(A-2Z)/Z$.
\tablabel{bg:rmc:experiments}}
\end{table}
}

\newcommand{\TabRMCEndPoints}{%
\begin{table}[tb]%
%\centering
\begin{tabular}{lS[table-format=2.6]SS}%
\hline
Reaction & \multicolumn{1}{C{3cm}}{Atomic Mass of Daughter (u)} & \multicolumn{1}{C{2cm}}{$\Delta{}M$ (MeV/c$^{2}$)}&\multicolumn{1}{C{2cm}}{$\max(E_e^\textrm{RMC})$ (MeV/c$^{2}$)}\\
\hline
${}^{27}$Al$(\mu,\gamma\nu){}^{27}  $Mg     & 26.984341 &  3.12  & 101.85 \\
${}^{27}$Al$(\mu,\gamma\nu2n){}^{26}$Mg     & 25.982593 &  9.56  &  95.41 \\
${}^{27}$Al$(\mu,\gamma\nu2n){}^{25}$Mg     & 24.985837 & 20.66  &  84.31 \\
${}^{27}$Al$(\mu,\gamma\nu{}p){}^{26}$Na    & 25.992633 & 18.13  &  87.37 \\
${}^{27}$Al$(\mu,\gamma\nu{}np){}^{25}$Na   & 24.989954 & 23.71  &  81.77 \\
${}^{27}$Al$(\mu,\gamma\nu{}d){}^{25}$Na    & 24.989954 & 21.49  &  84.00 \\
${}^{27}$Al$(\mu,\gamma\nu\alpha){}^{23}$Na & 22.994467 & 15.49  &  91.01 \\
\hline
\end{tabular}
\caption{%
Several potential daughter nuclei of nuclear muon capture in \textsuperscript{27}Al.
The mass of \textsuperscript{27}Al is 26.98153863~$u$, and one $u$ is taken as 931.494061~MeV/c$^2$~\cite{PDG2014}.
All masses come from~\cite{AUDI20033}.\tablabel{bg:rmc:massDifferences}}\end{table}%
\xspace}%

\newcommand{\FigRMCSimResults}{
\begin{figure}[tbp]
\centering
%\fbox{
\includegraphics[width=0.85\textwidth]{figs/backgrounds/RMC_simResults.pdf}
%}
\caption{
Observed electrons from a simulation of \num{6e7} \ac{RMC} photons.
The overlaid spectrum is normalised arbitrarily to fit on the plot.
\figlabel{bg:rmc:simulation}}
\end{figure}
}

\newcommand{\FigRPCData}{
\begin{figure}[btp]
\centering
\subfloat[][\figlabel{bg:rpc:data:ca}Calcium]  {\includegraphics[width=0.43\textwidth]{figs/backgrounds/RPC-data-calcium.png}}\hspace{0.2cm}%
\subfloat[][\figlabel{bg:rpc:data:mg}Magnesium]{\includegraphics[width=0.53\textwidth]{figs/backgrounds/RPC-data-magnesium.png}}
\caption{
Spectrum of photons coming from \acf{RPC}~\cite{Bistirlich:1972jy}.
The spectrum of manesium, which is adjacent to aluminium on the periodic table, was used as the basis of these studies.
\figlabel{bg:rpc:data}}
\end{figure}
}

\newcommand{\FigRPCSimulatedSpectrum}{
\begin{figure}[btp]
\centering
%\fbox{%
\includegraphics[width=0.73\textwidth,trim=1cm 0.5cm 2cm 1cm,clip]{figs/backgrounds/RPC_simulated_spectrum.pdf}%
%}
\caption{
Digitised (red) and smoothed (blue) spectrum of \ac{RPC} from magnesium (see \fig{bg:rpc:data:mg}) used as input to the Monte Carlo simulation.
\figlabel{bg:rpc:spectrum}}
\end{figure}
}

\newcommand{\FigPionStopDist}{
\begin{figure}[btp]
\centering
\subfloat[][\figlabel{bg:piStop:dist:x}X-direction]{\includegraphics[width=0.32\textwidth,trim=0.2cm 0 1cm 0.7cm,clip]{figs/backgrounds/Tidied_StoppedPi-X.pdf}}\hspace{0.1cm}%
\subfloat[][\figlabel{bg:piStop:dist:y}Y-direction]{\includegraphics[width=0.32\textwidth,trim=0.2cm 0 1cm 0.7cm,clip]{figs/backgrounds/Tidied_StoppedPi-Y.pdf}}\hspace{0.1cm}%
\subfloat[][\figlabel{bg:piStop:dist:z}Z-direction]{\includegraphics[width=0.32\textwidth,trim=0.2cm 0 1cm 0.7cm,clip]{figs/backgrounds/Tidied_StoppedPi-Z.pdf}}
\caption{
Stopping distributions of pions in the target.
These distributions have considerably different forms to the muon stopping distributions shown in \fig{sense:stops}, mostly due to the different momenta of muons and pions.
\figlabel{bg:piStop:dist}}
\end{figure}
}

\newcommand{\FigPiVsMuMomenta}{
\begin{figure}[btp]
\centering
%\fbox{%
\includegraphics[width=0.9\textwidth,trim=0 0.5cm 1.3cm 0.4cm,clip]{figs/backgrounds/Tidied_MuVsPiMomentum.pdf}%
%}
\caption{
The momentum of muons and pions for those that reach the target area and those that actually stop.
It is clear how the pion momenta are in general higher, including those that stop, although the maximum stopping momentum for pions is similar to that of muons.
\figlabel{bg:piVsMu:momenta}}
\end{figure}
}

\newcommand{\FigRPCSimResults}{
\begin{figure}[btp]
\centering
%\fbox{
\subfloat[][\figlabel{bg:rpc:sim:momVtime}Momentum Vs.\ Time]{%
\includegraphics[width=0.45\textwidth,trim=0.2cm 0 0cm 1.9cm,clip]{figs/backgrounds/Tidied_RPC_sim_mom-time.pdf}}
%}
%\begin{minipage}[b]{0.45\textwidth}
%\subfloat[][\figlabel{bg:rpc:sim:time}Arrival Time]{%
%\includegraphics[width=\textwidth,trim=0.9cm 0.3cm 1cm 0.5cm,clip]{figs/backgrounds/Tidied_RPC_sim_time.pdf}}\\
%\subfloat[][\figlabel{bg:rpc:sim:mom}Momentum]{%
%\includegraphics[width=\textwidth,trim=0.9cm 0.3cm 1cm 0.5cm,clip]{figs/backgrounds/Tidied_RPC_sim_mom.pdf}}%
%\end{minipage}\hspace{1ex}
\hspace{1em}%
\subfloat[][\figlabel{bg:rpc:sim:time}Arrival Time of High-$p$ Electrons]{%
\includegraphics[width=0.49\textwidth,trim=0 0 0 2.8cm,clip]{figs/backgrounds/RPC_lifetime.png}}%
%\fbox{
%}
\caption{
Detection of secondaries from RPC photons in the target.
Although many high-momentum electrons are detected \protect\subref{fig:bg:rpc:sim:momVtime}, they are all well before the time-gated detected window \protect\subref{fig:bg:rpc:sim:time}.
\figlabel{bg:rpc:sim}}
\end{figure}
}

\newcommand{\FigAntiprotonMeco}{
\begin{figure}[tbp]
\centering
\includegraphics[width=0.6\textwidth]{figs/backgrounds/Antiproton_Meco24_energy.pdf}
\caption{
Variation in the antiproton production rate as a function of incident proton energy, according to Meco note 24~\cite{Meco024} and used in the COMET TDR~\cite{TDR2016}.
For reference, protons with 8~GeV kinetic energy have 8.89~GeV/c momentum, whilst with 10.14~GeV kinetic energy their momentum is 11.038~GeV/c.
The vertical coloured lines have been added to indicate these energies, whilst the horizontal bands show the range of predicted cross sections for the models of proton-nucleon and proton-nucleus interaction.
\figlabel{bg:antiprotons:meco24}}
\end{figure}
}

\newcommand{\FigAntiprotonData}{
\begin{figure}[tbp]
\centering
\includegraphics[width=1.0\textwidth,trim=0 0 0 0,clip]{figs/backgrounds/Antiproton_RatePerPOT_data.pdf}
\caption{
Experimental data for antiproton production rates for 10~GeV protons~\cite{Boyarinov:1994tp,Kiselev:2012sj}.
Each line represents the cross section obtained for the four different target materials covered in those papers, scaled to match the number of nucleons of tungsten and with the additional factors of \eq{bg:antiprotons:rate} included.
\figlabel{bg:antiprotons:data}}
\end{figure}
}

\newcommand{\FigAntiprotonEndpoint}{
\begin{figure}[btp]
\centering
\subfloat[][\figlabel{bg:antiprotons:end-point:tungsten}Tungsten]{\includegraphics[width=0.49\textwidth,clip=true,trim=0 0 1cm 1.7cm]{figs/backgrounds/Antiproton_Tungsten_theta_lab.pdf}}%\hspace{0.5cm}%
\subfloat[][\figlabel{bg:antiprotons:end-point:carbon}Carbon    ]{\includegraphics[width=0.49\textwidth,clip=true,trim=0 0 1cm 1.7cm]{figs/backgrounds/Antiproton_Carbon_theta_lab}}
\caption{
The kinematic end-point for antiproton production as a function of the outgoing antiproton direction with respect to the incoming proton in the frame of the target nucleus (the lab frame).
The absolute end-point is only achieved when the nucleus and outgoing protons recoils coherently.
\figlabel{bg:antiprotons:end-point}}
\end{figure}
}

\newcommand{\FigAntiprotonFits}{
\begin{figure}[tbp]
\centering
%	\fbox{
\includegraphics[width=1.0\textwidth]{figs/backgrounds/AntiprotonFits.png}
%}
\caption{
Piecewise fitting to experimental data and kinematic end-points.
Inlays show a zoom around the experimental data points.
\figlabel{bg:antiprotons:fits}}
\end{figure}
}

\newcommand{\FigAntiprotonAngularDependence}{
\begin{figure}[b]
\centering
\includegraphics[width=0.8\textwidth,trim=0 0 1.4cm 1cm,clip]{figs/backgrounds/AntiprotonAngularDependence.pdf}
\caption{
The angular dependence of the rate of antiproton emission, integrated over all momenta.
The different lines represent the different fits to the high momentum part of the spectrum.
The relationship given in~\cite{Boyarinov:1994tp} would suggest the data here should fit a straight line.
The dashed lines represent instead a quadratic fit to these points, which looks like a better fit.
For reweighting events the interpolated (straight solid) lines were used to be conservative.
\figlabel{bg:antiprotons:angular}}
\end{figure}
}

\newcommand{\TabAntiprotonRegions}{
\begin{table}[t]
\centering
\sisetup{table-number-alignment=right,table-format=1.2e3}%
\begin{tabular}{a|r|SS}
\multicolumn{1}{c|}{\multirow{2}{*}{Region}} & \multirow{2}{*}{Data Source} & \multicolumn{2}{c}{Total $\bar{p}$ per POT in this region} \\ 
                                             &                              & {Linear Tail}       &  {Exponential Tail} \\ 
\hline
                  0-59\degree    & 10\degree \cite{Kiselev:2012sj}    & 9.13e-05 & 5.26e-05 \\ 
                  59-97\degree   & 59\degree \cite{Kiselev:2012sj}    & 2.64e-08 & 4.17e-09 \\ 
                  97-119\degree  & 97\degree \cite{Boyarinov:1994tp}  & 3.40e-12 & 1.74e-12 \\ 
                  119-180\degree & 119\degree \cite{Boyarinov:1994tp} & 2.58e-12 & 5.71e-13 \\ 
\hline
\end{tabular}
\caption{
Angular regions and the source of the data used to build the momentum spectrum for that region.
The integrated rate for the two different high-momentum tail descriptions are also given.
Note that these values do \emph{not} contain the correction for the different incident proton energies;
for the COMET proton beam the antiproton yield is expected to be a factor 0.12 times those given here.
%The values in the final column are result of converting to rates per POT and integrating the differential cross-sections measured in \cite{Boyarinov:1994tp,Kiselev:2012sj}.
%integrated the fitted and extrapolated spectra and then integrates over the fitted angular dependence.
\tablabel{bg:antiprotons:regions}}
\end{table}
}

\newcommand{\FigAntiprotonSimHeightsTwoDPbar}{%
\begin{figure}[ph]%
\centering %
\subfloat[][\figlabel{bg:antiprotons:sim:2D-antip:10}Production between 0 and 59\degree]{%
\includegraphics[width=1\textwidth,trim=3.7cm 0.3cm 1.8cm 0.8cm,clip]{figs/backgrounds/Antiproton_height2D_antiproton_10.png}}\\%
\subfloat[][\figlabel{bg:antiprotons:sim:2D-antip:59}Production between 59 and 97\degree]{%
\includegraphics[width=1\textwidth,trim=3.7cm 0.3cm 1.8cm 0.8cm,clip]{figs/backgrounds/Antiproton_height2D_antiproton_59.png}}\\%
\subfloat[][\figlabel{bg:antiprotons:sim:2D-antip:97}Production between 97 and 119\degree]{%
\includegraphics[width=1\textwidth,trim=3.7cm 0.3cm 1.8cm 0.8cm,clip]{figs/backgrounds/Antiproton_height2D_antiproton_97.png}}\\%
\subfloat[][\figlabel{bg:antiprotons:sim:2D-antip:119}Production between 119 and 180\degree]{%
\includegraphics[width=1\textwidth,trim=3.7cm 0.3cm 1.8cm 0.8cm,clip]{figs/backgrounds/Antiproton_height2D_antiproton_119.png}}%
\caption{
The heights of antiprotons passing along the beamline for the four different angular regions of productions.
Each antiproton trajectory is weighted by the probability of producing an antiproton at this angle.
The colour scale on all these plots is the same.%
\figlabel{bg:antiprotons:sim:2D-antip}}%
\end{figure}%
\xspace}

\newcommand{\FigAntiprotonSimHeightsTwoDPiMin}{%
\begin{figure}[ph]%
\centering %
\subfloat[][\figlabel{bg:antiprotons:sim:2D-pi:10}Production between 0 and 59\degree]{%
\includegraphics[width=1\textwidth,trim=3.7cm 0.3cm 1.8cm 0.8cm,clip]{figs/backgrounds/Antiproton_height2D_pi-_10.png}}\\%
\subfloat[][\figlabel{bg:antiprotons:sim:2D-pi:59}Production between 59 and 97\degree]{%
\includegraphics[width=1\textwidth,trim=3.7cm 0.3cm 1.8cm 0.8cm,clip]{figs/backgrounds/Antiproton_height2D_pi-_59.png}}\\%
\subfloat[][\figlabel{bg:antiprotons:sim:2D-pi:97}Production between 97 and 119\degree]{%
\includegraphics[width=1\textwidth,trim=3.7cm 0.3cm 1.8cm 0.8cm,clip]{figs/backgrounds/Antiproton_height2D_pi-_97.png}}\\%
\subfloat[][\figlabel{bg:antiprotons:sim:2D-pi:119}Production between 119 and 180\degree]{%
\includegraphics[width=1\textwidth,trim=3.7cm 0.3cm 1.8cm 0.8cm,clip]{figs/backgrounds/Antiproton_height2D_pi-_119.png}}%
\caption{
The heights of secondary pions passing along the beamline produced from antiprotons in each of the four different angular regions of productions.
Each trajectory is weighted by the probability of producing the parent antiproton in its initial direction at the target.
\figlabel{bg:antiprotons:sim:2D-pi}}%
\end{figure}%
\xspace}%

\newcommand{\FigAntiprotonSimFluxes}{
\begin{figure}[b]
\centering 
\subfloat[][\figlabel{bg:antiprotons:sim:fluxes:antip}Unweighted Antiproton Survival Probability]{
\includegraphics[width=1\textwidth,trim=0.7cm 0 1.9cm 0.2cm,clip]{figs/backgrounds/Antiproton_fluxes.pdf}}\\
\subfloat[][\figlabel{bg:antiprotons:sim:fluxes:pion}Unweighted Secondary Pion Transport Probability]{
\includegraphics[width=1\textwidth,trim=0.7cm 0 1.9cm 0.2cm,clip]{figs/backgrounds/Antiproton_fluxes-pions.pdf}}
\caption{
The surivival probability of antiprotons and secondaries pions per antiproton produced in the target as a function of distance along the beamline.
These plots are not weighted by the probability that an antiproton is produced at a particular angle.
From left to right the vertical gray lines indicate the production target, Torus1 entrance, and the Torus2 exit.
\figlabel{bg:antiprotons:sim:fluxes}}
\end{figure}
}

\newcommand{\FigAntiprotonSimTime}{
\begin{figure}[b!]
\centering 
\subfloat[][\figlabel{bg:antiprotons:sim:time:antip}Timing of Antiprotons]{
\includegraphics[width=0.485\textwidth,trim=0.5cm 0.9cm 0.5cm 0.9cm,clip]{figs/backgrounds/Antiproton_timing_antiprotons.pdf}}
\hspace{1ex}\subfloat[][\figlabel{bg:antiprotons:sim:time:pion}Timing of Secondary Pions]{
\includegraphics[width=0.485\textwidth,trim=0.5cm 0.9cm 0.5cm 0.9cm,clip]{figs/backgrounds/Antiproton_timing_pions.pdf}}
\caption{
%\CHECK{Add pion stopping timing to RPC section to be able to compare to it here}
	The arrival time of antiprotons~\protect\subref{fig:bg:antiprotons:sim:time:antip} and pions~\protect\subref{fig:bg:antiprotons:sim:time:pion}
	at various points along the beamline and for the different initial antiproton directions.
Note that the x-axis scales are different.
Whilst the timing of antiprotons themselves is very delayed, the timing of secondary pions, which are produced predominantly at the production target, is relatively prompt and will be effective at suppressing the induced backgrounds.
\figlabel{bg:antiprotons:sim:time}}
\end{figure}
}

%\newcommand{\FigAntiprotonSimFluxes}{
%\begin{figure}[ph]
%\centering 
%\subfloat[][\figlabel{bg:antiprotons:sim:fluxes:10}Production between 0 and 59\degree]{
%\includegraphics[width=1\textwidth,trim=0.7cm 0 1.9cm 0.2cm,clip]{figs/backgrounds/Antiproton_flux_10.pdf}}\\
%\subfloat[][\figlabel{bg:antiprotons:sim:fluxes:59}Production between 59 and 97\degree]{
%\includegraphics[width=1\textwidth,trim=0.7cm 0 1.9cm 0.2cm,clip]{figs/backgrounds/Antiproton_flux_59.pdf}}\\
%\subfloat[][\figlabel{bg:antiprotons:sim:fluxes:97}Production between 97 and 119\degree]{
%\includegraphics[width=1\textwidth,trim=0.7cm 0 1.9cm 0.2cm,clip]{figs/backgrounds/Antiproton_flux_97.pdf}}\\
%\subfloat[][\figlabel{bg:antiprotons:sim:fluxes:119}Production between 119 and 180\degree]{
%\includegraphics[width=1\textwidth,trim=0.7cm 0 1.9cm 0.2cm,clip]{figs/backgrounds/Antiproton_flux_119.pdf}}
%\caption{\figlabel{bg:antiprotons:sim:fluxes}
%The surivival probability of antiprotons and their secondaries per antiproton produced in the target as a function of distance along the beamline.
%From left to right the vertical magenta lines indicate the production target, Torus1 entrance, and the Torus2 exit.
%}
%\end{figure}
%}

\newcommand{\FigAntiprotonSimPiMom}{
\begin{figure}[tb]
\centering 
\includegraphics[width=0.9\textwidth,trim=2.0cm 0 0 0,clip]{figs/backgrounds/Antiproton_pion_momentum.pdf}
\caption{
Momentum of pions passing the exit of Torus1 (90\degree around the bent muon beam transport solenoid) compared to pions in the main muon beam (which is arbitrarily normalised).
\figlabel{bg:antiprotons:sim:piMom}}
\end{figure}
}

\newcommand{\HeaderPi}[1]{\multicolumn{1}{#1}{Torus1 $\pi^-$}}
\newcommand{\HeaderPBar}[1]{\multicolumn{1}{#1}{$\bar{p}$ Stop}}
%\newcommand{\TabAntiprotonResults}{
%\begin{table}[tb]
%\centering
%\begin{tabular}{a|SS|SS|SS|}
%\multicolumn{1}{c|}{Region} & \multicolumn{2}{c|}{Observed Events} & \multicolumn{2}{c|}{Weighted Mean per $\bar{p}$}& \multicolumn{2}{c}{Rate per \ac{POT}}  \\
%\multicolumn{1}{c|}{}       & \HeaderPi{r}    & \HeaderPBar{r|}    & \HeaderPi{r}      & \HeaderPBar{r|}      &\HeaderPi{r}     & \HeaderPBar{r}                  \\
%\hline
%   0-59\degree &50&0&2.5e-4&& \\
%  59-97\degree &31&0&1.6e-4&& \\
% 97-119\degree &36&0&1.8e-4&& \\
%119-180\degree &64&9&3.2e-4&& \\
%\hline
%\multicolumn{1}{c|}{Total} & & & & & \\
%\hline
%\end{tabular}
%\caption{\tablabel{bg:antiprotons:results}
%Results of the antiproton simulation.
%`Torus1 $\pi^-$' are the pions that pass the exit of Torus1, which is 90\degree round the muon beamline.
%`$\bar{p}$ Stop' refers to the number of antiprotons that stop in the muon Stopping Target.
%The weighted mean is the sum of the observed events weighted by the production probability given the initial antiproton direction, divided by the total number of input antiprotons.
%Finally the Rate per \ac{POT} is weighted mean scaled by the number of antiprotons produced for this region per POT (last column of \tab{bg:antiprotons:regions}).
%}
%\end{table}
%}

%\newcommand{\TabAntiprotonResultsPiSecond}{
%\begin{table}[tb]
%\centering
%\begin{tabular}{a|SSS}
%\multicolumn{1}{c|}{Region} & \multicolumn{1}{c}{Observed Events} & \multicolumn{1}{c}{Weighted Mean per $\bar{p}$}& \multicolumn{1}{c}{Rate per \ac{POT}}  \\
%\hline
%   0-59\degree &50&2.5e-4&& \\
%  59-97\degree &31&1.6e-4&& \\
% 97-119\degree &36&1.8e-4&& \\
%119-180\degree &64&3.2e-4&& \\
%\hline
%\multicolumn{1}{c|}{Total} & & & & & \\
%\hline
%\end{tabular}
%\caption{\tablabel{bg:antiprotons:results}
%Results of the antiproton simulation.
%`Torus1 $\pi^-$' are the pions that pass the exit of Torus1, which is 90\degree round the muon beamline.
%`$\bar{p}$ Stop' refers to the number of antiprotons that stop in the muon Stopping Target.
%The weighted mean is the sum of the observed events weighted by the production probability given the initial antiproton direction, divided by the total number of input antiprotons.
%Finally the Rate per \ac{POT} is weighted mean scaled by the number of antiprotons produced for this region per POT (last column of \tab{bg:antiprotons:regions}).
%}
%\end{table}
%}

%\newcommand{\TabAntiprotonResultsTrans}{
%\begin{table}[p]
%\centering
%\begin{tabular}{a|SSS|SSS|S|}
% &\multicolumn{3|}{c}{Raw count} &\multicolumn{3|}{c}{Weighted Probability} & \\
% &\multicolumn{1}{p{0.6cm}}{Entr.}&\multicolumn{1}{p{0.6cm}}{Mid.}&\multicolumn{1}{p{0.6cm}|}{Target}&\multicolumn{1}{p{0.6cm}}{Entr.}&\multicolumn{1}{p{0.6cm}}{Mid.}&\multicolumn{1}{p{0.6cm}}{Target}&\multicolumn{1}{p{0.6cm}}{$P(\textrm{Target}|\textrm{Entr.}$}\\
%\hline
%\multicolumn{1}{l}{Pions} & \multicolumn{7}{p{6cm}}{}\\
%   0-59\degree    & 16943  & 1452  & 1  & 5.77E-06 & 5.10E-07 & 9.98E-11 & 1.73E-05 \\ 
%   59-97\degree   & 87230  & 7157  & 18 & 4.64E-09 & 4.06E-10 & 2.99E-13 & 6.45E-05 \\ 
%   97-119\degree  & 157385 & 13041 & 25 & 4.10E-14 & 3.59E-15 & 5.02E-18 & 1.22E-04 \\ 
%   119-180\degree & 222486 & 12997 & 30 & 1.25E-17 & 9.59E-19 & 1.91E-21 & 1.53E-04 \\ 
%\hline
%\multicolumn{1}{l}{Antiprotons} & \multicolumn{7}{p{6cm}}{}\\
%0-59\degree    & 7      & 3     & 0    & 1.40E-09 & 5.14E-11 & 3.23E-12                  &          \\ 
%59-97\degree   & 270    & 58    & 8    & 1.70E-12 & 4.81E-13 & 1.16E-14                  & 2.41E-02 \\ 
%97-119\degree  & 2907   & 830   & 114  & 1.10E-16 & 3.31E-17 & 5.17E-19                  & 1.56E-02 \\ 
%119-180\degree & 278105 & 20787 & 2237 & 2.27E-19 & 5.20E-20 & 7.74E-21                  & 1.49E-01 \\ 
%               &        &       &      &          &          & \multicolumn{1}{c}{Mean=} & 6.29E-02 \\ 
%%\hline
%\end{tabular}
%\end{table}
%}

\newcommand{\TabAntiprotonResultsPiSecond}{%
\sisetup{table-number-alignment=right, table-format=2.2e2}%
\begin{table}[p]%
\centering%
%\begin{adjustwidth}{-0.7cm}{}%
\begin{tabular}{lr|S!{\VRule}S!{\VRule}S!{\VRule}S|}%
\multicolumn{2}{c|}{\multirow{2}{3cm}{Rates for Secondary $\pi^-$}} & \multicolumn{4}{c|}{Secondary $\pi^-$ from Angular Region for Antiproton Production} \\ 
                              &  & \multicolumn{1}{a!{\VRule}}{0-59\degree} & \multicolumn{1}{a!{\VRule}}{59-97\degree} & \multicolumn{1}{a!{\VRule}}{97-119\degree} & \multicolumn{1}{a|}{119-180\degree} \\ 
\hline
\multirow{3}{1.6cm}{Raw Counts} & Torus1 & \multicolumn{1}{r!{\VRule}}{16943} & \multicolumn{1}{r!{\VRule}}{87230} & \multicolumn{1}{r!{\VRule}}{157385} & \multicolumn{1}{r|}{222486} \\ 
                               % & TS3 & \multicolumn{1}{r!{\VRule}}{1452}  & \multicolumn{1}{r!{\VRule}}{7157}  & \multicolumn{1}{r!{\VRule}}{13041}  & \multicolumn{1}{r|}{12997}  \\ 
                                & Torus2 & \multicolumn{1}{r!{\VRule}}{208}   & \multicolumn{1}{r!{\VRule}}{1106}  & \multicolumn{1}{r!{\VRule}}{1995}   & \multicolumn{1}{r|}{1847}   \\ 
                                & Target & \multicolumn{1}{r!{\VRule}}{1}     & \multicolumn{1}{r!{\VRule}}{18}    & \multicolumn{1}{r!{\VRule}}{25}     & \multicolumn{1}{r|}{30}     \\ 
\hline
\multirow{3}{1.6cm}{Weighted Mean} & Torus1  &6.93E-07 & 5.57E-10 & 4.92E-15 & 1.50E-18\\ 
                                %  & TS3  &6.12E-08 & 4.87E-11 & 4.31E-16 & 1.15E-19 \\
                                  & Torus2  &1.03E-08 & 7.11E-12 & 6.95E-17 & 1.47E-20\\ 
                                  & Target  &1.20E-11 & 3.59E-14 & 6.02E-19 & 2.29E-22\\ 
%\hline
%\multicolumn{2}{r|}{$P(\textrm{Target}|\textrm{Torus2})$}  & 0.039 & 0.040 &  0.022 & 0.017 \\
\hline
\end{tabular}%
%\end{adjustwidth}%
\caption{
Secondary pion fluxes from antiprotons observed at key points along the beamline.
See caption to \tab{bg:antiprotons:sim:antip} for a description of the column and row contents.
%The raw counts are the total observed events from the simulation of 80M antiprotons for each of the four angular regions.
%The weighted sum then shows the weighted sum over every particle passing the point, where the weight is determined as described by equation~\eq{bg:antiprotons:factorisation}.
\tablabel{bg:antiprotons:sim:pi}}
\end{table}\xspace}

\newcommand{\TabAntiprotonResultsAntip}{%
\sisetup{table-number-alignment=right, table-format=2.2e2}%
\begin{table}[p]%
%\begin{adjustwidth}{-0.7cm}{}%
\begin{tabular}{lr|S!{\VRule}S!{\VRule}S!{\VRule}S|}%
\multicolumn{2}{c|}{\multirow{2}{3cm}{Rates for $\bar{p}$ Transport}} & \multicolumn{4}{c|}{Angular Region for Antiproton Production} \\ 
                              &  & \multicolumn{1}{a!{\VRule}}{0-59\degree} & \multicolumn{1}{a!{\VRule}}{59-97\degree} & \multicolumn{1}{a!{\VRule}}{97-119\degree} & \multicolumn{1}{a|}{119-180\degree} \\ 
\hline
\multirow{3}{1.6cm}{Raw Counts} & Torus1 &\multicolumn{1}{r!{\VRule}}{7} & \multicolumn{1}{r!{\VRule}}{270} & \multicolumn{1}{r!{\VRule}}{2907} & \multicolumn{1}{r|}{278105} \\ 
                                & Torus2 &\multicolumn{1}{r!{\VRule}}{3} & \multicolumn{1}{r!{\VRule}}{58 } & \multicolumn{1}{r!{\VRule}}{830 } & \multicolumn{1}{r|}{20787 } \\ 
                                %& Torus2 Collim. & \multicolumn{1}{r!{\VRule}}{1} & \multicolumn{1}{r!{\VRule}}{39 } & \multicolumn{1}{r!{\VRule}}{431 } & \multicolumn{1}{r|}{9024  } \\ 
                                & Target &\multicolumn{1}{r!{\VRule}}{0} & \multicolumn{1}{r!{\VRule}}{8  } & \multicolumn{1}{r!{\VRule}}{114 } & \multicolumn{1}{r|}{2237  } \\ 
\hline
\multirow{3}{1.6cm}{Weighted Mean} & Torus1 & 1.68E-10    & 2.04E-13 & 1.32E-17 & 2.72E-20 \\ 
                                   & Torus2 & 6.16E-12 & 5.77E-14 & 3.97E-18 & 6.23E-21\\
                                   %& Torus2 Collim. & 1.95E-15    & 6.18E-15 & 1.85E-18 & 3.15E-21 \\ 
                                   & Target & {*}4.39E-16 & 1.39E-15 & 6.20E-20 & 9.29E-22 \\ 
\hline                                                           
\multicolumn{2}{r|}{$P(\textrm{Target}|\textrm{Torus2})$} &{*}0.225& 0.225  & 0.034  & 0.295\\
\hline
\end{tabular}%
%\end{adjustwidth}%
\caption{
Antiproton stopping rates and fluxes at key points: at the entrance to the bent
solenoids (Torus1), just before the final beam collimator (Torus2), and in front of the
stopping target (Target).
Raw counts are the total number of particles seen in each simulation of 80
million antiprotons. The weighted mean is the average of all particle
weights, given by the energy correction factor $\xi$
and the antiproton production angle, $\Phi(\theta,\phi)$.
$P(\textrm{Target}|\textrm{Torus2})$, gives the survival
probability for an antiproton to reach the target, given that it reached the Torus2 collimator.
Since no antiprotons were seen stopping for the angular region from 0 to
59\degree, the value for the final entry in that column---indicated
with an asterisk (*)---are obtained by multiplying the antiproton rate
expected at the Torus2 collimator with the median collimator acceptance from
the other three angular regions.%
\tablabel{bg:antiprotons:sim:antip}}
\end{table}\xspace}

%\newcommand{\TabAntiprotonFactors}{
%\begin{table}[tb]
%\centering
%        \begin{tabular}{llm{0.5\textwidth}}
%	\hline
%        Parameter & \multicolumn{1}{l}{Value} & Description \\
%	\hline
%        $R_{\pi/p}$                & \VarPiStopsPerPOT & Pion stopping rate per \ac{POT}  \\ 
%        $\mathcal{B}_\textrm{RPC}$ & \num{2.27e-2} & Branching ratio of \ac{RPC} \\ 
%	$f_{e,\textrm{RPC}}$           & \VarDetectedEsPerRPC & Probability of an RPC photon producing signal-like electrons in the detector \\ 
%	$A_\textrm{time}$              & \VarRPCTimingEfficiency & Acceptance of timing window to secondary electrons from RPC \\ 
%        $\epsilon_\textrm{extinction}$ & \VarExtinctionFactor[2] &  Extinction factor\\ 
%	\hline
%\end{tabular}
%\caption{\tablabel{bg:antiprotons:factors}
%Parameters and their values in the determination of the \ac{RPC} background rate.
%}
%\end{table}
%}

\newcommand{\TabAntiprotonEstimates}{
\sisetup{table-number-alignment=right,table-format=2.2e3,table-comparator=true}%
\begin{table}[b]
%\begin{adjustwidth}{-0.9cm}{}
\begin{tabular}{la|S|SSS|}
            &         & \multicolumn{1}{c|}{Stopping Rates} & \multicolumn{3}{c|}{Background Rate per POT} \\ 
                           &          &            & {No Timing} & {Prompt} & {Delayed} \\%& {Unweighted}
\hline                                                                              %              
\multirow{4}{*}{$\bar{p}$} & 0-59    & 4.39E-16 & 2.03E-22 & {-}      & 1.04E-22 \\ % & <6.58E-13
                           & 59-97   & 1.39E-15 & 6.43E-22 & {-}      & 3.30E-22 \\ % & 4.17E-16
                           & 97-119  & 6.20E-20 & 2.87E-26 & {-}      & 1.47E-26 \\ % & 2.48E-18
                           & 119-180 & 9.29E-22 & 4.29E-28 & {-}      & 2.20E-28 \\ % & 1.60E-17
\hline
\multirow{4}{*}{$\pi^-$}   & 0-59    & 1.20E-11 & 3.54E-18 & 3.54E-29 & 1.95E-30 \\ % & 6.58E-13
                           & 59-97   & 3.59E-14 & 1.06E-20 & 1.06E-31 & 5.83E-33 \\ % & 9.38E-16
                           & 97-119  & 6.02E-19 & 1.78E-25 & 1.78E-36 & 9.78E-38 \\ % & 5.44E-19
                           & 119-180 & 2.29E-22 & 6.78E-29 & 6.78E-40 & 3.72E-41 \\ % & 2.14E-19
\hline                                            
 \multicolumn{3}{r|}{Sum (per POT)}              & 3.56E-18 & 3.56E-29 & 4.34E-22\\           
\hline
\end{tabular}
%\end{adjustwidth}
\caption{
Final estimated antiproton-induced background rates.
See text for column definitions.
\tablabel{bg:antiprotons:estimates}}
\end{table}\xspace}

\newcommand{\FigBgBeamMomVsTime}{
\begin{figure}[t]
\centering 
\subfloat[][\figlabel{bg:beam:MomVsTime:strawTrk}$e^-$ at the Straw Tracker]{
\includegraphics[width=0.45\textwidth,trim=0 0 0 1.7cm,clip]{figs/backgrounds/Beam_TimeVsMomentum_StrawTkt.pdf}}
\subfloat[][\figlabel{bg:beam:MomVsTime:stopTgt}$e^-$ after the Beam Blocker]{
\includegraphics[width=0.45\textwidth,trim=0 0 0 1.7cm,clip]{figs/backgrounds/Beam_TimeVsMomentum_StopTgt.pdf}}
\caption{
The timing and momentum of electrons detected at the straw tracker, \protect\subref{fig:bg:beam:MomVsTime:strawTrk}, and passing a plane immediately after the beam blocker, \protect\subref{fig:bg:beam:MomVsTime:stopTgt}.
\figlabel{bg:beam:MomVsTime}}
\end{figure}
}

\newcommand{\TabBgBeamFactors}{
\begin{table}[tb]
\centering
        \begin{tabular}{llm{0.5\textwidth}}
	\hline
        Parameter & \multicolumn{1}{l}{Value} & Description \\
	\hline
        $A_\textrm{geom}$ & \VarBeamBgGeometric &  Geometric acceptance to high-momentum electrons per \ac{POT}\\ 
        $A_\textrm{mom}$ & \VarBeamBgMomentum &  Fraction of beam electrons within the signal momentum cuts \\ 
        $A_\textrm{time}^\textrm{delayed}$ & \VarBeamBgTiming &  Fraction of beam electrons delayed until the detector window \\ 
	\hline
\end{tabular}
\caption{
Parameters and their values in the determination of the background rate due to high-energy particles in the beam.
\tablabel{bg:beam:factors}}
\end{table}
}

\newcommand{\FigBgBeamExtrapolate}{
\begin{figure}[t]
\centering 
\subfloat[][\figlabel{bg:beam:acceptance:momentum}Momentum of electrons]{
\includegraphics[width=0.49\textwidth,trim=0 1.5cm 0 0.7cm,clip]{figs/backgrounds/Beam_Acceptance_momentum.pdf}}
\subfloat[][\figlabel{bg:beam:acceptance:time}Time of high-$p$ electrons]{
\includegraphics[width=0.49\textwidth,trim=0 1.5cm 0 0.7cm,clip]{figs/backgrounds/Beam_Acceptance_time.pdf}}
\caption{
Projections of the momentum and arrival time of electrons.
\protect\subref{fig:bg:beam:acceptance:momentum} shows the momentum for electrons that arrive at any time, 
whereas \protect\subref{fig:bg:beam:acceptance:time} only shows the arrival time for electrons with momentum above 65~MeV/c.
Magenta lines in \protect\subref{fig:bg:beam:acceptance:time} indicate the mean arrival time after the beam blocker and the mean and maximum arrival time
at the straw tracker.
The black line is an exponential fit to the momentum and time of electrons after the beam blocker, with the equation and constants shown next to the fit.
%Below this momentum, electrons from muon decay-in-orbit dominate the timing.
\figlabel{bg:beam:acceptance}}
\end{figure}
}

\newcommand{\FigBgCosmicPrimary}{
\begin{figure}[t]
\centering
\subfloat[][\figlabel{bg:cosmics:primary:momentum}Momentum]{
\includegraphics[width=0.47\textwidth,trim=0 0 0 0,clip]{figs/backgrounds/Cosmics_primary_momentum.pdf}}
\subfloat[][\figlabel{bg:cosmics:primary:transverse}Transverse Direction Cosine]{
\includegraphics[width=0.47\textwidth,trim=0 0 0 0,clip]{figs/backgrounds/Cosmics_primary_transverse.pdf}}
\caption{
Distributions of momenta and the transverse direction for the cosmic muon and antimuon fluxes provided by the ND280 collaboration.
\figlabel{bg:cosmics:primary}}\end{figure}\xspace}

\newcommand{\FigBgCosmicBeam}{
\begin{figure}[t]
\centering
\includegraphics[width=0.98\textwidth,trim=2.5cm 0 2.3cm 0,clip]{figs/backgrounds/Cosmics_highP_electrons.png}
\caption{
Projection onto beamline coordinate system of electrons from cosmic rays with momenta greater than 100~MeV/c, from a simulation of about 170 million cosmic muons.
Three tracks can be seen to enter the electron spectrometer and pass along, although their momenta are all above 105~MeV/c as visible from the upwards drift of the trajectories.
\figlabel{bg:cosmics:beam}}\end{figure}\xspace}

\newcommand{\FigBgCosmicMomenta}{
\begin{figure}[t]
\centering
\includegraphics[width=0.9\textwidth,trim=0 0 1.6cm 0,clip]{figs/backgrounds/Cosmics_strawTrackHits.pdf}
\caption{
Momentum of electrons that hit one and two layers (planes) of the straw tracker.
The vertical axis shows the raw number of observed events based on the simulation of 170 million cosmic muons described in the text.
\figlabel{bg:cosmics:momenta}}\end{figure}\xspace}

\newcommand{\TabBgCosmicParameters}{
\begin{table}[t]
\centering
\begin{tabular}{llm{0.7\textwidth}}
	\hline
        Parameter & \multicolumn{1}{l}{Value} & Description \\
	\hline
        $R_\mu$                     & \num{4.17e5} & Cosmic ray flux through entire hall per second         \\ 
        $n_e^\textrm{obs}$          & 0.46         & Observed electrons per MeV/c                           \\ 
        $N_\textrm{sim}$            & \num{1.76e8} & Number of simulated events                             \\ 
        $(1-\epsilon_\textrm{CRV})$ & \num{1e-4}   & Miss rate of the \ac{CRV}                              \\ 
        $A_\textrm{duty}$           & 0.36         & Accelerator duty factor                                \\ 
        $A_\textrm{timing}$         & 0.51         & Fraction of time within gated-time detection window \\ 
        $A_\textrm{TDAQ,Recon}$     & 0.70         & Combined acceptance of TDAQ and Reconstruction         \\ 
	\hline
\end{tabular}
\caption{
Parameters and their values for the determination of the cosmic ray background rate.
\tablabel{bg:cosmics:params}}
\end{table}
}

\newcommand{\TabBackgroundFinalVals}{
\begin{sidewaystable}[b]
\centering
\sisetup{table-number-alignment=right,table-format=1.2e3}%
\renewcommand{\arraystretch}{1.5} % Default value: 1
%\begin{adjustwidth}{-0.7cm}{}%
\begin{tabular}{ll|SSS|S|L{7cm}}
\toprule
Type                         & Source                   & \multicolumn{3}{c|}{Background Rate}                                       & \multicolumn{1}{c|}{Total Events} & Comment \\ 
                             &                          & {per $\mu^-$ stop} & {per POT}              & {per second}            &            \\ 
\toprule
\multirow{2}{*}{Intrinsic}   & DIO                      & \NumDIOPerMuStop & \NumDIOPerPOT                    & \NumDIOPerSecond                 & \NumDIOTotal                     & \\ 
			     & RMC                      & \NumRMCPerMuStop & \NumRMCPerPOT                    & \NumRMCPerSecond                 & \NumRMCTotal                     & \\ 
\cmidrule{2-7}
\multirow{4}{*}{Delayed}     & RPC                      & {--}             & \NumRPCDelayedPerPOT             & \NumRPCDelayedPerSec             & \NumRPCDelayedTotal              &  \\ 
			     & Beam                     & {--}             & \NumBeamBgDelayedPerPOT          & \NumBeamBgDelayedPerSec          & \NumBeamBgDelayedTotal           & \multirow{1}{7cm}{\footnotesize Beam includes high-energy electrons from $\pi$, $\mu$, and $n$ capture.} \\ 
			     & Stopped $\bar{p}$        & {--}             & \NumBgAntiprotonsPerPOT          & \NumBgAntiprotonsPerSec          & \NumBgAntiprotonsTotal           & \multirow{2}{7cm}{\footnotesize Based on conservative interpolation and extrapolation of limited experimental $\bar{p}$ data.}\\ 
                             & $\pi^{-}$ from $\bar{p}$ & {--}             & \NumBgAntiprotonsPiDelayedPerPOT & \NumBgAntiprotonsPiDelayedPerSec & \NumBgAntiprotonsPiDelayedTotal  & \\ 
\cmidrule{2-7}
\multirow{3}{*}{Prompt}      & RPC                      & {--}             & \NumRPCPromptPerPOT              & \NumRPCPromptPerSec              & \NumRPCPromptTotal               &  \\ 
                             & Beam                     & {--}             & \NumBeamBgPromptPerPOT           & \NumBeamBgPromptPerSec           & \NumBeamBgPromptTotal            &  \\ 
                             & $\pi^{-}$ from $\bar{p}$ & {--}             & \NumBgAntiprotonsPiPromptPerPOT  & \NumBgAntiprotonsPiPromptPerSec  & \NumBgAntiprotonsPiPromptTotal   & \\ 
\cmidrule{2-7}
\multicolumn{2}{c|}{Cosmics}                            & {--}             & {--}                             & \NumCosmicRatePerSecond          & \NumCosmicRateTotal              & \footnotesize Dominated by conservative miss-rate. \\ 
\bottomrule                                                                                                   
\multicolumn{2}{c|}{Total}                              & {--}             & {--}                             & \NumTotalBgPerSecond             & \NumTotalBgPhasII                &  \\ 
\bottomrule
\end{tabular}
%\end{adjustwidth}
\caption{%
Final predicted background rates and events.
The dominant backgrounds are cosmic ray electrons and stopping antiprotons, which are predicted to have produce almost equal rates.
Assumes an extinction factor of \VarExtinctionFactor, a double-Gaussian momentum resolution with a core resolution of $\sigma=200$~keV/c and a tail probability of \VarResolutionProbFromRMCToSignal that events are misreconstructed with at least an additional 2.4~MeV/c. 
All other parameters are the same as stated in chapter~\sect{sense}.
 \tablabel{bg:summary}}
\end{sidewaystable}%
}

\newcommand{\FigBgVsResolution}{
\begin{figure}[t]
\centering
\includegraphics[width=0.9\textwidth,clip=true,trim=0 0 0 1cm]{figs/backgrounds/BG_vs_resolution.pdf}
\caption{
Dependence of the intrinsic background rates on $\sigma_2$, the width of the high/low momentum tail in the reconstruction distribution.
From this it can be seen that the background rates from DIO will always dominate, compared to RMC.
\figlabel{bg:resolution}}\end{figure}\xspace}

%% Important variables:
% Number of POT
% Number of muons stopped
% Total run time
\chapter{COMET \phaseII: Backgrounds}
\sectlabel{backgrounds}
Having optimised and evaluated the signal sensitivity, it is now important to check the expected background rates.
Interpreting the final result is only possible if the number of background events has been predicted.
The certainty of any observation or the stringency of a final confidence limit is determined by the relative background rate and signal sensitivity; ideally the background rate should be small, well below a single expected event.

The types of background that must be considered for COMET were outlined in \sect{detector:background}, but in this section they will be evaluated for the optimised experiment design using the improved simulation.

\FigDIOBackground
\section{Muon Decay in Orbit (\acs{DIO})}
\sectlabel{bg:dio}
%\begin{easylist}
%# Discussed in previous section \sect{sense:momentum}
%# DIO spectrum near end-point: which expansion is correct?
%# Final DIO rate above momentum threshold with resolution
%\end{easylist}
In order to decide the optimal threshold for the momentum cut, it was necessary, in section \sect{sense:momentum}, to study the rate of \ac{DIO} events.
With a low-momentum threshold of \VarMomThreshold, the number of expected DIO events per muon stop is \VarDIOPerMuStop{}.
Given the \VarTotalMuStops~muon stops that should take place during \phaseII, the total expected number of background events due to DIO is \VarDIOTotal.

However, it is important to note just how steeply falling the DIO rate is in this region.
\Fig{bg:dio:rates} shows how, given a fixed run time, the number of DIO background events is affected by changing the low-momentum threshold.
Changing the momentum threshold also affects the signal acceptance, so that for a fixed run time the \ac{ses} is also affected.
Simultaneously, with more muons stopped the \ac{DIO} rate will increase for a fixed momentum threshold.
\Fig{bg:dio:rates} was produced using the same DIO and signal spectra shown in \fig{sense:spectra:resolution}, which included energy losses in the material of the target, beamline, and detector, as well as a 200~keV/c Gaussian resolution function. 
In addition, all the acceptance parameters of the sensitivity chapter were held fixed, except for the momentum cut efficiency.

\FigDIOEndPointComparison
The fact that the run-time, signal acceptance, and background rate depend so strongly on the momentum threshold makes the theoretical prediction for the \ac{DIO} spectrum particularly important.
The two most recent calculations of the high-energy tail of the DIO spectrum are shown in \fig{bg:dio:spectra}, as well as the cruder function that fits the whole spectrum range, used in SimG4.
The more conservative spectrum from the 2011 paper~\cite{Czarnecki2011} has been used in this study.
However, if one were to use the more recent 2015 spectrum~\cite{Czarnecki2014}, background rates due to DIO would fall to 0.057 events at the same momentum threshold.
This would agree with that paper's statement that the inclusion of radiative corrections suppresses the DIO background by 15\% at the end-point.

In total, based on the 2011 DIO spectrum, \sci{3}{-18} electrons would be produced per muon stop with momentum greater than \VarMomThreshold.
Therefore, the DIO rejection efficiency, including both geometric effects and the threshold on the detected momentum, suppresses the detection of DIO end-point electrons by about 98\%.

\section{Radiative Muon Capture (\acs{RMC})}
\TabRMCEndPoints%
\sectlabel{bg:rmc}
During the process of nuclear muon capture, there is a finite probability of a hard photon being radiated from the muon, nucleus, or the exchanged $W$-boson.
This is known as \acf{RMC} and is distinguished from radioactive gamma-ray production during nuclear de-excitation or decay of the daughter nucleus.
The maximum energy this photon can take, $\max(E_e^\textrm{RMC})$, differs from the \mueconv signal and \ac{DIO} end-point by the minimum energy needed to change the nucleus from $N(A,Z)$ to $N(A,Z-1)$:
\begin{align}
\max(E_e^\textrm{RMC})=&\big(M_\mu - B_{\mu,\textrm{binding}} - E_\textrm{rec}\big)-\big(M_{N_2}+\sum_iM_{h_i}-M_{N_1}-\Delta{}Z\cdot{}M_e\big),\\
                      =&E^\textrm{Conversion}_e  - \Delta{}M,
\end{align}
with $M_{N_1}$ and $M_{N_2}$ the mass of the parent and daughter nuclei respectively, and $M_{h_i}$ the mass of the $i$-th hadron (proton, neutron, alpha, \etc).
The mass of the electron, $M_e$, must also be included since a free electron must be ejected from the atomic orbitals for every proton removed from the nucleus, equivalent to $\Delta{}Z$.
$\Delta{}M$ is then the total energy lost to changes in the atomic mass%
\footnote{
One has to think of the atomic mass difference, rather than the nuclear mass difference.
The distinction is important since consideration of the nuclear masses alone would ignore the effect on the atomic electrons.
Under nuclear muon capture (radiative or otherwise), the number of protons in the nucleus is reduced by at least one, and accordingly atomic electrons become unbound.
The notation here follows that of the COMET TDR~\cite{TDR2016} and CDR~\cite{CDRphase2}, where the electron mass is absorbed into the value of $\Delta{}M$ (although there it is called $\Delta_{Z-1}$).
}
of the nuclei and other emitted hadrons.

When an aluminium-27 nucleus captures a muon various daughter nuclei are possible.
If no other particles are emitted, as part of direct nucleon capture, the daughter nucleus will be magnesium-27.
In general, this could be left in an excited state, but to reach the end-point of the \ac{RMC} spectrum it will be left in the ground-state configuration.
The atomic mass difference between these two nuclei is shown in \tab{bg:rmc:massDifferences}, where it can be seen the RMC end-point is separated from the \mueconv signal energy by around 3~MeV.

\FigRMCExperiments
%\newcolumntype{d}[1]{D{.}{\cdot}{#1} }

\Fig{bg:rmc:experiments} shows the summary table of experimental data on Aluminium, taken from the summary by Gorringe~\cite{RevModPhys.76.31}.
It is interesting to note that in all experiments to date none of the empirical fits to RMC have suggested an end-point above 90.1~MeV~\cite{PhysRevC.37.1633,PhysRevC.46.1094,PhysRevC.59.2853}.
This is close to the end-point predicted by the various transitions besides ${}^{27}$Al$(\mu,\gamma\nu){}^{27}$Mg.
The two most recent of these experiments both measure the branching ratio for \ac{RMC} producing photons with $E>57~$MeV to ordinary muon capture to be \sci{1.43}{-5}. 

To produce a background event in the detector, the high-energy photons produced from RMC must be converted to a high-energy electron.
Asymmetric pair production from this photon is one mechanism by which this can occur, although producing the positron at rest is highly suppressed.
In addition, around 1~MeV is consumed by the mass of the electron-positron pair.
Compton scattering provides another mechanism by which the high-energy photon can convert to a high-energy electron.
In the limit where the photon is reflected directly back on itself, the resultant electron has momentum greater than the incoming photon by around $M_e/2$.
As such it is Compton scattering which is more of a concern for COMET.

Finally, given the maximum photon energy is below the momentum threshold of \VarMomThreshold, the electron's energy needs to be mis-reconstructed by around 2.4~MeV.
If the final resolution function were a pure Gaussian with a 200~keV/c width (at signal energies), then mis-reconstruction by such an amount would be an 11~$\sigma$ event, \ie $P(p_\textrm{recon}-p_\textrm{true}>\SI{2.4}{MeV/c})=\num{1.9e-28}$.
However, the resolution function is not likely to be a pure Gaussian.
On-going studies for reconstruction with the  \phaseI StrECAL currently suggest that the tails of the resolution function are much larger than for a pure Gaussian.
Although the \phaseI Straw Tracker reconstruction is still being finalised, early studies suggest 
$P(p_\textrm{recon}-p_\textrm{true}>\SI{2.4}{MeV/c})=\VarResolutionProbFromRMCToSignal$~\cite{YFujiiStrECALRecon}.
Although this is considerably larger than the pure Gaussian, this is based on the most preliminary results of reconstruction algorithms.
If need be, cuts on fit quality can be improved over the current \phaseI values; the sensitivity estimate presented here should be robust against this since it includes the CDR estimate for the signal efficiency of such a cut.

\subsection{Calculation and Simulation of RMC}
Conversion of the RMC photons can take place in any of the material around the target, such as the beam blocker, solenoids, or cryostat.
To estimate the acceptance of the electrons produced from this conversion a simulation was performed where RMC photons were input at the stopping target, using the realistic muon stopping distribution obtained previously.
Geant4 implements the process of photo-conversion via both of the afore-mentioned methods.

To build a realistic end-point spectrum, the same recipe as used in the \phaseI TDR was applied~\cite{TDR2016}.
The spectrum shape near the end-point is modelled by the equation~\cite{CHRISTILLIN1980331}:
\begin{equation}
	\eqlabel{bg:rmc:spect}
	\Gamma(\textrm{RMC})\propto(1-2x+2x^2)(x)(1-x)^2,
\end{equation}
where $x=E_\gamma/\max(E_\gamma^\textrm{RMC})$.
The largest observed branching ratio for \ac{RMC} compared to ordinary muon capture---\num{1.83e-5}---is used to set the normalisation.
The relative rate of \ac{RMC} resulting in photons with energies above 90~MeV is, based on the spectrum in equation~\eq{bg:rmc:spect}, \num{7.7e-3} of the experimentally observed rate.
\FigRMCSimResults

Given the total number of muon stops during \phaseII, the probability of ordinary muon capture, and the relative rates for \ac{RMC}, one expects some \VarRMCTotalPhotonPhaseII  photons to be produced in the target with energies above 90~MeV.
Since this is an intrinsic background, the lifetime of this process will be the same as for signal events. 
Assuming, therefore, the same timing window acceptance, and applying the same reconstruction and TDAQ efficiencies as for signal, this number is reduced to \VarRMCTotalPhotonPhaseIIAccepted.

To check the probability that these photons produce a high-momentum electron in the detector, 
some \num{6e7} RMC events were generated in the stopping target, with initial photon energies greater than 90~MeV and distributed according to the spectrum in equation~\eq{bg:rmc:spect}.
\Fig{bg:rmc:simulation} shows the momentum and rate at which these were detected.
Based on this, the fraction of events reaching the detector with momentum larger than 98~MeV/c was \num{6.2e-7}, whilst above 100~MeV/c only \num{3.3e-8} electrons per \ac{RMC}$|_{p>90}$ event were detected.

The total number of electrons arriving with momentum larger than 98~MeV/c, therefore, is \VarRMCTotalElectronPhaseIIAccepted.
Since the resolution function is somewhat uncertain, this number can be interpreted as setting the limit for the high-momentum mis-reconstruction tail.
To have fewer than 0.01 background events produced by \ac{RMC}, the tails of the resolution distribution should contain less than $0.01/\VarRMCTotalElectronPhaseIIAccepted = \scinot{0.01/\NumRMCTotalElectronPhaseIIAccepted}$ events.

Using the estimate of the resolution function's tail from \phaseI, mentioned above, we predict therefore that \phaseII will produce \VarRMCPerMuStop per muon stop or \VarRMCTotal during the entire \phaseII run.

%As a result, if the current \phaseI resolution were applied, we would expect around 8~background events during the entire run.
%However, since the resolution function at this stage is so poorly known, it seems too premature to define a value.
%Instead, the predicted RMC rate provides a constraint on the high-energy tail of the resolution function: to do better than 0.1 background events during \phaseII from RMC, fewer than 1 in \num{2e4} high-energy electrons can be reconstructed with a momentum more than 2.4~MeV/c larger than the true value.
%On the assumption that such a resolution is achievable, we predict here that the number of backgrounds events due to RMC is less than \VarRMCPerMuStop per muon stop, or \VarRMCTotal during the entire \phaseII run.

\subsection{Aluminium-26 and \acs{RMC}}
Based on the above energy calculation, the end-point for \ac{RMC} against ${}^{26}$Al (to ${}^{26}$Mg and with no other particles emitted) would be 108.5~MeV.
Clearly such photons would be extremely dangerous to COMET if they are produced.

Aluminium-26 comes in two isomers, one with a half-life of around 6 seconds~\cite{PhysRevLett.106.032501}, the other lasting around 700 thousand years~\cite{AUDI20033}.
Since Al-26 is unstable its abundance in natural aluminium is low.
However, it can be produced by various methods, such as proton and deuteron bombardment of magnesium and sodium, or photo-neutron emission of the aluminium-27 isotope~\cite{THOMPSON1965486}.
With such production mechanisms, it is likely that aluminium-26 will be produced in the COMET stopping target, via the interaction of daughter nuclei of muon capture (which are typically magnesium and sodium)
with protons, deuteron, gammas, and neutrons coming from either the beam or as products of muon capture.
The exact rate of Al-26 production, however, is a complicated value to estimate and one that unfortunately cannot be estimated here.

One can, however, set an acceptable rate of production if the induced background event rate is to be kept at the level of 0.1 events.
Based on the previous simulation of Al-27 \ac{RMC}, the probability of an electron being detected within 5~MeV of the \ac{RMC} end-point is \num{5.6e-18} per nuclear muon capture.
With \VarTotalMuStops muon stops during \phaseII, and assuming that the branching ratios for both ordinary and radiative muon capture are the same for Al-26 and Al-27 (61\%), then the
concentration of Al-26 in the stopping target must be less than 1\% (by number density) on average during the entire \phaseII run.
The branching ratio for radiative muon capture is, in reality, likely to be slightly more than in Al-27, based on the fact that the neutron excess of a given isotope ($N=A/2-Z$) seems a better indicator of the branching ratio than atomic number, and that these two numbers appear anti-correlated~\cite{RevModPhys.76.31} ( the greater the neutron excess, the lower the RMC branching ratio).

Clearly though, there is more work to be done on this, including a better understanding of \ac{RMC} events coming from Al-26 and of the rate of Al-26 production from muon beams.
Data from the \alcap experiment might be able to help with this, however, since the production rate could depend a lot on the exact beam conditions, measuring this on site is important.
This should therefore be one of the measurements prioritised at \phaseI.

Finally, if this does produce a sizeable background contribution, one can imagine several techniques to mitigate or reduce the challenge this poses.
Provided one can measure the concentration of Al-26 in the stopping target at the end of the run, and ideally at various stages whilst running, then the number of backgrounds can be predicted and possibly subtracted.
Additionally, it could be possible to remove the stopping target and replace it with a fresh one such that the Al-26 concentration never rises beyond an unsafe level.

%\CHECK{Mention that this could be a useful source of calibration}
\section{Radiative Pion Capture (\acs{RPC})}
\sectlabel{bg:rpc}
When low-energy negative pions are stopped in material they behave similarly to negative muons and form pionic atoms.
The probability that the pion is then captured by the nucleus rather than decays in orbit is, however, considerably larger than for a muon.
Furthermore, given the extra 30~MeV/c$^2$ of the pion mass and the lack of an outgoing neutrino, the end-point for \acf{RPC} by the nucleus is well above the \mueconv signal energy.
As for photons of \ac{RMC}, the \ac{RPC} photons can then be converted via Compton scattering or pair production to signal-like electrons.
Pion capture could, therefore, be a dangerous source of backgrounds and was one of the dominant sources of background events at \sindrumII, along with cosmic sources~\cite{sindrum2006}.

Since the pion interacts via the strong force, negative pions capture almost immediately in the nucleus, on the order of picoseconds~\cite{Engelhardt:1975ct}.
The timing of backgrounds caused by pion capture are therefore determined predominantly by the time of pion production, and the delay due to transportation of the pion and signal-like electron.
If a background arises from pion capture and the pion was produced in the main muon pulse, then the pion or resultant background electron must have been significantly delayed.
Delayed \ac{RPC} backgrounds are therefore suppressed by the time-gated detector window.
If however a background arises because the pion was produced outside of the main proton pulse, due to late-arriving protons (or from antiprotons, but we will treat these separately below), then this background is considered prompt.
These prompt \ac{RPC} background events are, therefore, suppressed by the extinction factor.

The background rate per \ac{POT} for prompt and delayed pions is therefore:
\begin{align}
	R_\textrm{delayed}=&R_{\pi/p}\mathcal{B}_\textrm{RPC}f_{e,\textrm{RPC}}A_\textrm{time}, \\
	R_\textrm{prompt}=&R_{\pi/p}\mathcal{B}_\textrm{RPC}f_{e,\textrm{RPC}}\epsilon_\textrm{extinction},
\end{align}
where $R_{\pi/p}$ is the pion stopping rate per \ac{POT} and $\mathcal{B}_\textrm{RPC}$ is the branching ratio of \ac{RPC} for stopped pions.
$\epsilon_\textrm{extinction}$ is the extinction factor, whereas $A_\textrm{time}$ is the acceptance of the time-gated detector window to electrons coming from delayed RPC.
$f_{e,\textrm{RPC}}$ is the probability that an RPC photon converts to an electron which reaches the detector with signal-like momentum.
To a reasonable approximation, this can be factorised as:
\begin{equation}
f_{e,\textrm{RPC}}=f_{\gamma\rightarrow{}e^-}A_\textrm{geom}A_\textrm{mom}
\end{equation}
where $f_{\gamma\rightarrow{}e^-}$ is the conversion rate of RPC photons to an electron, $A_\textrm{geom}$ and $A_\textrm{mom}$
are, respectively, the geometric acceptance and momentum cut efficiency for such electrons.
However, such a factorisation misses out various correlations, such as where in the experiment the conversion takes place and at what momentum the secondary electron is produced, so that only a single value for $f_{e,\textrm{RPC}}$ will be reported here.

\FigRPCData
\subsection{Photons from Radiative Pion Capture (\acs{RPC})}
%In much the same way as \ac{RPC}, when a pion is captured in a nucleus there is a chance that it will produce a prompt, hard photon.
There is a range of experimental and theoretical data on \ac{RPC}. 
\Fig{bg:rpc:data} shows what is perhaps the most useful data currently available: the observed spectrum of photons coming from RPC for magnesium and calcium.
Magnesium being adjacent to aluminium on the periodic table, this spectrum is a reasonable proxy for the spectrum of \ac{RPC} on aluminium.
The relative rate for \ac{RPC} compared to ordinary pion capture is discussed in Amaro \etal~\cite{Amaro:1997ed}.
For experimental and theoretical studies of the three isotopes summarised in that paper---carbon, oxygen, and calcium---the measured and predicted branching ratios are all within 1.19\% and 2.27\%.
To be conservative, we take here the branching ratio for RPC on aluminium to be the largest of these at 2.27\%.

\subsection{Pion Stopping Rate}
\FigPionStopDist
\FigPiVsMuMomenta
To simulate the pion stopping rate, the pions from the main production simulation were resampled multiple times to build up a large number of pion stops.
\Fig{bg:piStop:dist} shows the distribution of pions stopping in the target, in one-dimensional projections to the ICEDUST global coordinate system.
In that coordinate system, pions arrive from large values of Z at the target, so the pion beam in \fig{bg:piStop:dist:z} is going from right to left.
By comparison with the plots in \fig{sense:stops}, pions tend to stop further downstream in the target.
This is readily understood by the fact that pions reaching the stopping target tend to have much higher momentum than muons reaching the target, as shown in \fig{bg:piVsMu:momenta}.

In total, the number of pions stopping in the target per \ac{POT} is $R_{\pi/p}=$\VarPiStopsPerPOT.

\subsection{Simulating \acs{RPC}}
\FigRPCSimulatedSpectrum
\FigRPCSimResults
Using the realistic pion stopping distribution shown in \fig{bg:piStop:dist}, \ac{RPC} photons were generated in the target and Geant4 used to convert and track electrons resulting from this process.
To model the distribution of photon energies from \ac{RPC}, the experimentally obtained spectrum from magnesium of ref.~\cite{Bistirlich:1972jy} was used.
To build the model, the raw spectrum was first digitised and then smoothed, using `TGraphSmooth::SuperSmooth()' from the  ROOT library~\cite{ROOT}.
These steps are shown in \fig{bg:rpc:spectrum}.

Based on a simulation of \num{4e6} RPC photons, the distribution of electrons and positrons reaching the detector was obtained.
The timing and momenta of such electrons is shown in \fig{bg:rpc:sim}, where it is clear that although many signal-like electrons are detected, they all arrive well before the gated-time threshold of 600~ns.

The probability of an \ac{RPC} photon producing an electron that reaches the detector with momentum between \VarMomThreshold and 105~MeV/c is: $f_{e,\textrm{RPC}}=$\VarDetectedEsPerRPC.
To estimate the rate of delayed \ac{RPC} backgrounds, we also need to know the value of $A_\textrm{time}$.
As can be seen in \fig{bg:rpc:sim:momVtime}, the timing for electrons originating from RPC photons and detected with momentum greater than 30~MeV/c is independent of the momentum.
By fitting the tail of this distribution with a single exponential, the lifetime of the high-energy electrons is found to be \VarRPCLifetime, such that with the timing window between 600 and 1200~ns we find the acceptance of the timing window to be: $A_\textrm{time}=\VarRPCTimingEfficiency$.

\Tab{bg:rpc:estimates} summarizes these numbers, from which we find that the rate of backgrounds for delayed \ac{RPC} is \VarRPCDelayedPerPOT per POT, whilst prompt RPC occurs at \VarRPCPromptPerPOT per POT.
\begin{table}[tb]
\centering
        \begin{tabular}{llm{0.5\textwidth}}
	\hline
        Parameter & \multicolumn{1}{l}{Value} & Description \\
	\hline
        $R_{\pi/p}$                    & \VarPiStopsPerPOT & Pion stopping rate per \ac{POT}  \\ 
        $\mathcal{B}_\textrm{RPC}$     & \num{2.27e-2} & Branching ratio of \ac{RPC} \\ 
	$f_{e,\textrm{RPC}}$           & \VarDetectedEsPerRPC & Probability of an RPC photon producing signal-like electrons in the detector \\ 
	$A_\textrm{time}$              & \VarRPCTimingEfficiency & Acceptance of timing window to secondary electrons from RPC \\ 
        $\epsilon_\textrm{extinction}$ & \VarExtinctionFactor[2] &  Extinction factor\\ 
	\hline
\end{tabular}
\caption{\tablabel{bg:rpc:estimates}
Parameters and their values in the determination of the \ac{RPC} background rate.
}
\end{table}

\section{Antiprotons in the Beam}
\sectlabel{bg:antiprotons}
When the primary 8~GeV proton beam interacts with the production target, antiprotons might be produced as a proton--antiproton pair:
\begin{equation}
p + N(A,Z) \rightarrow p + N^*(A,Z) + p+\bar{p},
\end{equation}
which, given their negative charge, can remain on-axis along the muon beamline if they have momentum around 40~MeV/c.
Because their mass is relatively large, at a fixed value of momentum antiprotons travel much more slowly than other particles in the muon beam, causing the beam's time structure to be smeared for antiprotons and their secondaries.
The pulsed beam and time-gated detector window are, therefore, not effective at suppressing backgrounds created by antiprotons in the beam.

%Secondaries of antiproton interactions with matter include X-rays and pions, though pion production dominates~\cite{BobBernsteinExperimentersGuide,Mishustin:2004xa}.
%
The interaction of antiprotons with matter has a strong analogy with muons and negative pions, particularly at low energies, when the antiprotons stop.
%There is a strong analogy between antiproton stopping in matter and the phenomona around negative muon and pion stopping.
In matter, antiprotons with energies of a few tens of keV---similar to that of atomic electrons---can become bound in the Coulomb potential of the atom's nucleus.
X-rays emitted in the ensuing electromagnetic cascade are typically not more than a 100~keV~\cite{Aramaki201352}.
Unlike for muons, the antiproton will only rarely reach the atomic ground state, with interactions with the nucleus becoming dominant before this is possible.
Such interactions can take the form of an immediate annihilation or the formation of a composite nucleus where the antiproton becomes bound within the nuclear potential~\cite{Wong:1984fy,Mishustin:2004xa}.
The binding energy of some of these nuclear levels can reach up to hundreds of MeV, and so in the transition, pions and other hadrons can be readily produced.
Eventually the antiproton annihilates in the nucleus, producing large multiplicities of pions and other mesons.
Based on the data in ~\cite{Mishustin:2004xa}, about 1.56 negative pions should be produced for every stopped antiproton.

As such, although antiprotons themselves are not an immediate source of high-energy photons or electrons, they are an additional source of pions, which can produce backgrounds via \acf{RPC}.

\subsection{Antiproton Production Rate and Spectrum}
The literature on antiproton production with 8~GeV protons on a tungsten target and at large angles is somewhat lacking.
%There is really very little literature on the production of antiprotons from a tungsten target for a range of angles.
Accordingly most hadron models are particularly under-constrained when it comes to antiproton production.
The QGSP_BERT_HP model used as the basis for SimG4 is, in fact, completely unable to produce antiprotons.

\FigAntiprotonMeco
In the COMET TDR~\cite{TDR2016}, the yield of antiprotons per \ac{POT} is given as \num{4e-5}, based on the findings of a literature review and simulations undertaken for the MECO experiment~\cite{Meco024}.
\Fig{bg:antiprotons:meco24} shows the production cross section for antiprotons as a function of incident proton momentum, based on that study.
%The dependence of the antiproton production cross section on incident proton momentum based on that study is shown in \fig{bg:antiprotons:meco24}.

Simulations of antiprotons for \phaseI have used this information and shown that absorber foils would be needed along the beamline.  
Since the \phaseII geometry in ICEDUST re-uses most of the implementation for the Production Target Capture and Torus1 sections, the upstream absorbers are also contained in the geometry used to study antiprotons here,
although the absorber near 90\degree is not.

\FigAntiprotonData
Whilst tungsten targets have not been studied at the relevant angles and proton energies, a set of papers \cite{Boyarinov:1994tp,Kiselev:2012sj} do exist covering antiproton production up to 2 radians for tantalum (which is adjacent to tungsten on the periodic table), copper, aluminium, and beryllium targets and using protons with 10.14~GeV kinetic energy.
Between them, these two papers provide the invariant triple-differential cross section as a function of antiproton momentum for production angles of 10, 59, 97 and 119\degree, defined as:
\begin{equation}
\eqlabel{bg:antiprotons:lit}
	F(p,\theta,\phi)=E \frac{d^3\sigma}{dp^3}=E\frac{d^3\sigma(p,\theta,\phi)}{p^2dpd\Omega}
\end{equation}
where $E$ and $p$ are the antiproton energy and momentum.  The earlier Boyarinov paper actually reports $f=F/A$, where $A$ is the relative atomic mass of the nucleus.
To convert the differential invariant cross sections given in the literature (equation~\eq{bg:antiprotons:lit}) into a differential production rate per \ac{POT}, $d^3R(\theta)/dpd\Omega$, the following formula is used:
\begin{align}
	\eqlabel{bg:antiprotons:rate}
	\frac{d^3R(p,\theta,\phi)}{dp^3}&=\frac{F(p,\theta,\phi)}{E}\frac{\rho N_A l}{m_N}
\end{align}
where $\rho$ is the mass density of the target, 19.25~g/cm$^3$ for Tungsten, $m_N$ the atomic mass, 183.86~g/mol for tungsten, $l$ the length of the target, and $N_A$ is Avogadro's number.
\Fig{bg:antiprotons:data} shows the data from these papers, after scaling the cross sections for each material to the number of nucleons in tungsten.
The data has also been converted to a differential rate of production, given the \phaseII production target dimensions, using equation~\eq{bg:antiprotons:rate}.
From this data, it is clear that as one moves to larger angles, the spectrum becomes considerably softer whilst the overall rate falls quickly, in much the same way as for pion production.

\subsection{Modelling Antiproton Production}
To develop the data in the Boyarinov~\etal papers into a model for antiproton production, three dependencies must be considered:
the incident proton's momentum, the outgoing antiproton's momentum, and the outgoing antiproton's direction with respect to the incident proton.

To account for these dependencies, the cross section was factorised according to the equation:
\begin{align}
\eqlabel{bg:antiprotons:factorisation}
	\frac{d^3R(p,\theta,\phi)}{dp^3}=&\xi{}P(p,\theta_\mathrm{min})\Phi(\theta,\phi),
\end{align}
where $\xi$ is a scale factor to correct for the different incoming proton energy, $P(p,\theta_\textrm{min})$ models the momentum distribution at a given angle and is normalised to unity (\ie $\int{}P(p,\theta_\textrm{min})dp=1$), and $\Phi(\theta,\phi)$ encapsulates the dependence on the antiproton's direction.  Each of these factors are discussed in-depth below.

\subsubsection*{Dependence on the Incident Proton's Momentum: $\xi$}
The factor $\xi$ is included in the antiproton production model to account for the difference in proton energy between the available data and the COMET proton beam.
It is defined as the ratio of the total production cross sections at the different energies:
\begin{equation}
\xi=\frac{\sigma_\mathrm{tot}(E_p=8 GeV)}{\sigma_\mathrm{tot}(E_p=10.14 GeV)}
\end{equation}

The data in the Boyarinov papers was taken using protons with kinetic energies of 10.14~GeV.
The threshold for antiproton production against a single, free nucleon is 5.6~GeV; the cross-section increases rapidly above this energy, as indicated in \fig{bg:antiprotons:meco24}.
Because of this, COMET has chosen to use an 8~GeV (kinetic energy) beam, and accordingly it is important to scale the antiproton cross sections from the Boyarinov~\etal papers to account for this.

\Fig{bg:antiprotons:meco24} suggests that the cross section at 8~GeV is about one order of magnitude smaller than at 10.14~GeV, although the data and models included in that plot are for a lead target as opposed to tungsten.

In a paper by Danielewicz~\cite{PhysRevC.42.1564}, a theoretical calculation is performed to estimate the rate of antiproton production.
Equation (4.3) of that paper gives the antiproton production cross section for nucleon-proton interactions as:
\begin{equation}
\sigma_{NN}^{\bar{p}}=C(\sqrt{s}-4m)^{7/2}
\end{equation}
where $\sqrt{s}$ is the Mandelstam variable for the total invariant incoming energy, and $m$ the mass of a nucleon.
$C$ depends on properties of the nucleus and is fixed in that paper using experimental data.
Based on this equation, the cross section for antiproton production at 8~GeV should be a factor of 0.12 times that at 10.14~GeV.
Since this agrees well with \Fig{bg:antiprotons:meco24}, we take $\xi=0.12$.

\subsubsection*{Dependence on the Outgoing Antiproton's Momentum: $P(p,\theta_\textrm{min})$}
The momentum spectrum of antiprotons at a given angle is captured in the factor $P(p,\theta_\textrm{min})$.
We base the shape of this spectrum on the data from the Boyarinov papers.

\FigAntiprotonEndpoint
In order to complete the data and cover the full range of energies at which antiprotons might be produced, we add two additional points to the experimental data for the maximum and minimum momentum.
Including the minimum momentum only adds the constraint that at zero momentum the cross section also be zero.
On the other hand, to calculate the maximum momentum, the kinematic end-point was found by considering the entire nucleus and the two outgoing protons to recoil coherently against the antiproton.
The value of the end-point kinetic energy and longitudinal and transverse momenta for both tungsten and carbon are shown in \fig{bg:antiprotons:end-point} using the formulae derived in appendix~\sect{appendix:antiprotonEndpoint}.
This end-point will be a highly conservative estimate, since in reality not all the nucleus will recoil coherently;  the de Broglie wavelength for a proton with 8~GeV kinetic energy is about 0.15~fm, compared to the 7~fm or so of a tungsten nucleus%
\footnote{Antiproton production from protons in this energy regime was historically referred to as `sub-threshold production', since it is close to the threshold for single nucleon interactions, around 5.6~GeV.
The fact antiprotons are observed below this energy is therefore proof that it is not a single, stationary nucleon that interacts with the proton; the nuclear environment as a whole must be considered.
Either the nucleons are, themselves, in motion or clusters of bound nucleons are involved.
  The older literature on this topic in fact refers to a parameter called the `cumulative number' which was related to the number of nucleons that would need to be involved to produce secondary hadrons (\eg antiprotons or pions) with such outgoing energies, given the incoming proton energy.  
These sub-threshold secondary particles were themselves sometimes called cumulative particles.}.
Additionally, achieving this end-point configuration would be highly phase-space suppressed.

\FigAntiprotonFits
With the addition of these two end-points, to interpolate and extrapolate the data, a polynomial of order 4 to 6 (depending on the number of available data points) was used to fit from zero up to the last experimental data-point.
For the high-momentum tails of each spectrum, three fits were tried: a straight line fit between the last experimental data-point and the kinematic end-point described above; an exponential fitted to the kinematic end-point and the last two experimental data points; and an average of both of these fits.
\Fig{bg:antiprotons:fits} shows the results of this fitting procedure, where it can be seen that a low-momentum peak is visible and a high-momentum tail well described.
Whilst these spectra are likely a poor representation of the true production spectrum, they serve as useful upper bounds which can be used as inputs for the antiproton background rate estimation.

The momentum distribution will, of course, depend on the angle of production relative to the proton beam.
With only four production angles available in the experimental data, however, for a given angle, $\theta$, we assign the momentum distribution based on the data at the next smallest production angle, $\theta_\textrm{min}$.
Since the spectrum becomes softer at larger values of $\theta$, this will tend to further overestimate the high-energy tail of the momentum distribution.
In the case of angles smaller than 10\degree, given the lack of other data, we fix the  shape of the momentum distribution to be the same as that at 10\degree.

The function $P(p,\theta_\textrm{min})$ is the piecewise function that is fit to the data and kinematic end-points and then normalised so that the integral is one.
The overall normalisation is instead absorbed into the angular dependence factor, $\Phi(\theta,\phi)$.

\subsubsection*{Dependence on the Outgoing Antiproton's Direction: $\Phi(\theta,\phi)$}
\FigAntiprotonAngularDependence
The function $\Phi(\theta,\phi)$ accounts for the variation of the cross section as the angle between the incoming proton and outgoing antiproton changes.
By symmetry, it can be seen that there should be no dependence on $\phi$ (both the proton beam and target are assumed to be unpolarised).

To estimate the cross section at a given angle, the fitted distributions described in the last section were integrated.
The results of these integrations is shown in \fig{bg:antiprotons:angular}.

%When these distributions are used to generate particles they are essentially normalised so their integral is unity.
%For a given value of $\theta$ the momentum distribution $dP/dp$ is chosen to match the first spectrum measured at a larger value of $\theta$.
%This keeps the model conservative since the true momentum spectrum should soften as one moves to higher angles which this will not do.
%$d\Phi/d\theta$ is then used to vary the overall normalisation as a function of theta.
Boyarinov \etal state~\cite{Boyarinov:1994tp} that the angular dependence of antiproton production should take the form:
\begin{equation}
%\frac{d\Phi(\theta)}{d\theta}
\Phi(\theta)=\alpha e^{\beta\cos\theta}
\end{equation}
where $\alpha$ and $\beta$ are constants. 

It is for this reason that the axes of \fig{bg:antiprotons:angular} are $\cos\theta$ and $\ln(R)$: if such a function were valid, the points would form a straight line.
In fact it seems the variation in these parameters are better described by a quadratic function, although the fitted quadratic functions all reach a minimum for $\cos\theta\ne-1$, which seems unphysical.

However, to keep things conservative, rather than use  a linear or quadratic fit to the integrated data, we use the linear interpolation between each adjacent datum to provide the function, $\Phi(\theta,\phi)$.
For the same reason, we use the integrated cross sections obtained for the fits that use a linear function for the high-momentum tail.

\TabAntiprotonRegions
\subsection{Simulating Antiproton Backgrounds}
\FigAntiprotonSimFluxes
To study backgrounds from antiprotons, four separate simulations were performed, one for each of the angular regions described in \tab{bg:antiprotons:regions}, with \num{8e7} antiprotons fed in to each simulation.
Within each angular region, antiprotons were generated uniformly across the production target, with momenta distributed according to $P(p,\theta_\textrm{min})$.
Given how steeply the angular distribution tails off, in order to check the acceptance of antiprotons produced at large angles, events were generated isotropically within the angular region, and then re-weighted during analysis based on $\xi\Phi(\theta,\phi)$, ie. the probability that an antiproton is produced in this direction.

The beamline's acceptance to antiprotons and the secondary pions that they produce are shown in \fig{bg:antiprotons:sim:fluxes}.
The events in this plot are unweighted by the antiproton angle.
From these plots, it is clear that the material at the entrance to the Torus1 solenoid removes a large fraction of the antiprotons---about 99\% of the antiprotons produced between 97 and 119\degree.
It can also be seen that the acceptance in the bent solenoids for antiprotons is about two orders of magnitude greater for antiprotons produced at angles greater than 119\degree, compared to those produced between 97 and 119\degree.

For pion production and acceptance, however, the unweighted rate for the four angular regions is very similar, particularly towards the end of Torus2.
This, and the fact that the transmission probability has a negative gradient for each angular region along the whole beamline, implies that the main source for pions is the interaction of antiprotons in the primary target.
Pions produced by antiprotons hitting collimators or the beam-pipe will quickly be removed since they will tend to stay close to this material.
If pions were produced at the foil at the entrance to the Torus1 then they would be in the centre of the beam-pipe and have a greater chance of being transported, although this does not seem to be a major source.

\FigAntiprotonSimHeightsTwoDPbar
\FigAntiprotonSimHeightsTwoDPiMin
\TabAntiprotonResultsAntip
\TabAntiprotonResultsPiSecond
\afterpage{\clearpage}

\Fig{bg:antiprotons:sim:2D-antip} shows the trajectories of antiprotons along the beamline, where each trajectory has been weighted according to the antiproton's production direction with respect to the proton beam axis.
Similarly, \fig{bg:antiprotons:sim:2D-pi} shows the same information but for pions instead.
The colour scale and range is the same for each of the four plots in these figures.
One can see, then, that although the acceptance per antiproton is considerably larger for larger angle production (119--180\degree), the production rate is so small that the final probability of observing an antiproton is quite low.
As such, it is the antiprotons produced in the forwards direction and scattered backwards by the production target that will dominate the antiproton background rate.

\FigAntiprotonSimTime
\Tab{bg:antiprotons:sim:antip} and \tab{bg:antiprotons:sim:pi} show the number of antiprotons and pions that enter the Torus1 (`Torus1'), reached the final collimator (`Torus2'), and arrived at the target (`Target').
Also shown is the sum over all events with each event weighted by the initial antiproton direction and proton energy scale correction ($\xi\cdot\Phi(\theta,\phi)$), and normalised to the number of antiprotons introduced into the simulation.
No antiprotons are seen stopping in the target for the angular region from 0 to 59\degree.
To estimate the stopping rate of antiprotons produced in this direction, 
the weighted sum for antiprotons just before the collimator was multiplied by the median acceptance of the collimator, $P(\textrm{Target}|\textrm{Torus2})$, which occurs for the region from 59 to 97\degree.

The timing distributions for antiprotons and pions are shown in \fig{bg:antiprotons:sim:time}, from where it is clear that timing information from the pulsed proton beam is not meaningful for antiprotons.
However, since the majority of secondary pions that are successfully transported to the target originate in the production target, the pulsed beam timing is a strong way to suppress the backgrounds these can produce via \ac{RPC}.
In fact, given that these pions originate in the production target, the background rate due to RPC of pions directly produced by the proton beam can be reused, including the timing cut.
For antiprotons, the rate of RPC backgrounds per pion stop is also used, however the mean number of negative pions produced per antiproton stop is included, as well as the overall efficiency of the timing window, assuming a flat timing distribution.
%All of these factors are summarised in \tab{bg:antiprotons:factors}.
For both pions and antiprotons, the efficiency of the TDAQ and reconstruction are also included, with the same values as for signal electrons.
Based on these factors, the numbers in \tab{bg:antiprotons:sim:antip} and \tab{bg:antiprotons:sim:pi} can be converted to background estimates, which are shown in \tab{bg:antiprotons:estimates}.
\TabAntiprotonEstimates

The total predicted antiproton rate per \ac{POT} is therefore \VarBgAntiprotonsPerPOT, which is completely dominated by antiprotons produced in the forward direction but scatter within the production target and are successfully transported to the stopping target.
This number is considerably larger than would be desirable for the COMET experiment.
It is important, therefore, to remember that this is an especially conservative estimate:
\begin{itemize}
\item The kinematic end-point is a large overestimate for the true end-point of the spectrum since the incoming proton will never interact with the entire nucleus.  
\item The straight-line-fit to the tail, that was used to generate antiprotons in the simulation will produce higher momentum antiprotons than are likely to occur, which changes the penetration depth in material.
      In addition, the straight-line-fit increases the integrated yield of antiprotons significantly, which impacts on the event-by-event weight.
      For example, for the region from  59--97\degree, there are up to two orders of magnitude difference on the predicted antiproton yield depending on the high-energy description.
%      \CHECKmargin{Add to the first table the integrals with the other tail descriptions.}
\item Antiproton events were generated with a uniform distribution in the production target. Given that the proton beam profile will be close to a double Gaussian, with most of the probability concentrated in the centre of the target, there will typically be more material for the antiprotons to pass through as they leave the target.  This will increase the rate of absorption and thus reduce the antiproton rate further.
%\item The data this model is based on uses 10~GeV initial protons. The antiproton yield is expected to grow rapidly at these sorts of energies.
%      According to the model in~\cite{Meco024}, the cross section at 8~GeV is at least an order of magnitude lower than for incident protons with 10~GeV.
\item The final value assumes that all the negative pions produced from stopping antiprotons are captured by nuclei in the target.  In reality, many of these will escape from the target and decay in flight or be stopped and captured elsewhere.  The pion multiplicity factor used here is likely an overestimate, therefore.
\end{itemize}
At the same time, the fact that these results are somewhat statistically limited, given the small number of antiproton stops seen from the forward going antiprotons, one can expect a sizeable uncertainty on this value.
In addition, some of the other factors that have been included, such as from the RPC background estimation, also increase the uncertainty since there might be other correlations that such a factorisation cannot include.

On the other hand, the fact that it is scattering from forward produced antiprotons has some advantage, since there is experimental data in the forward region---admittedly with materials other than tungsten or with proton energies other than 8~GeV---that could be used to develop further understanding of this background.

\subsection{Reducing Antiproton Backgrounds Further}
%\begin{easylist}
%# Ways to improve this / further optimisations:
%## Antiprotons close to the beam axis after the production target: on-axis blocker (activation?)
%## Absorber material in lower half of TS3
%# Thicker absorber at the muon beamline entrance
%# Need to check the muon beam intensity
%\end{easylist}
If the background rate due to antiprotons presented here proves correct, then changes to the current beamline design should be implemented to reduce it.
%Given that the antiproton with the current beamline design appears larger than desirable, it would be necessary to take steps to reduce it.
There are several options that can be pursued, but each must be checked for its impact on the muon beam intensity.

Firstly, one can simply increase the thickness of the foil at the entrance to the Torus1.
In a similar way, one can introduce additional foils along the beamline, such as at the TS3 solenoid in between the Torus1 and Torus2.
The downside of both of these methods is that they will be in the path of the muon beam, and since the $dE/dx$ curve for negative muons below 100~MeV has a negative gradient, this material will likely increase the mean muon momentum (but, of course, decrease the overall intensity).

If material is added, then it would be beneficial if this material were low-Z.  This would maximize the ratio of the rate of antiproton loss to the muon loss rate.
In fact, a bug in the geometry during initial simulations of antiprotons left air at the exit to the production section.  
The impact on the muon and pion momentum and acceptance was found to be negligible, but the antiproton acceptance into the Torus1 was zero at the limit of the sample size of the simulations.

It may be possible to identify a region in the beamline close to the production target where a separation occurs between antiprotons and low-energy muons and pions.
This would be highly favourable as it would reduce the impact on signal sensitivity, although in the studies performed here such a region has not been apparent.
Additionally, material close to the production target would need to be considered from an activation and safety perspective.

Finally, as well as adding material along the beamline, one could consider a slight reduction in the proton beam energy.
Even a reduction by 200 MeV would reduce the total antiproton production cross-section by between 40 and 50\% based on the methods described above to determine the $\Xi$ weighting factor.
Reducing the beam power further, by around 500 MeV would reduce the cross section by 75 to 90\% of the current estimate, \ie reduce the production rate by about one order of magnitude lower.

\section{Direct Beam-Related Backgrounds}
\sectlabel{bg:beam}
%\begin{easylist}
%	# Number of POT simulated
%	# Number of observed high-momentum electrons or muons
%	# Lifetime of muons at the target (need to smear this a bit ?): Fit the projection of high-momentum particles to the time with an exponential and extrapolate
%	# Survival probability per high-momentum particle before the target to reach the straw tracker
%\end{easylist}
Although neither a muon or pion in their rest frame can decay to electrons greater than 55~MeV, if these particles are boosted sufficiently the resulting electron can be produced with signal-like momentum.
For muons at 78~MeV/c, Michel electrons can occur at 105~MeV/c if produced in the direction of the muon's momentum (and neglecting muon polarisation).
For pions, the threshold for signal-like electrons to be produced during decay is 58~MeV/c, although the branching fraction for this process is helicity suppressed to around \num{1e-4}.

The bent solenoids and the pulsed beam help to suppress these sorts of backgrounds, whilst the beam blocker, downstream of the target disks, reduces the geometric acceptance.
Following the results of chapter~\sect{phaseII-optimisation}, the inner radius of the collimator at the exit of Torus2 is set to 12~cm, whilst the beam blocker has a radius of 25~cm. 
As the beam passes along the stopping target region, the beam envelope grows due to the reduction in the solenoidal field strength.
The radius of the beam blocker, however, is large enough that even with this growth there should be no direct line of sight between the muon beam and the spectrometer.
As such, the only way for beam particles to reach the detector are by a hard scatter off the target or beam blocker, or to be produced via decay of another particle close to the target itself.

To check the probability that signal-like electrons are produced in the beam and arrive at the detector in the time window,
the output of the large production target simulation was resampled five times, so that around \num{1.1e9}~\ac{POT} events were studied, equivalent to about 8~\phaseII bunches.
\Fig{bg:beam:MomVsTime:strawTrk} shows the momentum and timing of electrons at the straw tracker, from which it is clear that no background events were observed within the gated-time detector window.
One electron is observed at 98~MeV/c, but this occurs around 100~ns after the proton bunch.
This immediately sets a limit on the background rate---less than 1 in \num{1.4e9} per POT.
\FigBgBeamMomVsTime

To be able to understand this background better, it is necessary to apply some extrapolation.
Indeed, it will never be feasible to directly simulate the $10^{21}$ or so POT that would be necessary to deduce the background rate without some extrapolation.
To make this extrapolation, we look at the same distribution immediately after the beam blocker.  
At this point there are far more electrons so it is easier to fit and model the distributions.
\Fig{bg:beam:MomVsTime:stopTgt} shows the momentum vs.\ time plot for electrons immediately after the beam blocker.

The rate of observing delayed, signal-like electrons due to high-energy beam particles is modelled by factorisation of the acceptances:
\begin{align}
R_\mathrm{beam}=&A_\textrm{geom}A_\textrm{mom}A_\textrm{time},\\
=&\left(\frac{N_\textrm{obs}}{N_\textrm{POT}}\right)\left(\frac{\int\limits_{104.2}^{105.5}f(p)dp}{\int\limits_{p_\textrm{low}}^{105.5}f(p)dp}\right)\left(\frac{\int\limits_{600}^{1200}g(t)dt}{\int\limits_{30}^{1200}g(t)dt}\right),
\eqlabel{bg:beam:factors}
\end{align}
where $A_\textrm{geom}$, $A_\textrm{mom}$, and $A_\textrm{time}$ are geometric acceptance and efficiencies of the momentum and timing thresholds, respectively.
$N_\textrm{obs}$ is the number of events actually seen at the straw tracker, whilst $N_\textrm{POT}$ is the number of simulated \acp{POT}.
In the case of prompt beam electrons, the above equations are the same, except that $A_\textrm{time}$ becomes the extinction factor multiplied by the fraction of time spent in the gated-time detection window.
\FigBgBeamExtrapolate

The functions $f(p)$ and $g(t)$ represent the momentum and time projections of electrons after the beam blocker. 
Their factorisation in this manner is valid due to the lack of strong correlation visible in \fig{bg:beam:MomVsTime:stopTgt}.
It is reasonable to assume that, close to the signal region, the momentum distribution observed immediately after the beam blocker will have the same shape as the momentum distribution at the straw target.
Since the acceptance function is essentially flat at the high-energy limit (see \sect{optim:stopTgtPos} and \sect{sense:momentum}), only the normalisation of the spectrum should change, and this we have encapsulated already in the geometric acceptance.

Finally, the timing distribution, $g(t)$, is also taken by fitting the observed electron distribution after the beam blocker.
Here, however, one needs to be careful about extrapolating the distribution from after the beam blocker to the straw tracker.
The extrapolation must account for two effects: an overall translation of the distribution to later times, and a broadening of the distribution (an increase in the fitted lifetime).
Both effects are determined by the time needed for particles at the beam blocker to propagate through the spectrometer.
With a radius of 2~m, the typical distance travelled along the Spectrometer will be at least about 6~m, and since the electrons are relativistic and travelling at close to the speed of light, transport through the spectrometer can be expected to add about 20~ns.
This assumes that the electron's pitch angle is equal to zero so that they follow the beam axis perfectly.
With a large pitch angle the distance travelled, and therefore transportation time, will be longer, although since larger pitch angles drift more in the bent solenoids, these particles will tend to be removed more effectively.
The effect of broadening on the distribution arises from lower energy electrons moving more slowly, however for the relativistic signal-energy electrons this effect will be negligible.

\Fig{bg:beam:acceptance} show the momentum and time functions based on the distribution of electrons after the stopping target.
The equations for the lines of best fit are also shown, and the region used to fit indicated by the solid black fitted line.
The low-momentum threshold for this study, $p_\textrm{low}$, is set to 65~MeV/c, since this removes electrons from \ac{DIO} of muons along the beamline, which have relatively long life-times but cannot be produced at signal energies.
Above this threshold, there are 65 events at the straw tracker, giving a geometric acceptance of \VarBeamBgGeometric.
For the efficiency of the momentum cut, using the fit to the projection shown in \fig{bg:beam:acceptance:momentum}, the ratio of events expected between 104.2 and 105.5~MeV/c compared to the integrated rate from 65 to 105.5~MeV/c is \VarBeamBgMomentum.
For the timing efficiency, the fit given in \fig{bg:beam:acceptance:time} is first shifted so that the mean of timing distribution at the beam blocker sits at the time of the latest arriving electron above 65~MeV/c, as indicated
by the vertical magenta lines. 
This amounts to a shift of 110.5~ns, which is about twice as large as if one were to use the mean of the straw tracker electrons.
Having applied this shift (which we implement as a shift in the start and stop time for the integration), we find the suppression of electrons due to the time-gated window gives a factor \VarBeamBgTiming.

\TabBgBeamFactors
All of these factors are summarised in \tab{bg:beam:factors}.
Finally, then, we can pull all these factors together using \eq{bg:beam:factors}, from which we find that the predicted prompt beam-related background rate per POT is~\VarBeamBgPromptPerPOT, whilst the delayed background rate due to beam particles is \VarBeamBgDelayedPerPOT.

Given the extrapolation that has been necessary, this result comes with a large degree of uncertainty.
However, in principle this includes all sources of electrons in the beam: high-energy muons, pions, and electrons coming directly from the production target with high energy.
This also includes electrons produced from neutron capture, provided the neutron itself originated downstream of the production target section, which were not included here in order to accelerate the simulation speed.
In addition these neutrons should be studied carefully using multiple hadron production models but, as discussed in \sect{bg:notCovered}, these neutrons are left to future studies.

\section{Cosmic Ray Background}
\sectlabel{bg:cosmics}
\FigBgCosmicPrimary
Cosmic rays at sea level consist primarily of positive and negative muons.
These are produced when high-energy particles originating in the sun or beyond the solar system strike the upper atmosphere and produce a particle shower.
About one muon per square centimetre per minute is expected at sea level, with the average muon energy around 4~GeV~\cite{PDG2014}.
At these energies, not only can such muons readily produce high-momentum electrons, but they are weakly ionising and therefore penetrate material relatively easily.

This makes cosmic muons a potential source of backgrounds, as the high-energy muons can pass through the experiment hall to the beamline whereupon they can produce signal-like electrons.
In addition, this background would lack any timing structure.
As a result, a  \ac{CRV} is included in both \phaseI and \phaseII, made with four layers of scintillator bars as the active material.
When a cosmic ray passes this, events in the detector can either be tagged, or any trigger on this event vetoed and not read out, such that signal-like electrons produced from cosmic rays should not be included in the analysis.

At the time of writing, the \ac{CRV} has been optimised and studied for \phaseI~\cite{TDR2016,CM18:CRV,CM20:CRV}.
These studies suggest that a miss-rate of fewer than 1 in $10^4$ events is achievable.
For \phaseII however, given the changes to the experiment geometry, the situation is a little different.
%If the muon penetrates into the detector and decays from there, or if a high-energy electron produced outside of the detector but close to it, is 
Principally, the stopping target in \phaseI is surrounded by the \ac{CDC}, such that cosmic muons must escape detection in both the \ac{CRV} and \ac{CDC}, and then decay close to the target.
In \phaseII, the stopping target is located well away from the detector, which means cosmic muons can only be vetoed by the \ac{CRV}.
Furthermore, the \ac{CRV} geometry in \phaseII should cover both the target and the intermediate beamline, in addition to the detector solenoid itself.

\phaseII is expected to run about \VarRunTime~s as suggested in section~\sect{sense}.
During COMET operation, a single accelerator cycle is expected to last 2.48~s, with about 0.7~s for extraction, giving a duty factor of 0.28.
In addition, with a time-gated window lasting 600~ns and bunches separated by 1.170~$\mu$s, the overall duty factor drops to about 0.14.
As such, during \phaseII running there are really only about \num{2.20e6}~s when a cosmic muon can produce backgrounds.
During this time then, \num{4.8e4} cosmic muons should pass through each horizontal square centimetre of the COMET hall.
If we assume the miss-rate for the \phaseII will be no worse than for \phaseI, \ie \num{1e-4}, then for each for each square centimetre of the experiment, more than 48 muons should be studied.

\subsection{Simulating Cosmic Muons}
To simulate the cosmic ray background, input muon events were provided by the ND280 experiment~\cite{George:2010aca}.
These contained positive and negative muons that had been produced in an air shower simulation using the CORSIKA simulation package~\cite{Corsika}, for the atmosphere above J-PARC.
\Fig{bg:cosmics:primary} shows the distributions of momentum and directions that these files contained.

Muons were generated in a horizontal plane directly above the experiment hall (and centred on the origin of the ICEDUST coordinate system, the production target).
A 50$\times$50~m${}^2$ square within in this plane was used to generate the events.
This immediately adds a limitation to this simulation that very horizontally-moving muons cannot be studied, but it is needed to increase the simulation efficiency.
%Given the angular distribution shown in \fig{bg:cosmics:primary} and the offset of the detector solenoid with respect to the production target, between XX and XX of all cosmic events can be simulated this way.
About \num{1.76e8} cosmic events were generated in this way; for the 2500~m${}^2$ region above the experiment where the simulation was started, more than \num{1.2e8} events should be studied, according to the above calculations.

Material above the beamline can degrade and stop cosmic muons so it is important this is modelled well.
The concrete and iron shielding for the entire experimental hall is implemented, as well as the shielding necessary for the \ac{CRV}.
However, at this stage the geometry in the simulation only has the CRV design intended for \phaseI around the detector solenoid (in its \phaseII position).
Since the CRV will need to be extended to include all the way up to the stopping target in \phaseII, there will likely be more shielding material inserted above the beamline.

\subsection{Results of the Simulated Cosmic Background Rate}
\FigBgCosmicBeam
\FigBgCosmicMomenta
The tracks of high-energy electrons that enter the detector region are shown in \fig{bg:cosmics:beam}.
\Fig{bg:cosmics:momenta} shows the momentum distribution of observed electrons.  
There are large fluctuations across the momentum distribution, but this is most likely due to statistical fluctuations rather than truly physical.
From 80 to 120 MeV/c, 140 electrons hit at least one layer in the straw tracker, although only 19 of these hit a second plane.
The number of background events due to cosmic rays can therefore be estimated as:
\begin{align}
	R_\textrm{cosmics}=R_\mu\frac{n_e^\textrm{obs}}{N_\textrm{sim}} \Delta{}p (1-\epsilon_\textrm{CRV}) A_\textrm{duty}A_\textrm{timing}A_\textrm{TDAQ,Recon},
\end{align}
where $R_\mu$ is the rate of cosmic muons expected through the experiment hall per second, $n_e^\textrm{obs}$ is the observed number of events per MeV/c, $N_\textrm{sim}$ is the total number of simulated cosmic ray muons, and $\Delta{}p$ is the width of momentum thresholds, from \VarMomThreshold to 105.5~MeV/c.
$\epsilon_\textrm{CRV}$ is the \ac{CRV} efficiency, such that $(1-\epsilon_\textrm{CRV})$ is the expected miss-rate. 
The factors $A_\textrm{duty}$,$A_\textrm{timing}$, and $A_\textrm{TDAQ,Recon}$ represent the accelerator duty factor, fraction of time spent during the time-gated detection window, and the combined TDAQ and reconstruction efficiency, respectively.  
The values for each of these parameters are given in \tab{bg:cosmics:params}.
From this we can estimate that there will be \VarCosmicRatePerSecond backgrounds per second due to cosmic rays.
\TabBgCosmicParameters

To reduce the cosmic ray background further, the most important parameter is the \ac{CRV} efficiency, the complement of which multiplies the entire background rate.
Since this background is completely uncorrelated with the beam, the shorter the running time, the lower this background rate will be.
If this cannot be reduced further, then it might be necessary to add more material above the beamline or above the experiment hall.
Fortunately, rather than depend on simulation for this, \phaseII has the benefit of \phaseI, which will have measured and characterised the cosmic ray flux on the experiment hall floor.

%%Since all previous studies have only required a hit in one or more straw tracker planes for a track to be `accepted', in the interest of consistency we shall do the same here.
%The 19 events between 80 and 120 MeV/c coming from 176 million cosmic ray simulated events imply \VarCosmicRatePerMuPerMeV signal-like electrons per primary muon per MeV/c, or \VarCosmicRatePerMu within the momentum thresholds.
%As such, for the rate of cosmic rays passing through the 50$\times$50 m${}^2$ patch of sky above the hall, we can expect to be 
%With the signal momentum range of \VarMomThreshold to 105.5~MeV/c, and assuming 1 muon per minute per cm${}^2$, 
%
%Of these events, about XXX electrons with signal-like momenta were observed to hit a straw tracker plane, which are shown in the event display in \fig{bg:cosmic:eventDisplay}.
%No events were observed hitting more than one of the straw planes with the necessary momentum to be considered a background, and as such no events would have been reconstructed.
%Accordingly, we can only set a limit on the rate of backgrounds due to cosmic events coming from overhead muons.
%At 90\% C.L.\ the rate must be smaller than XXX to expect to see no background events in \num{1.7e8}

\section{Summary of Background Rates}
\sectlabel{bg:summary}
\TabBackgroundFinalVals
\Tab{bg:summary} gives an overview of the predicted rates.
The overall expected number of background events during \phaseII is \VarTotalBgPhasII events, assuming \VarRunTime~s of run time.
This is dominated by the cosmic ray background and events from antiprotons that stop in the target.
It is clear, though, that further work must be undertaken on these predictions in the future.
In light of the various uncertainties, discussed in more depth below, we have tried to be as conservative as reasonably possible in this estimate; improving the uncertainties would hopefully reduce the estimated number of events.

\clearpage
\subsection{Backgrounds Not Considered}
\sectlabel{bg:notCovered}
\subsubsection*{Neutrons from the Production Target}
Although some of the methods above are different, compared to the \phaseI TDR and \phaseII CDR there is only one source of backgrounds that has not been covered here:
the production of high-energy electrons from the capture of neutrons that originate in the production target.
These neutrons were killed in order to reach a large number of muons and pions in the simulation, although neutrons produced along the muon beam (from the collimators, stopping target, \etc) have been included with the beam background.
Whilst there is a large amount of concrete and iron material in between the production target and the detector, the production target is likely a major source of neutrons in the experiment, and this needs to be studied rigorously using multiple hadron production and neutron transport models.

Such studies have advanced for \phaseI, with the background rate from neutrons estimated to be about \num{3.1e-29}~per \ac{POT}~\cite{TDR2016}.
With a different experiment layout, production target material, and higher proton beam power, separate studies should be performed  for \phaseII.
A na\"{i}ve scaling of this to the \phaseII number of protons on target would suggest that \num{2.7e-8} events can be expected in \phaseII.
Since the detector sits further away from the production target and since there is additional beamline then it is likely that the neutron rate is lower.
%The increased proton power would already be included in the naive scaling applied above to first order.
The change of target material between \phaseI and \phaseII from low to high-Z material will increase the neutron production rate per unit mass of target.
However, to compensate for the equivalent reduction in muon yield that this produces, the target in \phaseI will be longer and as such the neutron rate per POT might well be similar.
%The different target material will change the neutron production cross section, but since the target is also made longer in \phaseI to maximize the muon stopping rate, the increase in neutron yield due to the different target geometry and material might not be significant.

\subsubsection*{Detector Performance}
It has been implicitly assumed in the preceding work that the StrECAL is able to perfectly distinguish electrons from positrons, muons, and pions.
This seems a reasonable assumption given the charge and relativicities of the particles, but it would be important to consider the rate of mis-identification in the event that these particles are able to be transported to the detector.
These particles are already suppressed by the long beamline, however, and only muons have been seen to be transported to the detector in the simulations for this thesis.

In a similar vein, the impact of `pile-up events' should be considered.
Pile-up occurs when two separate particles enter or strike the same region of a detector in a shorter space of time than the detector can resolve.
In the ECAL, pile-up could be an issue for triggering but since the ECAL's primary purpose in the analysis is to provide a time stamp, this should not be an issue for backgrounds.

Pile-up in the Straw Tracker should be easier to resolve because of the finer granularity of the detector, and the use of multiple hits on multiple straws to reconstruct the trajectory.
Instead, it is the occupancy of the Straw Tracker that should be considered.  
How often is a 105~MeV/c electron found due to the mistaken combination of hits in a high-occupancy event?
Machine learning techniques for \phaseI's \ac{CyDet} have been developed to significantly suppress hits from tracks other than signal electrons~\cite{TDR2016} and it is likely that a similar approach will be possible in \phaseII.
In addition, although the beam power and production target material in \phaseII increase the multiplicity of particles in the beam, the longer beamline may well bring the occupancy down to lower levels than observed in \phaseI.

%Other sources of background might exist---putting a limit on all processes that produce fewer than a single electron from $10^{21}$~POT is no easy feat---and it 

\subsection{Impact of Momentum Resolution}
\FigBgVsResolution
The momentum resolution is the key factor in determining the background rates for the intrinsic backgrounds (\acf{RMC} and \acf{DIO}).
However, since the distribution of residuals (the resolution function) is not well known for \phaseII, there is a degree of uncertainty over the size of these backgrounds.

To understand how robust these estimates are to changes to the resolution function, a simple parametrisation for the resolution function was used.
In this simplified scheme, the probability density for mis-reconstructing an electron's momentum by a certain amount, is given by the function:
\begin{equation}
	P(p_\textrm{reco}-p_\textrm{true}) = P(\Delta{}p) \propto e^{(\Delta{}p/\sigma_1)^2}+\kappa e^{(\Delta{}p/\sigma_2)^2}
\end{equation}
where $\kappa$ is dimensionless and used to set the relative weight of the second Gaussian, and $\sigma_1$ and $\sigma_2$ define the widths of the respective Gaussian.

\Fig{bg:resolution} shows the background rates due to \ac{DIO} and \ac{RMC} as a function of $\sigma_2$ for different values of $\sigma_1$ and $\kappa$.
The current \phaseI results suggest $\kappa$ to be about 0.08, and $\sigma_1$ to be less than 200~keV/c, and so these values are used as baseline estimates.
To produce the plot in \fig{bg:resolution}, we fix the momentum thresholds to be \VarMomThreshold and \VarMomThresholdHigh, as found from the sensitivity estimates.
Over the range of $\sigma_2$ values shown in \fig{bg:resolution}, the signal acceptance, given these momentum thresholds, changes by less than 4.2\%.
  
From this it is clear that provided the value of $\sigma_2$ is kept below 1~MeV/c, fewer than 0.01 background events would be expected from \ac{RMC}, whilst the background rate from RMC grows exponentially above this point.
However, it is also clear that the rate from DIO electrons will dominate regardless of the resolution.
This should not be surprising given that the \ac{DIO} spectrum reaches up to the \mueconv end-point whilst the RMC end-point is separated by some 2.4~MeV/c.

\subsection{Principal Uncertainties}
Whilst it is not possible to make a quantitative estimate of the uncertainty at this time, it is clear that for some backgrounds the uncertainty is large.

\begin{description}
\item[Statistical Uncertainty and Use of Extrapolation]
In many instances the generated sample sizes for these background predictions were too small to directly evaluate the background rate.
In these cases it was necessary to extrapolate from parameters and regions with higher statistical power, such as lower energy parts of a spectrum, earlier in time, or use a more upstream part of the experiment.
In addition, it was often necessary to resample particles produced from an earlier stage of the simulation, such as for the muon and pion stopping distributions, or the beam fluxes leaving the production target region.
This introduces a different kind of error, which is very difficult to quantify, namely that there will now be some correlation between the events seen downstream.
In the future, generating larger Monte Carlo samples using the continually increasing computing power available to the COMET experiment will be important in reducing this uncertainty.

\item[StrECAL Reconstruction and Resolution]
In the case of \ac{DIO} and \ac{RMC}, the two intrinsic sources of background, the main source of uncertainty is in the resolution function.
The exact distribution for the residuals after reconstruction is not well known for the \ac{StrECAL}, especially for \phaseII.
In particular, for \ac{RMC}, the high-energy tail needs to be studied better. 
It has been assumed that fewer than \VarResolutionProbFromRMCToSignal events will be reconstructed with momentum larger by at least 2.4~MeV/c than the true value.
This assumption must be confirmed, however.
RMC on Al-26 was also discussed as a potential concern, however Al-26 has a tiny abundance in natural aluminium, whilst the rate of production in the COMET muon beam is difficult to predict.

\item[RPC Spectrum]
In the case of \acf{RPC}, the dominant uncertainty comes from the spectrum of photons produced during RPC.
The study here and in the \phaseI TDR depend entirely on a single set of experiments, which do not cover RPC from aluminium, only on magnesium and calcium.
This will be difficult to improve further without additional observations, and it would be useful to study this at \phaseI if possible, perhaps by changing the Torus1's dipole field to improve pion acceptance.

\item[Antiproton Production]
Given the lack of experimental evidence for antiproton production in the `sub-threshold' energy regime of the backwards direction, there is large uncertainty over the expected rate.
Whilst the model developed here was kept conservative, as one of the dominant backgrounds it is important to study this further.
A dedicated study comparing multiple different hadron production codes to all the possible experimental datasets (including forward and backward produced antiprotons from the different targets with different proton momentum) could reduce our uncertainty if no new experimental data is forthcoming.
Such a study would then be able to extrapolate better for the target material, proton energy, and production direction used in COMET.

In addition to the production rates, there is reasonable uncertainty on the mechanisms by which antiprotons can produce background electrons.
Here it has been assumed that the primary mechanism is via pion production, but the rates for such production in aluminium is also not well known.
Such an assumption must also be tested.
%Direct emission of high-energy electrons is supposed to be relatively small~\cite{}, but it has not been rig

\phaseI will be able to play a sizeable role in reducing this uncertainty as one of the first experiments to look at antiproton production under the \COMET conditions.
The fact that the production target will likely be made from graphite in \phaseI as opposed to the tungsten of \phaseII reduces the direct usefulness of such measurements however.
On the other hand, given that there is more (although still rather limited) data for graphite will allow \phaseI measurements to better understand the transportation of antiprotons along the COMET beamline.
Furthermore, in addition to antiproton measurements during the main \phaseI runs, additional runs are under consideration to specifically study the antiproton background.
These could involve changing the production target material is to something other than graphite or changing the beam energy and studying the change in antiproton yields.

\item[Cosmic Ray Veto Geometry and Shielding]
Without a complete \ac{CRV} design, it is tough to be confident that the energy loss of cosmic muons is well captured in the \phaseII simulation.
In addition, the miss-rate of \num{1e-4} should be experimentally confirmed and ideally improved.

Again, after \phaseI the uncertainty on this background should be significantly reduced.
The muon flux in the experiment hall, at the beamline position will have been measured to a high precision in all directions.
Furthermore, since the CRV design will be largely unchanged the miss-rate will be well understood.
\end{description}

\section{Further Studies and Improvements}
\sectlabel{bg:improvements}
This chapter has presented an almost comprehensive set of estimates for background sources in \phaseII.
It should not, however, be treated as the final set of predictions, but rather as an initial set of studies that should be developed further.
In addition to the previous suggestions regarding uncertainties in the present studies,
the following list gives a summary of possible methods that could be pursued to improve and expand the preceding set of estimates.

\begin{description}
	\item[Explore techniques for improving the statistical power of limited Monte Carlo data]
		In addition to the use of resampling and extrapolation techniques used here, cross checks with various other methods should be explored.
		Together these would provide an increase in the statistical power of a limited set of simulated data, and provide a validation and uncertainty estimate for one another.
		Methods such as particle weighting and biasing could also be employed.
		Additionally, the transmission of each section of the beamline could be studied separately, building a set of multi-dimensional (particle type, time, momentum, position, \etc) transfer matrices.
		By taking the product of these the overall transmission of the beamline can be deduced as a function of input particle type.
	\item[Use different hadron production models] Primarily this is needed to cross check the muon and pion yields, and the production target optimisation.
		This is already well accommodated in ICEDUST, but should be exercised properly.
	\item[High-momentum electron acceptance vs. beamline position] Many background estimates would benefit from knowing the exact acceptance to high-momentum electrons at a given location in the beamline.
		This could be in the form of a multi-dimensional histogram or matrix, and would ideally also contain the acceptance's dependence on the electron direction, and transverse location.
		With this information, simulations for the cosmic ray backgrounds, photon conversion after RPC and RMC, and beam simulations would only need to know the probability of creating a high-momentum electron at a given location to deduce the resulting background rate.
		This would also benefit the optimisation studies, as mentioned in chapter~\sect{phaseII-optimisation}.
\end{description}

The timescale for these studies is fairly varied: whilst \phaseII is several years away still, in some cases the true deadline will be \phaseI, which is now little more than a year away.
Decisions for the design of \phaseI will impact, in some cases, the design of \phaseII, whilst the run plan for \phaseI will include dedicated studies aimed at reducing key uncertainties of \phaseII.
For both of these reasons it is important that there is a continuous effort to maintain and refine the understanding of COMET \phaseII.
%\begin{easylist}
%# Neutron background with different hadron codes and \alcap neutron spectrum from mu capture
%# Improved cosmic ray veto geometry
%# High-energy electron acceptance
%# Pion stops elsewhere in the beamline
%# StrECAL resolution function
%# Improved proton beam timing structure 
%# RMC from Al-26 that is produced from Al27 via neutrons. Photons can be made up to 108 MeV
%# Assumption that the PID abilities of the StrECAL are 100\%, is this true? At what level would muons be reconstructed at electrons?
%# RPC from pions stopping elsewhere in the beam -> Map of acceptance to RPC electrons vs. position of pion stop in the beamline
%\end{easylist}


  %% To ignore a specific chapter while working on another,
  %% making the build faster, comment it out like this:
  %\input{chap4}
\end{mainmatter}

% Produce the appendices
\begin{appendices}
\chapter{Drifts in a Bent Solenoid}
The Lorentz force:
\begin{equation}
\frac{d\vec{p}}{dt}=\frac{q}{m}\vec{p}\times\vec{B}
\end{equation}

\section{Uniform Solenoidal Field}
B field is uniform and parallel to axis of solenoid.
Define the Larmor frequency, $\omega$, and radius, $a$, as:
\begin{align}
a=&\frac{\gamma m \vec{v}\times\vec{B}}{qB^2} = \frac{p_\mathrm{T}}{qB} \\
\omega=&\frac{qB}{m}
\end{align}

\section{Field in a Bent Solenoid}
Producing a cylindrical solenoid  channel can be imagined as directly bending that of a normal uniform and linear one.
By symmetry it can be seen that any gradient introduced to the magnetic field can only be radially, in the plane of the bending.
Further, by considering Ampere's law with a current loop in the plane of bending formed by two radial straight lines (with length $|r-R|<L$, where $R$ and $L$ are the bending and aperture radii of the solenoid channel) and an arc,
it can be seen that the variation in the field is given by:
\begin{equation}
\frac{\vec{\nabla B}}{B}=\frac{1}{r}\hat{r}
\end{equation}
%\CHECK{is this accurate?}
\CHECK{Sketch or figure?}

\section{Drift Calculation}
There are two sources of drift in a bent solenoid: the gradient in the field, and the centrifugal force arising from the circular coordinate system needed to describe the field lines.
The two can be treated separately in the sense that the motion of a particle moving through a field with straight field lines but with a transverse gradient given by $\nabla B/ B \propto 1/r$ would be described by an equation of motion equivalent to that from the first source of drift in the bent solenoid system. 
Similarly, a system with a uniform field but field lines that follow circular paths would exhibit drift equivalent to the second component mentioned above.

\subsection{Gradient Drift}
``Grad-B'' drift is well described in text books, but in the interest of completeness a short derivation shall be given here.
The drift arising due to the gradient in the field can be treated as a perturbation of the motion of the particle in a uniform solenoidal field.
The total velocity $\vec{V}$, is given by:
\begin{equation}
\vec{V}=\vec{v}+\vec{v_g},
\eqlabel{velocity-grad}
\end{equation}
where $\vec{v}$ is the unperturbed velocity of the particle in the transverse plane, and $\vec{v_g}$ is the velocity arising due to the gradient in the field.

Treating the field as a Taylor expansion:
\begin{equation}
\vec{B}(\vec{r})=\vec{B_0}+\left.(\vec{r}\cdot\nabla)\right|_{\vec{r}=0}\vec{B}+\ldots
\eqlabel{B-field-taylor}
\end{equation}
and substituting equations \eq{velocity-grad} and \eq{B-field-taylor} into the Lorentz force, gives:
\begin{align}
m\frac{d(\vec{v}+\vec{v_g})}{dt}=&q(\vec{v}+\vec{v_g})\times\left(\vec{B_0}+\left.(\vec{r}\cdot\nabla)\right|_{\vec{r}=0}\vec{B}\right) \\
%=&q\vec{v}\times\vec{B_0} \\
%&+q\left(\vec{v}\times\left.(\vec{r}\cdot\nabla)\right|_{\vec{r}=0}\vec{B}+\vec{v_g}\times\vec{B_0}\right)
\end{align}
so that to first order, the perturbing velocity is given by:
\begin{equation}
\frac{d\vec{v_g}}{dt}=\frac{q}{m}\left(\vec{v}\times\left.(\vec{r}\cdot\nabla)\right|_{\vec{r}=0}\vec{B}+\vec{v_g}\times\vec{B_0}\right)
\end{equation}
Since we are only interested in steady-state solutions where $dot{v_{g}}$ is close to zero, the above equation gives:
\begin{equation}
\vec{v_g}=\frac{q}{m}\frac{\vec{B_0}\times\left(\vec{v}\times\left.(\vec{r}\cdot\nabla)\right|_{\vec{r}=0}\vec{B}\right)}{B^2_0}
\end{equation}
which by considering the form 	after time averaging becomes:
\begin{equation}
\langle\vec{v_g}\rangle_t=\frac{\vec{B_0}\times\left(\vec{v}\times\left.(\vec{r}\cdot\nabla)\right|_{\vec{r}=0}\vec{B}\right)}{B^2_0}
\end{equation}

\chapter{Kinematic End-point for Anti-proton Production}
\sectlabel{appendix:antiprotonEndpoint}

\newcommand{\TabTgtImprovDiskParams}{
\begin{table}[tb]
\centering
\sisetup{
table-number-alignment=center
}
\begin{tabular}{|r|S[table-format=5.1]S[table-format=1.5]S[table-format=3.3]|}
\multicolumn{1}{|L{1.1cm}}{Disk Index} & \multicolumn{1}{|L{1.8cm}}{Beamline Position (mm)} & \multicolumn{1}{L{1.5cm}}{Relative Gyroradius} & \multicolumn{1}{L{2cm}|}{Disk Radius in Simulation} \\
\hline
0 & 19542.0 & 1.09832 & 144.758 \\
1& 19592.1 & 1.11405 & 148.932 \\
2& 19642.2 & 1.13064 & 153.400 \\
3& 19692.3 & 1.14836 & 158.248 \\
4& 19742.4 & 1.16753 & 163.576 \\
5& 19792.5 & 1.18847 & 169.495 \\
6& 19842.6 & 1.21130 & 176.068 \\
7& 19892.7 & 1.23601 & 183.326 \\
8& 19942.8 & 1.26227 & 191.198 \\
9& 19992.9 & 1.28943 & 199.516 \\
10 & 20043.0 & 1.31668 & 208.037 \\
11 & 20093.1 & 1.34304 & 216.451 \\
12 & 20143.2 & 1.36771 & 224.474 \\
13 & 20193.3 & 1.39005 & 231.867 \\
14 & 20243.4 & 1.40984 & 238.518 \\
15 & 20293.5 & 1.42688 & 244.317 \\
16 & 20343.6 & 1.44127 & 249.272 \\
17 & 20393.7 & 1.45330 & 253.450 \\
\hline
\end{tabular}~%
\begin{tabular}{|r|S[table-format=5.1]S[table-format=1.5]S[table-format=3.3]|}
\multicolumn{1}{|L{1.1cm}}{Disk Index} & \multicolumn{1}{|L{1.8cm}}{Beamline Position (mm)} & \multicolumn{1}{L{1.5cm}}{Relative Gyroradius} & \multicolumn{1}{L{2cm}|}{Disk Radius in Simulation} \\
\hline
18 & 20443.8 & 1.46330 & 256.950 \\
19 & 20493.9 & 1.47151 & 259.842 \\
20 & 20544.0 & 1.47823 & 262.219 \\
21 & 20594.1 & 1.48347 & 264.080 \\
22 & 20644.2 & 1.48745 & 265.500 \\
23 & 20694.3 & 1.49056 & 266.612 \\
24 & 20744.4 & 1.49296 & 267.471 \\
25 & 20794.5 & 1.49478 & 268.125 \\
26 & 20844.6 & 1.49612 & 268.604 \\
27 & 20894.7 & 1.49703 & 268.933 \\
28 & 20944.8 & 1.49730 & 269.030 \\
29 & 20994.9 & 1.49730 & 269.030 \\
30 & 21045.0 & 1.49714 & 268.972 \\
31 & 21095.1 & 1.49692 & 268.891 \\
32 & 21145.2 & 1.49671 & 268.817 \\
33 & 21195.3 & 1.49667 & 268.804 \\
34 & 21245.4 & 1.49673 & 268.823 \\
   &         &         &         \\
\hline
\end{tabular}
\caption
[Parameters of the target disks used in the proof of principle simulation.]{
Parameters of the target disks used in the proof of principle simulation.
Since the beamline axis is parallel to the $-z$ direction in the stopping target region, the Z-coordinate can be obtained from the beamline position via: $z=2164-(s-19542)$.
\tablabel{app:tgtImprov:diskParams}}
\end{table}
}

\newcommand{\FigTgtImprovFieldVsBeamline}{
\begin{figure}[tb]
\includegraphics[width=0.9\textwidth,trim=1.0cm 0cm 0.8cm 0,clip]{figs/appendixStopTgtImprov/Plot_FieldMagnitude.pdf}
\caption
[The magnitude of the magnetic field around the stopping target.]{
The magnitude of the magnetic field around the stopping target.
The vertical grey lines indicate the position of the stopping target disks and beam blocker in the target design used in chapter~\sect{sense}.
The muon beam arrives from the left, and signal electrons should leave to the right.
\figlabel{app:tgtImprov:field}}
\end{figure}
}

\newcommand{\FigTgtImprovMirrorVsBeamline}{
\begin{figure}[tb]
\subfloat[][\figlabel{app:tgtImprov:mirror:thetaMin}Minimum Pitch Angle]{%
\includegraphics[width=0.9\textwidth,trim=1.5cm 0 0.8cm 0,clip]{figs/appendixStopTgtImprov/Plot_ThetaMin.pdf}}\\
\subfloat[][\figlabel{app:tgtImprov:mirror:phaseSpace}Fraction of Space Which Can be Mirrored]{%
\includegraphics[width=0.9\textwidth,trim=1.5cm 0 0.8cm 0,clip]{figs/appendixStopTgtImprov/Plot_PhaseSpaceFraction.pdf}}
\caption
[The ability for a particle to be mirrored downstream, if it is initially heading upstream at this point.]{
The ability for a particle to be mirrored downstream, if it is initially heading upstream at this point.
For an electron produced at a given point along the beamline, if its pitch angle is too small it will not be mirrored.
The minimum pitch angle is shown in \protect\subref{fig:app:tgtImprov:mirror:thetaMin}.
Since electrons are produced isotropically, this minimum pitch angle can be used to define the fraction of phase space that will be mirrored back, \protect\subref{fig:app:tgtImprov:mirror:phaseSpace}.
\figlabel{app:tgtImprov:mirror}}
\end{figure}
}

\newcommand{\FigTgtImprovGyroradiusVsBeamline}{
\begin{figure}[tb]
\includegraphics[width=0.9\textwidth,trim=1.2cm 0.2cm 0.8cm 0.2cm,clip]{figs/appendixStopTgtImprov/Plot_GyroradiusGrowth.pdf}
\caption
[Growth in the beam envelope as a function of the distance along the beamline.]{
As the muon beam progresses along the beamline, the beam envelope varies proportional to the square root of the field strength.
The maximum field occurs at the exit of the Torus2 solenoid, with a field strength of about 2.88~T, and this plot shows the relative change in the beam envelope compared to that point.
\figlabel{app:tgtImprov:gyroradius}}
\end{figure}
}

\newcommand{\FigTgtImprovSenseVsBeamline}{
\begin{figure}[tb]
\includegraphics[width=0.9\textwidth,trim=1.1cm 0.3cm 0.8cm 0.3cm,clip]{figs/appendixStopTgtImprov/Plot_ApproxEfficiency.pdf}
\caption
[The effective signal sensitivity achieved by placing a disk of fixed radius at a given location along the beamline.]{
The effective signal sensitivity achieved by placing a disk of fixed radius at a given location along the beamline.
It suggests an optimal target position slightly upstream of the location found in the optimisation chapter, but the difference is easily explained by simplifications in this model.
\figlabel{app:tgtImprov:senseVsBeamline}}
\end{figure}
}

\newcommand{\FigTgtImprovNewGeom}{
\begin{figure}[tb]
\includegraphics[width=0.7\textwidth]{figs/appendixStopTgtImprov/NewGeometry.png}
\caption
[The new geometry of the stopping target, with the beam blocker removed.]{
The new geometry of the stopping target, with the beam blocker removed.
The blue vertical planes are the virtual monitors used in the simulation and not a physical piece of material.
\figlabel{app:tgtImprov:newGeom}}
\end{figure}
}

\newcommand{\FigTgtImprovMuStops}{
\begin{figure}[ptb]
\begin{minipage}[b]{0.35\textwidth}
\subfloat[][\figlabel{app:tgtImprov:muStops:XZ}Above Target (X-Z)]{%
\includegraphics[width=\textwidth,trim=0 0 0 0,clip]{figs/appendixStopTgtImprov/Tidied_Mu_StopPosition_XZ.png}}\\
\subfloat[][\figlabel{app:tgtImprov:muStops:YZ}Side View (Z-Y)]{%
\includegraphics[width=\textwidth,trim=0 0 0 0,clip]{figs/appendixStopTgtImprov/Tidied_Mu_StopPosition_ZY.png}}
\end{minipage}
\subfloat[][\figlabel{app:tgtImprov:muStops:Z}Z Projection]{%
\includegraphics[width=0.5\textwidth,trim=0 0 0 0,clip]{figs/appendixStopTgtImprov/Tidied_Mu_StopPosition_Z.png}}
\caption
[Projections of the muon stopping distribution with the improved target design, to be compared to \fig{sense:stops}.]{
Projections of the muon stopping distribution with the improved target design, to be compared to \fig{sense:stops}.
Although the disks are symmetric, it is clear from \protect\subref{fig:app:tgtImprov:muStops:YZ} how there are fewer muon stops in the upper parts of the disks.
Removing the material there will improve the signal acceptance further.
\figlabel{app:tgtImprov:muStops}}
\end{figure}
}

\newcommand{\FigTgtImprovMuMomentum}{
\begin{figure}[ptb]
\subfloat[][\figlabel{app:tgtImprov:muMomentum:lin}Linear Scale]{%
\includegraphics[width=0.49\textwidth,trim=0 0 0 0,clip]{figs/appendixStopTgtImprov/Tidied_muon_momentum_withDetectedMuons_lin.png}}
\subfloat[][\figlabel{app:tgtImprov:muMomentum:log}Logarithmic Scale]{%
\includegraphics[width=0.49\textwidth,trim=0 0 0 0,clip]{figs/appendixStopTgtImprov/Tidied_muon_momentum_withDetectedMuons_log.png}}
\caption
[The momentum of muons that reach the stopping target and eventually stop.]{
The momentum of muons that reach the stopping target and eventually stop.
By comparison with \fig{sense:muMomenta}, the new target design is clearly able to stop much higher momentum muons, although without the beam blocker
it is also clear that high energy muons reach the detector.
\figlabel{app:tgtImprov:muMomentum}}
\end{figure}
}

\newcommand{\FigTgtImprovSignalDistribution}{
\begin{figure}[tb]
\subfloat[][\figlabel{app:tgtImprov:signalMom:lin}Linear Scale]{%
\includegraphics[width=0.49\textwidth,trim=0 0 0 0,clip]{figs/appendixStopTgtImprov/Tidied_StrawTrk_DetectedSignal.png}}
\subfloat[][\figlabel{app:tgtImprov:signalMom:log}Logarithmic Scale]{%
\includegraphics[width=0.49\textwidth,trim=0 0 0 0,clip]{figs/appendixStopTgtImprov/Tidied_StrawTrk_DetectedSignal_log.png}}
\caption
[The observed  momentum distribution and acceptance for signal electrons.]{
The observed  momentum distribution and acceptance for signal electrons.
Given the extra material, the energy losses in the target are larger than the previous design, although the overall geometric acceptance is significantly increased.
\figlabel{app:tgtImprov:signalMom}}
\end{figure}
}

\chapter{Revisiting the Stopping Target Region}
\sectlabel{appendix:stopTgtImprove}
Of the various optimisations presented in chapter~\sect{phaseII-optimisation}, the stopping target was only optimised for its position, keeping the shape, thickness and size as in the CDR.
It was also observed in that chapter that the beam blocker had a large impact on the signal acceptance.
To investigate ways in which the sensitivity of \phaseII can be further improved we develop here the idea of a beam blocker-free stopping target design, as well as investigate the impact of changing the target shape.

The overall result of this study suggests an improvement in the signal sensitivity of \phaseII of a factor of 2.5 should be achievable.

\section{Understanding the Stopping Target Region}
\FigTgtImprovFieldVsBeamline
\FigTgtImprovMirrorVsBeamline
\FigTgtImprovGyroradiusVsBeamline
\FigTgtImprovSenseVsBeamline
One of the challenges when optimising the stopping target is the complexity of the region, and correlations between the various parameters.

The magnetic field around the stopping target, for example, is not a simple solenoidal field.
\Fig{app:tgtImprov:field} shows the magnitude of the field along the beamline axis in the stopping target region.
At the exit of the bent solenoids and for the first straight section immediately downstream, the field is kept close to 3~T.
By the entrance of the electron spectrometer, the field on axis is closer to 1~T, although a transverse gradient is introduced from the bent solenoids.
The transition from 3 to 1~T happens over 1.5~m across three solenoid sections.

The stated purpose of the reduction in field strength is to mirror those signal electrons that are produced in the target heading upstream.
This, then, is expected to increase the geometric acceptance of the detector.
Magnetic mirroring can be understood in various ways, such as via the adiabatic constants.
However, qualitatively and intuitively, since magnetic field lines never end (`there are no magnetic monopoles'), when the field strength is reduced, the field lines spread out.
This can be seen as a normal solenoidal field, with a radial component superimposed.
  Normal gyration around the solenoidal field, therefore, generates a Lorentz force that points back along the direction of the solenoidal field, and reduces the longitudinal velocity of the particle.
Thus a particle initially with only transverse momentum in the region of tapering field strength will acquire momentum in the longitudinal direction towards the lower field strength.
Similarly a particle initially heading from a low-magnetic field to a high-magnetic field can be mirrored back if the difference in the field is sufficient.
The limit for this to occur is given by the condition~\cite{Jackson}:
\begin{equation}
\tan^{-1}(\theta)=\left|\frac{v_\mathrm{L}}{v_\mathrm{T}}\right|<\left(\frac{B_\textrm{max}}{B}-1\right)^{1/2},
\end{equation}
where $B_\textrm{max}$ and $B$ are the maximum field strength and current field strength.
$v_\textrm{L}$ and $v_\textrm{T}$ are the longitudinal and transverse velocity of the particle at this point, the ratio of which is the arctangent of the pitch angle, $\theta$, between the particle's velocity and the solenoidal axis.

Given that the maximum field strength occurs at the exit of the Torus2, with a magnitude of $B_\textrm{max}=2.88$~T, one can calculate the minimum pitch angle an upstream-heading particle must have in order for it to be mirrored back downstream.
This is shown in \fig{app:tgtImprov:mirror:thetaMin} where the minimum pitch angle varies from around 85\degree to about 43\degree at the end of Stopping Target Section.
Since signal electrons are produced isotropically, this limit can be turned into the fraction of phase space for signal electrons produced at a point along the beamline that will be mirrored back, and this is shown in \fig{app:tgtImprov:mirror:phaseSpace}.
This clearly motivates putting the stopping target further downstream in order to maximize the mirroring effect.

On the other hand, if the target is located downstream, the muon beam is required to traverse the tapered field.
In that case there is another effect, namely the adiabatic growth of the beam envelope, which must be considered.
Since particles in the beam follow the field lines, as the field is reduced, the beam profile must grow.
This growth is governed by the expression:
\begin{equation}
Ba^2=\textrm{const}
\end{equation}
where $B$ is the field strength at a given point, and $a$ the radius of gyration.
Thus if the field strength reduced by a factor of 4, for example, the beam envelope can be expected to grow by a factor of 2.
The radius of gyration is shown in \fig{app:tgtImprov:gyroradius}, relative to the radius at the exit of the Torus2, as a function of the distance along the beamline.

The position optimisation of the stopping target shown in chapter~\sect{phaseII-optimisation} kept the target disks at a fixed radius and moved the whole set of target disks (and beam blocker) upstream and downstream.
With a fixed radius set of disks, as the target was moved downstream, the fraction of the muon beam expected to hit the target was reduced.
This favoured putting the stopping target more upstream in order to increase muon stopping rate.
%In addition, although it would seem that increasing the target disk radius will mean introducing more material to interfere with signal electrons, the electrons will also experience the same growth in beam size,

Thus there were two competing effects: move the target downstream and the mirroring of signal events increases, but move the target upstream and the muon stopping rate is increased.
These factors can be combined to give a proxy for the expected signal sensitivity for a target disk of fixed radius placed at one point along the beamline.
The muon stopping rate will be proportional to the fraction of the beam profile that overlaps with the target.
Meanwhile, the fraction of upstream-pointing phase space that would be mirrored downstream, plus the fraction of phase space pointing downstream in the first place (half of it), should be proportional to the signal acceptance.
This combination is shown in \fig{app:tgtImprov:senseVsBeamline}.

Indeed, the maxima of this function, at 19700 mm in beamline coordinates, is very close to the optimal location found during the position scan.
This neglects many factors, such as the fact that disks upstream will remove muons from the disks downstream, and that not all electrons that head downstream will be accepted into the detector.
Nevertheless, it seems to capture the underlying phenomena that determine the optimal target location.

\section{Improving the Target Design}
In order to improve the signal acceptance, and taking into account the previous discussion, we make three key changes to the target design.
Firstly we allow the target radius to vary as a function of the beamline distance, thus keeping a constant overlap with the muon beam profile.
Secondly we remove the beam blocker, which will hugely increase the acceptance of signal events.
Finally, to compensate for the missing beam blocker, and to increase the number of stopped muons, we add more target disks into the target.
From \fig{sense:stops} in section~\sect{sense:stops} it can be seen that the muon stopping rate per disk has not reached zero by the last disk, which implies more muons could be stopped.
In addition, the increasing target radius can be used to remove the line of sight between the Torus2 exit and the electron spectrometer that the beam blocker previously provided to remove beam particles.

\FigTgtImprovNewGeom
\TabTgtImprovDiskParams
Although each of these changes should each be studied carefully, as a proof of principle the design shown in the event display in \fig{app:tgtImprov:newGeom} was tested in a simulation.
The disk radii and positions are given in \tab{app:tgtImprov:diskParams}.

To check the sensitivity and stopping rates with this configuration, 75~million POT events were fed in to a simulation that had muon \acf{DIO} and nuclear capture disabled, to be able to study the sensitivity in a single simulation.

\subsection{Muon Stopping Rate}
\FigTgtImprovMuStops
\FigTgtImprovMuMomentum
\FigTgtImprovSignalDistribution
From this simulation, the muon stopping rate was found to be \num{2.8e-3}~per POT, and increase of 1.75 times compared to the stopping rate described in chapter~\sect{sense}.
The stopping distributions in the target are shown in \fig{app:tgtImprov:muStops}.
It can be seen from the projection to the Z-axis that we catch much more of the tail of the distribution.
It can also be seen from the YZ projection how there are very few muons stopping in the upper half of the target. 
This suggests one would really be able to remove this material which would give an additional increase to the signal acceptance.

The extra disks and the enlarged radius also mean that the higher energy muons sitting lower down in the beam are more effectively stopped.
Being at the bottom of the beam these muons were more likely to miss the old target design, but also penetrated further, requiring more target material to be stopped.
This can be seen by comparison of the muon momentum distribution shown in \fig{app:tgtImprov:muMomentum} and comparing this with the momenta shown in \fig{sense:muMomenta}.
Whereas before we only stopped muons up to around 50~MeV/c, now muons up to 60~MeV/c are consistently stopped.
As a result of all of these changes, the stopping efficiency for muons that reach the target has increased from 43\% to about 77\%.

Unfortunately, it can also be seen from \fig{app:tgtImprov:muMomentum} that some of the high-momentum muons that do not stop in the target reach through to the detector.
This means the background rate with the design shown here from beam backgrounds would be significantly larger.
However, further optimisations of the target based on this concept might be able to remove these high-momentum muons using additional disks or adjusting the upstream collimators.

\subsection{Signal Acceptance}
\Fig{app:tgtImprov:signalMom} shows the detected momentum distribution of electrons.
Of the signal events that occur, 48\% are detected at the detector, which is the geometric acceptance with this design, and a factor of 2.2 times better than the result in chapter~\sect{sense}.
The momentum cut however is less efficient now, accepting only 50\% of events, compared to the previous 70\%.  
Overall though, the signal acceptance with this target design is a factor 1.55 times better than the old design.

\subsection{Summary and Future Improvements}
By removing the beam blocker, adding more disks and increasing the disk radius along the beam axis, the muon stopping rate and geometric acceptance increase, although the momentum cut efficiency is decreased.
The overall sensitivity can be expected to increase by a factor of 2.5.

Backgrounds from such a design need to be studied, in particular, high-energy electrons and muons in the beam which are no longer removed by the beam blocker, but should be lost in the target.
In addition the \ac{DIO} background, which given the poorer momentum resolution and greater signal acceptance, might produced more background events with this design, and create a significantly higher hit rate.
Other backgrounds like those from antiprotons in the beam, cosmic rays and neutrons will not be changed by the different stopping target.
Radiative pion capture will be increased slightly since the pion stopping rate will also likely increase.

In addition to background studies with this design there is further optimisation that can be performed.
The most obvious area for improvement is to remove the material at the top part of the target disks, since in this region no muons stop here.
In addition one can add more target disks to the lower part to try and catch even more of the high-momentum tail.

The target disk thickness is another parameter that could be studied.
In fact, somewhere between these two options is the possibility to use wedge-shaped target disks, that put less material higher up in the beam.
If one can control this so that the low-momentum muons stop as far downstream as possible, rather than in the first target disk, as currently occurs, then one can improve both the signal acceptance \emph{and} the momentum resolution.

\chapter{Summary of the AlCap Experiment}
\sectlabel{alcap}
The following was originally submitted for the proceedings of NuFact `15 and has been published online from the meeting's home page.
Although peer reviewed, it has not appeared in print, and so I reproduce it here in its entirety.

\end{appendices}

%% Produce the un-numbered back matter (e.g. colophon,
%% bibliography, tables of figures etc., index...)
\begin{backmatter}
  %\input{backmatter}
\chapter{List of Simulation Configurations}
\chapter{List of Acronyms}
%:'<,'>sort '\\acro{\ze.' i
\begin{acronym}
\acro{bsm}[BSM]{Beyond the Standard Model}
\acro{CDC}{Cylindrical Drift Chamber}
\acro{CDR}[CDR]{Conceptual Design Report}
\acro{CyDet}{Cylindrical Detector}
\acro{DIO}[DIO]{decay in orbit}
\acro{JFY}{Japanese Fiscal Year}
\acro{JPARC}[J-PARC]{the Japanese Proton Accelerator Research Complex}
\acro{LOI}[LoI]{Letter of Intent}
\acro{LYSO}{Cerium-doped Lutetium Yttrium Oxyorthosilicate}
\acro{MR}[MR]{J-PARC Main Ring}
\acro{np}[NP]{New Physics}
\acro{PID}{Particle Identification}
\acro{POT}[POT]{Protons on Target}
\acro{RCS}[RCS]{J-PARC Rapid Cycling Synchrotron}
\acro{ses}[SES]{Single-Event Sensitivity}
\acro{sm}[SM]{Standard Model}
\acro{StrECAL}{Straw tube tracker and Electromagnetic Calorimeter}
\acro{TDR}[TDR]{Technical Design Report}
\end{acronym}


\end{backmatter}

\bibliographystyle{plain}
%\bibliographystyle{bensbib}
\bibliography{library/record}

\end{document}
